\documentclass[11pt,a4paper]{article}
\usepackage{fontspec}
\setmainfont{Arial}
\setsansfont{Arial}
\usepackage[ngerman]{babel}
\usepackage[margin=20mm,headheight=14pt]{geometry}
\usepackage{graphicx,xcolor,longtable,array,booktabs,needspace}
\usepackage{amsmath}
\usepackage[unicode,hidelinks]{hyperref}
\usepackage{fancyhdr}
\definecolor{accent}{HTML}{17364A}
\definecolor{teal}{HTML}{087C83}
\pagestyle{fancy}\fancyhf{}
\fancyhead[L]{\small Dissertation | Fragenkatalog und Lernskript}
\fancyfoot[L]{\small Überarbeitete Fachbegriffe | 24.09.2026}
\fancyfoot[R]{\thepage}
\renewcommand{\headrulewidth}{0.3pt}
\setlength{\parindent}{0pt}\setlength{\parskip}{6pt}
\setlength{\emergencystretch}{4em}
\raggedright
\renewcommand{\contentsname}{Inhaltsverzeichnis}
\setcounter{tocdepth}{1}
\newcommand{\topic}[1]{\clearpage\section*{\color{accent}#1}\addcontentsline{toc}{section}{#1}}
\newcommand{\question}[1]{\Needspace{6\baselineskip}\par\vspace{6pt}{\large\bfseries\color{accent}#1}\par\nobreak}
\begin{document}
\begin{titlepage}
\vspace*{22mm}{\Huge\bfseries\color{accent}Dissertation\par}
\vspace{6mm}{\LARGE Fragenkatalog und Lernskript\par}
\vspace{12mm}{\large Vorbereitung für Florian Rzepka\par}
\vspace{16mm}
Einheitlicher LaTeX-Satz mit englischen Fachbegriffen und deutschen Erklärungen.

Enthalten sind die Fragen und Antworten, Rechenübungen, der Abbildungsatlas, die Tabellenprüfung, die Modellübersicht und ein Begriffsschlüssel.

Frage F075 wurde verständlicher formuliert. Die Modellübersicht und ihre offenen Nachweispunkte bleiben erhalten. Die Abbildungen zeigen unveränderte Originalausschnitte der Dissertation.

\vfill Stand der Überarbeitung: 24. September 2026.
\end{titlepage}
\tableofcontents
\topic{Verteidigung der Dissertation}
Fragenkatalog, begründete Antwortvorschläge, Rechenübungen und Abbildungsatlas für Florian Rzepka. Arbeitsstand 16. September 2026, nach Aktualisierung auf Git-Commit 2d8fdc9.

\textbf{Gegenstand:} Deployment-Oriented Neural Network State-Estimation Methods for Battery Management Systems: Ageing Estimation, Robustness Benchmarking, Embedded Optimization, and Hardware Integration.

\textbf{Quellenbasis:} Die angegebene main.pdf mit 230 PDF-Seiten, main.tex, das eingebundene Robustheitskapitel, dessen Anhang und die eingebundenen Ergebnistabellen. Das PDF trägt das Erstellungsdatum 3. September 2026. PDF-Seiten und gedruckte Seiten sind nicht identisch. Im Haupttext entspricht beispielsweise PDF-Seite 134 der gedruckten Seite 104. Der Atlas verwendet ausdrücklich PDF-Seiten und die tatsächlich gedruckten Abbildungsnummern, nicht historische Dateinamen wie Figure\_\allowbreak{}09.

\textbf{Wozu dieses Dokument dient:} Antworten üben, Rechenwege verstehen, Grenzen offen benennen und passende Belege finden. Es sagt nicht vorher, was bestimmte Mitglieder der Kommission persönlich fragen werden. Die Fragen sind fachliche Simulationen aus Batteriephysik, Statistik, maschinellem Lernen und eingebetteten Systemen.

\textbf{Evidenzkennzeichnung:} BELEGT bezeichnet Angaben der Dissertation, nicht eine hier erneut durchgeführte Messung. HERLEITUNG bezeichnet eine nachvollziehbare Rechnung oder fachliche Einordnung. HYPOTHESE bezeichnet einen möglichen Mechanismus ohne isolierenden Nachweis. OFFEN bezeichnet Angaben, die vor der Verteidigung anhand von Skripten, Messprotokollen oder Originalarbeiten geklärt werden sollten. Diese Vorbereitung wiederholt keine Simulation und zertifiziert weder Firmware noch Versuchsdaten.

\textbf{Aufbau:} Teil 1 behandelt roten Faden und Grundbegriffe. Teil 2 behandelt Daten und Referenzen. Teil 3 behandelt den SOH-MLP. Teil 4 behandelt Robustheit und Statistik. Teil 5 behandelt Mikrocontroller, Pruning und Quantization. Teil 6 enthält kritische Rückfragen, Rechenübungen und einen Lernplan. Der anschließende Atlas behandelt alle 67 nummerierten Abbildungen einschließlich Anhang und zeigt die Originalseiten.

\textbf{Lernmethode:} Erst die Frage ohne Antwort lesen. Dann eine Antwort von 30 bis 60 Sekunden geben. Anschließend die ausführliche Begründung und die genannte Grenze erklären. Zahlen immer mit Modell, Datensatz, Einheit und Auswertungsfenster nennen. Ein plausibler Mechanismus ist noch kein experimentell identifizierter Mechanismus.

\topic{1. Roter Faden und wissenschaftlicher Beitrag}
\question{F001. Was ist Ihre Dissertation in drei Sätzen?}
\textbf{Antwort:} Ich untersuche Batteriezustandsschätzung als durchgängige Aufgabe von der zeitlichen Eingangsrepräsentation bis zur Ausführung im BMS. Zuerst zeige ich, wie ein feedforward SOH-Modell durch History Features zeitliche Information nutzen kann, dann vergleiche ich vier SOC-Repräsentanten unter nominalen und gestörten Messungen, anschließend untersuche ich die Kompression und Mikrocontroller-Ausführung rekurrenter SOC- und SOH-Modelle. Die gemeinsame Aussage ist, dass geringer Testfehler allein nicht genügt, weil Störungsverhalten, Recovery, Zustandshandhabung, Speicher und Laufzeit eigenständige Anforderungen sind.

\textbf{Vertiefung:} Die Beiträge sind keine lückenlose Versuchsreihe mit einem einzigen Modell. Kapitel 5 verwendet NMC und ein MLP, Kapitel 6 LFP und unter anderem eine GRU, Kapitel 7 LFP und zwei LSTM-Modelle. Verbunden werden sie durch die Evaluationslogik. \textbf{Beleg:} Kapitel 1 und 8. \textbf{Nicht behaupten:} Alle Schlussfolgerungen seien am selben Modell oder in einem vollständig integrierten Feldsystem gemessen worden.

\question{F002. Was ist tatsächlich neu, wenn LSTM, EKF und Pruning bereits bekannt sind?}
\textbf{Antwort:} Der Beitrag liegt in der anwendungsspezifischen Repräsentation, der kontrollierten Vergleichsmethodik und der nachvollziehbaren Implementierungs- und Messkette. Ich beanspruche weder die Erfindung einer neuen rekurrenten Zelle noch eines allgemein neuen Pruningverfahrens. Der Mehrwert entsteht beispielsweise durch die Trennung von Accuracy, Robustness und Recovery, die zellweise gepaarte Auswertung und den Nachweis, dass identische GRU-Gewichte bei unterschiedlicher Zustandshandhabung sehr unterschiedliche Kosten und Transienten erzeugen.

\textbf{Nachfrage:} Ist das nur Engineering? Eine sinnvolle Antwort benennt prüfbare Erkenntnisse: Wann hilft explizite Historie, welche Fehlermechanismen verändern einen nominalen Vergleich und welche Optimierung adressiert tatsächlich den Hardwareengpass? \textbf{Grenze:} Die Breite des Rahmens ersetzt keine isolierende Ablation jedes Mechanismus. Die eigene Rolle bei Datenaufnahme, Hardware, Software, Analyse und Manuskripten sollte anhand persönlicher Arbeitsnachweise erläutert werden, nicht aus diesem Text erfunden werden.

\question{F003. Welche drei Forschungsfragen beantworten welche Ergebnisse?}
\textbf{Antwort:} Forschungsfrage 1 betrifft die zeitliche Repräsentation mit begrenzter architektonischer Komplexität. Kapitel 5 untersucht kumulierten Strom und Lag-Sequenzen. Forschungsfrage 2 betrifft praktische Robustheit. Kapitel 6 vergleicht nominale Genauigkeit, Messstörungen und Initialisierungs-Recovery. Forschungsfrage 3 betrifft Embedded-Umsetzung. Kapitel 7 vergleicht Base, Pruned und Quantized hinsichtlich Schätzfehler, Speicher und Laufzeit.

\textbf{Prüfpunkt:} Die Energieantwort ist eingeschränkt. Kapitel 7 verwendet eine aus Laufzeit und angenommenen 0,5 W berechnete Energiegröße, keine unabhängige Energiemessung. Die Forschungsfrage wird damit bezüglich Laufzeit und Speicher direkt, bezüglich Energie nur unter einer konstanten Leistungsannahme beantwortet. \textbf{Beleg:} Kapitel 1.1, 7.6 und 8.1.

\question{F004. Warum ist die Arbeit keine Suche nach dem universell besten Modell?}
\textbf{Antwort:} Ein Ranking hängt von Referenzdefinition, Input Features, Training Data, Störungen, Auswertungsfenstern und Hardware ab. DD hat in Kapitel 6 den kleinsten nominalen MAE und führt mehrere Robustheitszusammenfassungen an. HECM hat die bessere beobachtete Recovery. DM und HDM sind wesentlich billiger in der Ausführung. Schon diese Ergebnisse verhindern eine eindimensionale allgemeine Siegerbehauptung.

\textbf{Vertiefung:} Ein vollständiger Familienvergleich müsste mehrere Vertreter, Architekturen, Parametrierungen und Trainingsinitialisierungen einschließen. Hier werden vier festgelegte Repräsentanten verglichen. Die übertragbaren Aussagen betreffen ihre Mechanismen, zum Beispiel Integration ohne kontinuierliche Spannungskorrektur, nicht jede mögliche Implementierung einer Klasse.

\question{F005. Was bedeutet deployment-oriented konkret?}
\textbf{Antwort:} Eingänge müssen zur Laufzeit verfügbar und kausal berechenbar sein. Skalierung, Zustände und Gewichtsexport müssen zwischen Training und Firmware konsistent sein. Die Implementierung muss in den Speicher passen und rechtzeitig fertig werden. Außerdem müssen Störungen und Wiederanlauf betrachtet werden. Deployment ist deshalb mehr als das Umwandeln einer Modelldatei in ein C-Array.

\textbf{Beispiel:} Ein Modell mit 2024 Samples Historie kann nominell kompakt sein und trotzdem bei jedem Ausgabeschritt 2024 rekurrente Schritte wiederholen. Ein dauerhaft fortgeschriebener Zustand vermeidet diese Wiederholung, verändert aber die zeitliche Semantik. Genau diese gekoppelte Entscheidung untersucht Kapitel 6. \textbf{Grenze:} Ausführbarkeit auf einem Entwicklungsboard ist noch keine funktionale Sicherheit oder Feldqualifikation.

\question{F006. Warum drei Anforderungsdimensionen statt eines einzigen Scores?}
\textbf{Antwort:} Performance beschreibt die Brauchbarkeit der Schätzung, Hardware die Ressourcen und Deployment die Umsetzbarkeit und Wartbarkeit. Ein kleiner Fehler kann fehlende Echtzeitfähigkeit nicht kompensieren. Ein kleiner Speicher kann gefährliche Transienten nicht kompensieren. Manche Bedingungen sind harte Grenzen und gehören vor eine gewichtete Nutzwertoptimierung.

\textbf{Vertiefung:} Erst Anforderungen prüfen, etwa maximal zulässige Latenz und sichere Fehlerbehandlung. Danach innerhalb zulässiger Kandidaten Trade-offs vergleichen. Ein gewichteter Score ist eine Entscheidungshilfe, keine Naturkonstante und keine Sicherheitsfreigabe. \textbf{Beleg:} Abbildung 2.1, Kapitel 6.5 und 7.11.

\topic{2. Batteriephysik und Zustandsschätzung}
\question{F007. Was unterscheidet SOC, SOH und verfügbare Energie?}
\textbf{Antwort:} SOC ist ein relativer Ladungszustand gegenüber einer festgelegten nutzbaren Kapazität. Kapazitäts-SOH setzt aktuell gemessene Kapazität ins Verhältnis zu einer nominalen oder anfänglichen Referenz. Energie hängt zusätzlich vom Spannungsverlauf und den Betriebsgrenzen ab. Gleicher SOC bedeutet deshalb nicht automatisch gleiche verfügbare Energie bei unterschiedlichen Zellen oder Alterungszuständen.

\textbf{Herleitung:} Ladung ist das Integral des Stroms, elektrische Energie das Integral von Spannung mal Strom. Kapazitätsverlust und Widerstandsanstieg wirken unterschiedlich auf Energienutzung und Leistungsfähigkeit. Die Dissertation untersucht primär SOC und kapazitätsbezogenen SOH, nicht eine vollständige Leistungs- oder Sicherheitszustandsschätzung. \textbf{Beleg:} Kapitel 3.1 bis 3.2 und 4.1.3.

\question{F008. Warum ist LFP für spannungsbasierte SOC-Schätzung schwierig?}
\textbf{Antwort:} Auf dem breiten OCV-Plateau ändert sich die Gleichgewichtsspannung nur wenig mit SOC. Damit erzeugt ein SOC-Unterschied nur ein kleines Spannungssignal, das von Sensorfehlern, Polarisation und Hysterese überlagert werden kann. Eine lokale Näherung lautet: SOC-Unsicherheit ist ungefähr Spannungsunsicherheit geteilt durch die OCV-SOC-Steigung.

\textbf{Grenze:} Eine kleine Steigung bedeutet eingeschränkte lokale Spannungsinformation, nicht dass LFP grundsätzlich unschätzbar wäre. Zeitliche Dynamik, Stromintegration, informative Randbereiche und zusätzliche Zustandsinformation können helfen. Der Datensatz belegt keine eindeutige Zerlegung aller dieser Informationsquellen. \textbf{Beleg:} Kapitel 3.1.1 und 3.3. \textbf{Übung:} Erkläre, warum eine große Spannungsänderung unter Last nicht automatisch eine große SOC-Änderung ist.

\question{F009. Warum kann gemessener SOH zwischendurch steigen?}
\textbf{Antwort:} Das Diagramm zeigt verfügbare Kapazität aus diskreten Check-ups, nicht direkt die Menge irreversibel verlorenen aktiven Materials. Konditionierung, Relaxation, Temperatur, Abschaltkriterien und Messunsicherheit können einen späteren Kapazitätstest höher ausfallen lassen. Die lineare Interpolation übernimmt diese Schwankungen zwischen den Ankern.

\textbf{Wichtige Grenze:} Aus einer ansteigenden Linie lässt sich weder eine Selbstheilung noch genau eine Ursache ableiten. Die Arbeit nennt mehrere plausible Beiträge, isoliert sie aber nicht. Eine saubere Prüfung würde Check-up-Bedingungen, Rohstromintegration und Temperatur vergleichen und Wiederholungsmessungen auswerten. \textbf{Beleg:} Abbildungen 4.1 und 4.4. Dies ist besonders für die U-förmige NMC-Zelle 9 wichtig.

\question{F010. Wie unterscheiden Sie Kalenderalterung und zyklische Alterung?}
\textbf{Antwort:} Kalenderalterung hängt unter anderem von verstrichener Zeit, Temperatur und Lager-SOC ab. Zyklische Alterung hängt zusätzlich von Durchsatz, Stromraten, Entladetiefe und Betriebsgrenzen ab. In einem laufenden Zyklierexperiment treten beide Beiträge gleichzeitig auf.

\textbf{Vertiefung:} Eine Darstellung über äquivalente Vollzyklen entfernt den Zeiteinfluss nicht. Schnell zyklierte Zellen erreichen denselben Durchsatz in anderer Zeit als langsam zyklierte. Um getrennte kausale Beiträge zu bestimmen, braucht man geeignete Kontrollbedingungen, etwa Kalenderlagerung und ein identifizierbares Versuchsdesign. Die vorliegenden Kurven erlauben zunächst einen Vergleich unter den gewählten Gesamtbedingungen.

\question{F011. Wie funktioniert Coulomb Counting und woher kommt Drift?}
\textbf{Antwort:} Das Verfahren integriert den gemessenen Strom und normiert die Ladungsänderung durch die angenommene Kapazität. Bei positiver Laderichtung steigt SOC, bei positiver Entladerichtung fällt er. Eine Gleichung ist nur zusammen mit dieser Vorzeichenkonvention sinnvoll. Fehler entstehen aus Anfangszustand, Stromoffset, Gain, Zeitbasis, Integrationsnäherung und Kapazitätsannahme.

\textbf{Herleitung:} Ein konstanter Stromoffset b erzeugt ohne Korrektur einen SOC-Fehler mit Betrag b mal Zeit in Stunden geteilt durch Kapazität in Ah. Ein Gainfehler wirkt dagegen auf das vorzeichenbehaftete Stromintegral. Daher müssen beide Fehlerarten getestet werden. Die Dissertation verwendet in Kapitel 6 einen ladungspositiven Qc-Zustand und SOC = 1 + Qc/\allowbreak{}C. \textbf{Lernquelle:} Movassagh et al., Coulomb-Counting-Fehleranalyse, Quellenabschnitt.

\question{F012. Warum begrenzt ein Full-Charge-Reset die Drift nicht vollständig?}
\textbf{Antwort:} Ein Reset setzt einen Zustand an einem erkannten Ladeanker zurück. Danach beginnt die Integration gestörter Ströme erneut. Zudem kann die Erkennung selbst durch Spannung, Timing und Messfehler beeinflusst werden. Die Schätzung kann nach einem Anker korrekt sein und später wieder abweichen.

\textbf{Vertiefung:} Clipping auf 0 bis 1 ist kein physikalisch nachgewiesener Reset. Clipping kann Fehler im Ausgang unsichtbar machen, obwohl intern noch ein Ladungsfehler existiert. In einem rekurrenten Modell kann außerdem eine weitere fehlerbehaftete Historie bestehen. \textbf{Beleg:} Kapitel 6.4.2 und Abbildung 6.6. Deshalb nicht sagen, ein einzelner Volladepunkt lösche jeden über die Lebenszeit entstandenen Fehler.

\question{F013. Was bedeuten die beiden RC-Zweige im HECM?}
\textbf{Antwort:} Sie approximieren dynamische Spannungsanteile mit unterschiedlichen Zeitkonstanten. Ein ohmscher Widerstand beschreibt den unmittelbaren Strom-Spannungs-Anteil, die RC-Zweige verzögerte Polarisation. Im diskreten Modell steht exp(-Δt/\allowbreak{}τ) für das Abklingen eines vorherigen RC-Zustands.

\textbf{Grenze:} Zwei RC-Zweige sind eine reduzierte terminale Beschreibung. Sie identifizieren nicht eindeutig jeweils einen einzelnen elektrochemischen Prozess. Ob ein oder drei Zweige besser wären, müsste unter gleichen Daten- und Ressourcenbedingungen untersucht werden. \textbf{Übung:} Für τ = 10 s und Δt = 1 s bleiben nach einem Schritt rund 90,5 Prozent eines ungetriebenen Zustands übrig. Nach 10 s sind es rund 36,8 Prozent.

\question{F014. Erklären Sie den EKF ohne Formelsammlung.}
\textbf{Antwort:} Zuerst sagt das Modell aus dem bisherigen Zustand und dem Strom den nächsten Zustand und die zu erwartende Spannung voraus. Dann wird die vorhergesagte mit der gemessenen Spannung verglichen. Der Kalman-Gain verteilt dieses Residuum auf die Zustandskorrektur, gewichtet mit geschätzter Modell- und Messunsicherheit.

\textbf{Vertiefung:} P beschreibt Zustandsschätzunsicherheit, Q Prozessunsicherheit und R Messunsicherheit. Die Jacobimatrizen linearisieren die nichtlinearen Zusammenhänge lokal. Ein größeres R verringert bei sonst gleichen Bedingungen den Einfluss der Messung. Das ist keine Garantie für einen besseren Filter. Falsch spezifizierte Unsicherheiten oder ein falsches OCV-Modell können systematische Fehler erzeugen. \textbf{Beleg:} Kapitel 3.3, Gleichungen 3.3 bis 3.4 im PDF prüfen.

\question{F015. Warum nicht UKF, elektrochemisches Modell oder neuronaler Observer?}
\textbf{Antwort:} Der 2RC-EKF bietet einen verbreiteten, überschaubaren Repräsentanten mit expliziter Spannungskorrektur. Ein UKF vermeidet explizite Jacobimatrizen, benötigt aber mehrere Modellfortschreibungen. Ein elektrochemisches Modell kann reichere Zustände darstellen, verlangt mehr Parametrierung und meist mehr Rechenaufwand. Ein neuronaler Observer wäre ein weiterer Hybrid mit eigenem Training.

\textbf{Grenze:} Die Wahl ist eine Begrenzung des Vergleichsraums. Aus den Ergebnissen lässt sich keine Überlegenheit gegenüber nicht untersuchten Beobachtern ableiten. Eine gute Anschlussstudie vergleicht alternative Observer mit demselben Referenzziel, denselben Zellgrenzen und derselben Hardwaremessung.

\topic{3. Daten, Labels und experimentelle Aussagekraft}
\question{F016. Welche Datensätze verwenden Sie genau?}
\textbf{Antwort:} Kapitel 5 verwendet 14 NMC-Zellen mit 2,6 Ah nominaler Kapazität, sieben Alterungsszenarien und je zwei Zellen. Variiert werden mittlerer SOC, DOD, Entladerate und 35 beziehungsweise 45 °C. Kapitel 6 und 7 verwenden 15 LFP-Zellen mit 1,8 Ah aus der MG-FARM-DoE-Kampagne, mit nominal 25 °C Umgebung sowie variierter Lade- und Entladerate und DOD.

\textbf{Vertiefung:} Die gemessene Zelltemperatur ist nicht auf 25 °C beschränkt. Beispielsweise dokumentiert die Split-Tabelle für C29 bis 49,4 °C. Umgebungssollwert und Sensortemperatur dürfen nicht verwechselt werden. Die unterschiedlichen Datensätze erlauben keine direkte Rangfolge ihrer Modell-MAE. \textbf{Beleg:} Tabelle 4.1, Tabelle 4.3 und Kapitel 4.1.

\question{F017. Wie lauten die Zellgrenzen des Robustheitsbenchmarks?}
\textbf{Antwort:} Training: C01, C03, C05, C11, C17, C23. Validation: C07, C19, C21. Test: C09, C13, C15, C25, C27, C29. Die Entwicklungsgrenze gilt für neuronale Modelle, Scaler, Pruning und Lookup-Identifikation.

\textbf{Nachfrage:} Sind das automatisch auch die Splits jeder anderen Studie? Nein. Die Dissertation dokumentiert den sechszelligen Split ausdrücklich für Kapitel 6. Kapitel 5 hat ein anderes NMC-Schema. Für Kapitel 7 muss die genaue Zuordnung aus der zugehörigen Trainingskonfiguration belegt werden, bevor die Robustheitssplits darauf übertragen werden. \textbf{OFFEN:} Eine Backup-Folie sollte für jede Studie separat Training, Validation, Test und gezeigte Beispielzelle enthalten.

\question{F018. Ist der Referenz-SOC wirklich Ground Truth?}
\textbf{Antwort:} Er ist die operational definierte Dataset-Referenz, die aus Stromintegration, Kapazitätsinformationen und bestätigten Ladeankern aufgebaut wurde. Er ist keine direkte, fehlerfreie Messung des inneren SOC. Der Begriff Ground Truth kann als Bezeichnung des Auswertungsziels dienen, wenn seine Konstruktion und Unsicherheit ausdrücklich erklärt werden.

\textbf{Vertiefung:} Modell und Ziel können gemeinsame Messfehler oder Integrationsannahmen teilen. Ein niedriger MAE belegt Übereinstimmung mit dieser Referenz, nicht automatisch höchste physikalische Wahrheit. Unabhängige Referenzvalidierung könnte hochgenaue Strommessung, Unsicherheitsrechnung und kontrollierte Endpunkte umfassen. \textbf{Beleg:} Kapitel 4.1.3. \textbf{Nicht sagen:} Weil nichts Besseres verfügbar ist, sei die Referenz unabhängig oder fehlerfrei.

\question{F019. Warum zwei SOH-Normierungen und welche Falle entsteht dadurch?}
\textbf{Antwort:} NMC und die Embedded-Studie beschreiben Kapazität relativ zur Nennkapazität. Der Robustheitsbenchmark normiert auf die erste gemessene Referenzkapazität der Trajektorie. Beide Größen sind kapazitätsbezogen, aber numerisch nicht automatisch identisch.

\textbf{Herleitung:} Bei Cnom = 1,8 Ah, Cref,0 = 1,9 Ah und C = 1,52 Ah ist nominaler SOH 0,8444 und relativer SOH 0,8. Für die Rekonstruktion der Kapazität muss der passende Nenner wieder verwendet werden. Cnom mal relativer SOH wäre hier 1,44 Ah statt 1,52 Ah. \textbf{OFFEN:} Kapitel 6 schreibt für HDM Ceff = Cnom mal SOH. Vor der Verteidigung klären, auf welcher Basis das tatsächlich eingespeiste LSTM-SOH normiert ist. Unterschiedliche Konventionen können korrekt sein, müssen aber konsistent ineinander überführt werden.

\question{F020. Ist eine offline interpolierte SOH-Referenz ein Kausalitätsverstoß?}
\textbf{Antwort:} Nicht automatisch. Ein Label darf retrospektiv aus später verfügbaren Referenzmessungen erstellt werden. Eine Online-Schätzung darf diese späteren Informationen hingegen nicht als Eingang erhalten. Entscheidend ist die Trennung zwischen Zielkonstruktion und Featurepipeline.

\textbf{Vertiefung:} Interpolation glättet den Zielverlauf und begrenzt, welche Dynamik überhaupt bewertet wird. Eine Schätzung, die zwischen seltenen Check-ups glatt verläuft, ist damit noch nicht als Echtzeitsensor für schnelle Kapazitätsänderungen validiert. \textbf{OFFEN:} Für jedes Modell dokumentieren, welche Kapazitätstests maskiert wurden, welche Diagnosephasen als Eingänge verbleiben und wie initialer SOH verfügbar gemacht wird. Die Initialausrichtung von Kapitel 7 verwendet y0 und ist eine zusätzliche Annahme.

\question{F021. Warum zufällige Zeitfenster nicht als unabhängige Zellen zählen?}
\textbf{Antwort:} Benachbarte Zeitfenster teilen Zellchemie, Fertigungseigenschaften, Historie und oft große Teile derselben Signale. Sie sind korreliert. Millionen Samples oder viele Seeds ersetzen deshalb keine zusätzlichen unabhängigen Zellen.

\textbf{Vertiefung:} Zell-disjunkte Test Data prüfen Transfer auf neue Zellen innerhalb der abgedeckten Domäne. Ein zusätzlicher betriebsbedingungs-disjunkter Split würde stärker prüfen, ob neue Last- oder Temperaturregime generalisiert werden. Auch verschiedene Zellen aus derselben Charge sind nicht automatisch repräsentativ für die gesamte Produktpopulation. \textbf{Beleg:} Kapitel 6.2.6 und Grenzen in Kapitel 8.

\question{F022. Sind Low, Middle und High physikalische Alterungsklassen?}
\textbf{Antwort:} Nein. Es sind deskriptive Lastgruppen anhand des gemessenen P95 des absoluten C-Rates über die ganze Trajektorie. Low umfasst C25 und C27, Middle C09, C13 und C15, High nur C29. Das beschreibt Belastungsabdeckung und ist keine isolierte Kausalanalyse der Alterung.

\textbf{Vertiefung:} P95 ist weniger empfindlich gegenüber einzelnen Maximalwerten als der absolute Peak, hängt aber von der Mischung aus Zyklierung und Diagnosephasen ab. C25 und C27 unterscheiden sich auch innerhalb Low deutlich. Für High gibt es keine zwischenzellige Streuung, weil n = 1. SOH, Last, Temperatur und Betriebsdauer können miteinander gekoppelt sein. \textbf{Beleg:} Tabelle 4.3 und Abbildung 4.5.

\question{F023. Warum gibt es 16 und nicht 18 Zell-SOH-Fenster?}
\textbf{Antwort:} Sechs Zellen mal drei SOH-Bereiche ergäben 18 Kombinationen. C27 bleibt im betrachteten Datensatz vollständig im frischen Bereich, sodass zwei Kombinationen fehlen. Es werden nur beobachtete Zustände ausgewertet und keine künstlichen aged-Fenster ergänzt.

\textbf{Vertiefung:} Dadurch sind die Altersgruppen nicht vollständig balanciert. Ein Unterschied zwischen fresh und aged kann zum Teil daran liegen, dass unterschiedliche Zellen beitragen. Eine alterskausale Aussage verlangt innerhalb derselben Zellen gepaarte Vergleiche und die Kontrolle anderer Betriebsänderungen. \textbf{Beleg:} Kapitel 4.1.2 und 6.3.

\question{F024. Wie wurde verhindert, dass besonders günstige Testfenster ausgewählt wurden?}
\textbf{Antwort:} Innerhalb jeder vorhandenen Zell-SOH-Kombination wird ein gemessenes Fenster anhand von SOH, Temperatur, P95-C-Rate, Durchsatz und Anteil niedrigen SOCs gewählt. Der Abstand zur robust skalierten typischen Feature-Lage entscheidet, nicht der Modellfehler. Damit ist die Auswahl nicht direkt auf einen Estimator optimiert.

\textbf{Kritische Präzisierung:} Die dargestellte Formel wählt den Punkt mit geringstem Abstand zum Coordinate-wise Median, nicht zwingend den klassischen Medoid mit minimaler Summe aller paarweisen Abstände. Das sollte terminologisch erklärt werden. \textbf{OFFEN:} Wie behandelt das Skript eine MAD von null? Welche SOH-Grenzen gelten genau? Warum repräsentieren typische Fenster auch seltene gefährliche Zustände? Letzteres tun sie nicht vollständig. Daher ergänzen Störungsversuche und Ereignisanalysen die Auswahl.

\question{F025. Was bringt das DoE und was beweist es nicht?}
\textbf{Antwort:} Das DoE verteilt endliche Versuchskapazität strukturiert über Lade-C-Rate, Entlade-C-Rate und Entladetiefe. Dadurch entstehen systematisch verschiedene Betriebsbedingungen statt zufälliger Einzelprofile. Die Arbeit nutzt diese Diversität für Zustandsschätzung und Robustheitsbewertung.

\textbf{Grenze:} Die Darstellung eines DoE-Würfels ist noch keine statistische Identifikation aller Haupteffekte und Interaktionen. Dafür braucht man den genauen Designplan, Wiederholungen, Auswertungsmodell und gegebenenfalls eine Alias-Struktur bei fraktionellen Designs. Die drei Faktoren sind nicht die einzigen physikalischen Einflussgrößen. Alterung, Eigenerwärmung und Fertigungsstreuung bleiben relevant.

\topic{4. SOH-MLP und zeitliche Features}
\question{F026. Wie erhält ein speicherloses MLP zeitliche Information?}
\textbf{Antwort:} Die Historie wird vor dem Netz in einen festen Featurevektor geschrieben. Lag-Sequenzen enthalten vergangene Spannungs-, Temperatur- und kumulierte Stromwerte. Das Netz verarbeitet diese gemeinsam, ohne einen internen Recurrent State zwischen zwei Aufrufen zu benötigen.

\textbf{Vertiefung:} Speicherlos ist nur das Netz als Funktion. Das Gesamtsystem benötigt den Historienpuffer, Resampling und kumulative Zustände. Der Vergleich lautet daher nicht mit oder ohne Gedächtnis, sondern explizit repräsentierte versus intern gelernte Historie. \textbf{Beleg:} Abbildungen 5.2 und 5.4. \textbf{Nachfrage:} Warum nicht aktueller Strom allein? Throughput Features enthalten längerfristige Betriebsinformation, können aber auch als Zeit- oder Altersproxy fungieren.

\question{F027. Warum getrennte kumulierte Lade- und Entladeströme?}
\textbf{Antwort:} Ein vorzeichenbehaftetes Gesamtintegral kann nach einem vollständigen Zyklus wieder nahe null liegen, obwohl die Zelle erheblichen Durchsatz erfahren hat. Getrennte oder betragsbezogene Integrale erhalten diese Beanspruchungsinformation. Die physikalische Relevanz liegt in der Vorgeschichte, nicht in einer direkten Gleichsetzung von Ah und Kapazitätsverlust.

\textbf{Grenze:} Korrelation mit SOH kann durch gemeinsame Zeitabhängigkeit entstehen. Für einen Nachweis des zusätzlichen Feature-Nutzens wären Ablationen gegen Zeit, Zykluszahl, kumulierten Durchsatz und einfache Regressionsbaselines hilfreich. \textbf{Beleg:} Kapitel 5.2.1 bis 5.2.2. Die Vorzeichen und Resetregeln dieser NMC-Features nicht mit Qc des SOC-Benchmarks vermischen.

\question{F028. Bedeutet die rote Korrelationsfarbe positive Korrelation?}
\textbf{Antwort:} In Abbildung 5.1 ausdrücklich nicht. Die Farbstärke kodiert den Betrag des Pearson-Koeffizienten. Sowohl -1 als auch +1 liegen am roten Ende, null am blauen. Das Vorzeichen muss aus der gedruckten Zahl gelesen werden.

\textbf{Vertiefung:} Pearson misst linearen Zusammenhang und ist empfindlich gegenüber Trend, Ausreißern und Mischungen verschiedener Zellen. Hohe Korrelation beweist weder Ursache noch Nichtredundanz. Ein Netz kann aus zwei stark korrelierten Features dennoch unterschiedliche Informationen gewinnen, oder nur dieselbe Zeitentwicklung doppelt sehen. \textbf{Übung:} Erkläre an einer negativen Durchsatz-SOH-Korrelation, warum ein positives Durchsatzsignal einen sinkenden Zielwert begleiten kann.

\question{F029. Warum gerade 16 Lag-Schritte und 10 Minuten?}
\textbf{Antwort:} Das ist das im untersuchten Raster günstige Setting. Es verbindet eine relativ breite Historie mit einer zeitlichen Aggregation, die noch Informationen über typische Lade- und Entladephasen bewahrt. Sehr grobe Aggregation kann Dynamik verwischen, sehr kurze Intervalle können stark redundante Eingänge liefern.

\textbf{Präzisierung:} Die Arbeit nennt 160 Minuten. Bei 16 punktförmigen Stützstellen inklusive aktuellem Zeitpunkt beträgt der Abstand vom ältesten bis zum neuesten Punkt dagegen 15 mal 10 = 150 Minuten. 160 Minuten kann korrekt die Abdeckung von 16 aufeinanderfolgenden 10-Minuten-Bins bezeichnen. \textbf{OFFEN:} Die genaue Bin- und Indexkonvention am Skript prüfen. Das ist eine Definitionsfrage, keine Widerlegung des Befunds. \textbf{Beleg:} Abbildung 5.5 und Kapitel 5.5.

\question{F030. Ist Ihr MLP tatsächlich klein?}
\textbf{Antwort:} Es ist strukturell einfach, weil es ein Feedforward-Netz ohne rekurrente Zustandsverwaltung ist. Die finale Tabelle listet aber neun Hidden Layers mit 128, 256, 256, 256, 128, 512, 128, 256 und 256 Neuronen. Deshalb sollte einfach nicht mit nachgewiesen geringer Parameterzahl verwechselt werden.

\textbf{Herleitung:} Unter der Annahme von 48 Eingängen, also drei Features mal 16 Lag-Werte, ergeben die dichten Schichten einschließlich Bias 434561 Parameter. Bei FP32 sind das rund 1,74 MB reine Parameterdaten. Der Wert ist eine Architektur-Rechnung, kein verifizierter Modellexport. \textbf{OFFEN:} Tatsächliche Eingangsdimension und gespeichertes Modell prüfen. Der anfängliche Suchkandidat mit acht Neuronen beschreibt nicht das finale Modell. \textbf{Beleg:} Tabelle 5.1.

\question{F031. Was macht Hyperband und warum wurde es verwendet?}
\textbf{Antwort:} Hyperband verteilt ein begrenztes Trainingsbudget auf viele Konfigurationen und verwirft schwache Kandidaten früh. Es kombiniert verschiedene Startbudgets und aufeinanderfolgende Ausleseschritte. Gesucht wurden unter anderem Historienlänge, Auflösung, Tiefe, Breite und Trainingsparameter.

\textbf{Grenze:} Frühe Trainingsleistung ist nicht immer ein zuverlässiger Prädiktor späterer Leistung. Ein langsamer konvergierender guter Kandidat kann zu früh ausscheiden. Die Auswahl muss an Validation erfolgen und das finale Testset unberührt lassen. \textbf{OFFEN:} Wenn Abbildung 5.5 Testzellen zeigt, darf deren Optimum nicht als ausschließlich validation-basierte Auswahl dargestellt werden, bevor die historische Auswertung geklärt ist. \textbf{Lernquelle:} Hyperband-Originalarbeit im Quellenabschnitt.

\question{F032. Wie zuverlässig ist die Validation mit 20 Prozent der Training Sequences?}
\textbf{Antwort:} Die vier vollständig zurückgehaltenen NMC-Testzellen schaffen eine wichtige unabhängige Testgrenze. Innerhalb des Trainingspools können zufällig getrennte, überlappende Sequenzen jedoch eine optimistische Validation erzeugen. Das hängt von der tatsächlichen Aufteilung und Fensterüberlappung ab.

\textbf{OFFEN:} Die Dissertation nennt keine vollständige genaue Sequenz-Partition in diesem Absatz. Zu prüfen sind zellweise oder zeitblockweise Trennung, Überschneidungen und zeitlicher Abstand zwischen Fenstern. Für eine strengere Wiederholung wären gruppierte Splits oder zeitliche Sperrzonen sinnvoll. Nicht behaupten, dass eine potenzielle Validierungsleckage tatsächlich vorlag, ohne die Skripte zu prüfen.

\question{F033. Warum ist die U-förmige Zelle 9 schlecht geschätzt?}
\textbf{Antwort:} Das Modell folgt nur Teilen der Referenzentwicklung und hat größere absolute und systematische Abweichungen. Die U-Form allein erzeugt keinen Fehler. Würde sie exakt reproduziert, wären MAE und MSE klein. Der Befund spricht für eine Grenze der Generalisierung unter dieser ungewöhnlichen Trajektorie.

\textbf{Beleg:} Kapitel 5.5 nennt MAE 2,528 Prozentpunkte und MSE 8,327 Quadrat-Prozentpunkte für Zelle 9. Mögliche Ursachen sind unzureichende Trainingsabdeckung, Label-/\allowbreak{}Check-up-Besonderheiten und eine unpassende Feature Representation. Welche davon dominiert, ist nicht isoliert belegt. \textbf{Nicht sagen:} Das Modell versagt, weil SOH niemals steigen dürfe.

\question{F034. Was zeigen die guten NMC-Fehlerwerte wirklich?}
\textbf{Antwort:} Für günstige Einzelzellen nennt die Arbeit MAE 0,487 und 0,253 Prozentpunkte, im Mittel 1,10 Prozentpunkte. Das ist eine Genauigkeitsaussage gegenüber der kapazitätsbasierten Referenz unter den untersuchten Bedingungen. 1,10 Prozentpunkte entsprechen 0,0110 auf der normierten Skala.

\textbf{Vertiefung:} MSE skaliert quadratisch. 2,95 Quadrat-Prozentpunkte entsprechen 0,000295 in normierten Einheiten, nicht 0,0295. Die dargestellte Glättung darf nicht mit ungefilterter Netzleistung gleichgesetzt werden. \textbf{OFFEN:} Vortragen können, auf welchen Zellen und welchen Roh-/\allowbreak{}Filterausgaben die angegebenen Mittel beruhen. Drei gezeigte Beispielzellen sind nicht automatisch das vollständige Testset aus vier Zellen.

\question{F035. Warum verbessert Entfernen problematischer Trainingszellen nicht alles?}
\textbf{Antwort:} Das Entfernen von C3/\allowbreak{}C4 im NMC-Versuch reduziert mögliche unregelmäßige Daten, entfernt aber auch Variabilität. Der mittlere MAE sinkt laut Arbeit nur von 1,10 auf 1,07 Prozentpunkte, während die atypische Zelle 9 instabiler geschätzt wird. Datenbereinigung und Repräsentativität können gegensätzliche Effekte haben.

\textbf{Grenze:} Ohne Wiederholungen über Trainingsseeds ist diese kleine Differenz nicht automatisch statistisch belastbar. Eine klare Antwort lautet, dass die Untersuchung keinen einfachen Zusammenhang weniger problematische Daten gleich bessere Generalisierung zeigt. Dafür bräuchte man gepaarte Wiederholung, dokumentierte Auswahlregeln und Unsicherheiten.

\question{F036. Ist dies Alterungsprognose oder Zustandsschätzung?}
\textbf{Antwort:} Das Netz schätzt den aktuellen SOH aus bis dahin vorliegenden Operating Features. Eine über die Zeit gezeichnete Folge solcher Schätzungen ist noch keine Vorhersage zukünftiger Restlebensdauer. Eine Prognose braucht einen Zukunftshorizont und Annahmen über die zukünftige Nutzung.

\textbf{Nachfrage:} Kann es für eine neue Chemie benutzt werden? Die Verfahrensidee ist übertragbar, die gelernten Gewichte und optimalen Zeitskalen nicht ohne Validierung. NMC- und LFP-Spannungsverläufe und Alterungsmechanismen unterscheiden sich. Neu trainieren oder anpassen und zell-disjunkt testen wäre der saubere Weg.

\topic{5. Robustheitsbenchmark und Recovery}
\question{F037. Erklären Sie die vier Repräsentanten in einer Minute.}
\textbf{Antwort:} DM integriert Strom mit fester Kapazität und einem Volladeanker. HDM benutzt denselben Kern, skaliert die Kapazität aber mit geschätztem SOH. HECM ergänzt eine 2RC-Zustandsbeschreibung, SOH-abhängige Tabellen und EKF-Spannungskorrektur. DD ist eine geprunte GRU, die Messkanäle und abgeleitete Features gemeinsam aus einer kausalen Sequenz auf SOC abbildet.

\textbf{Präzisierung:} HDM, HECM und DD erhalten denselben kausalen SOH-Verlauf. Das kontrolliert einen gemeinsamen Eingang, macht die Modelle aber nicht identisch empfindlich gegen SOH-Fehler. DD ist kein raw-signal-only-Modell, weil Qc, Ableitungen, Zeitabstand und SOH bereits Vorwissen enthalten. \textbf{Beleg:} Kapitel 6.2.1 bis 6.2.3.

\question{F038. Wie sieht die DD-GRU exakt aus?}
\textbf{Antwort:} Der Ausgangspunkt hat 96 Recurrent Units. Strukturelles Pruning entfernt 29 vollständige Einheiten und zugehörige Verbindungen, sodass 67 bleiben. Danach folgen ein MLP mit 96 Hidden Units und ein Sigmoid-Ausgang. Die Eingänge sind U, I, T, SOH, Qc, dU/\allowbreak{}dt, dI/\allowbreak{}dt und Δt. Im primären Benchmark verwendet jede Ausgabe die letzten 2024 Samples und beginnt die Fensterrechnung mit zurückgesetztem Recurrent State.

\textbf{Grenze:} Das ist nicht das 64er SOC-LSTM aus Kapitel 7. Die 30,2 Prozent Breitenreduktion entspricht auch nicht exakt einer 30,2-prozentigen Laufzeit- oder Parameterreduktion. Quadratische rekurrente Matrizen verändern sich stärker. \textbf{Beleg:} Kapitel 6.2.2.

\question{F039. Warum bekommt DD Qc, obwohl daraus direkt SOC berechnet werden kann?}
\textbf{Antwort:} Qc ist ein kausal integriertes Charge Feature und damit ein starker, physikalisch sinnvoller SOC-Prädiktor. DD kann diese Information mit Spannungs-, Temperatur- und Zeitkontext verbinden. Es ist legitim, dieses Feature zu nutzen, solange seine Konstruktion online verfügbar ist und dieselben Störungen erfährt wie die Rohmessungen.

\textbf{Kritische Grenze:} Wenn das eingegebene Qc bereits offline driftkorrigierte Zielinformation enthielte, wäre der Vergleich problematisch. Deshalb trennt die Arbeit Online-Qc von der Dataset-Referenz. Ohne Feature-Ablation ist nicht gemessen, welcher Anteil des Erfolgs allein auf Qc oder auf die GRU zurückgeht. Eine sinnvolle Vergleichsbaseline wäre ein einfacher Regressor mit denselben Features.

\question{F040. Was bedeutet gemeinsame kausale SOH-Spur?}
\textbf{Antwort:} Aus fünf Basissignalen werden stündlich Mittelwert, Standardabweichung, Minimum und Maximum gebildet. Diese 20 Features gehen in ein gemeinsames SOH-Netz mit Feature-Embedding, zweilagigem LSTM und MLP-Blöcken. Seine Zustände laufen kausal weiter, und seine Ausgabe wird zwischen Stundenupdates gehalten. Alle SOH-abhängigen SOC-Zweige erhalten diese identische Information.

\textbf{Vertiefung:} Vor dem Test gibt es 192 Stunden ungestörten Kontext. Nach Störungsbeginn werden betroffene abgeleitete Features aus dem gestörten Stream rekonstruiert. Das ist wichtig, weil andernfalls künstlich perfekte Nebenkanäle verbleiben könnten. \textbf{Grenze:} Gleiche SOH-Spur kontrolliert den Eingang, beseitigt aber keine SOH-Unsicherheit. Im Hardwareteil von Kapitel 6 wird sie extern geliefert, nicht auf dem Controller berechnet.

\question{F041. Wie unterscheiden sich Accuracy, Robustness und Recovery?}
\textbf{Antwort:} Accuracy misst die Abweichung vom Referenz-SOC unter nominalen Messungen. Robustness beschreibt die zusätzliche Verschlechterung unter korrumpierten oder fehlenden Eingängen. Recovery beschreibt, wie eine initial gestörte Schätzung wieder an die entsprechend korrekt initialisierte Schätzung herankommt.

\textbf{Wichtig:} Recovery in dieser Arbeit ist eine gepaarte Modell-zu-Modell-Abweichung, keine Garantie kleiner Abweichung vom Dataset-SOC. Zwei gleich falsche Trajektorien können vollständig zueinander recovern. Deshalb muss die Accuracy separat gezeigt werden. \textbf{Beleg:} Kapitel 6.2.6 und Abbildung 6.9.

\question{F042. Wieso haben alle Modelle denselben Auswertungsbeginn?}
\textbf{Antwort:} Die Rolling-GRU liefert erst nach ausreichender Historie eine gültige Ausgabe. Damit frühe schwierige oder leichte Samples nicht nur bei manchen Modellen in den Fehler eingehen, beginnt die gemeinsame Auswertung erst mit der ersten gültigen DD-Ausgabe. Alle Modelle werden auf demselben Zeitintervall verglichen.

\textbf{Grenze:} Das blendet einen Teil des realen Kaltstartproblems aus. Recovery-Zeit ist weiterhin auf die Intervention bezogen, während der gemeinsame beobachtete Bereich später beginnt. Die 2024 Samples sind bei 1 Hz ungefähr 33,7 Minuten. Sehr frühe Unterschiede und Recovery müssen deshalb hinsichtlich Beobachtungsgrenze vorsichtig interpretiert werden.

\question{F043. Was genau wird bei Initialisierungsfehlern verändert?}
\textbf{Antwort:} Es wird ein SOC-äquivalenter Fehler von -0,10 eingebracht. Bei DM/\allowbreak{}HDM betrifft dies den zugänglichen Integrationszustand, bei HECM den expliziten SOC-Zustand, bei DD einen Qc-Offset von -0,10 mal nominaler Kapazität. DD besitzt keinen separaten SOC-Speicher, der direkt um zehn Prozentpunkte verschoben werden könnte.

\textbf{Grenze:} Die Interventionen sind funktional vergleichbar, aber nicht identisch in ihren latenten Zuständen oder in ihrer unmittelbaren Ausgangswirkung. Bei HDM hängt die SOC-Wirkung eines Ladungsoffsets zusätzlich vom Kapazitätsnenner ab. Deshalb ist es ein operationeller Re-Anker-Test und kein universeller Konvergenzvergleich identischer Anfangsfehler. \textbf{Beleg:} Kapitel 6.3 und 6.4.5.

\question{F044. Was bedeuten 0,02, 300 Sekunden und persistent?}
\textbf{Antwort:} Zunächst muss der Betrag der Differenz zwischen gestörter und sauber initialisierter Ausgabe höchstens 0,02 SOC, also zwei Prozentpunkte, für mindestens 300 Sekunden betragen. Das ist ein erster qualifizierter Bandeintritt. Persistent ist die Rückkehr erst, wenn die Differenz danach bis zum Ende des 24-Stunden-Horizonts innerhalb des Bandes bleibt.

\textbf{Beispiel:} DD tritt nach einer Stunde für zehn Minuten ein, verlässt das Band wieder und bleibt erst nach drei Stunden dauerhaft darin. First entry ist eine Stunde, persistent recovery drei Stunden. Läuft der Versuch ohne dauerhafte Rückkehr aus, ist er rechtszensiert. \textbf{Grenze:} Das nachträgliche persistent-Kriterium ist eine Auswertedefinition. Online kann man noch nicht wissen, ob später ein Rückfall kommt.

\question{F045. Warum kann DD kurz recovern und danach wieder schlechter werden?}
\textbf{Antwort:} DD kombiniert Qc mit wechselnden Messverläufen innerhalb eines rollenden Fensters. Die Vorhersage kann zeitweise wenig und später wieder stärker vom fehlerhaften Feature abhängen. Außerdem ändern sich SOC-Region und Dynamik. Der erste Bandeintritt beweist keine dauerhafte Entfernung aller Ursachen.

\textbf{Beleg:} DD hat eine Rückfallfraktion von 0,78 gegenüber 0,17 bei HECM. Die beobachteten persistenten Recovery-oder-Zensur-Mittel sind 2,60 h für DD und 1,20 h für HECM. \textbf{Grenze:} Die Feature-Erklärung ist plausibel, aber keine gemessene Attribution ohne Ablation. DM/\allowbreak{}HDM haben jeweils eine hohe persistente Zensurfraktion von 0,89.

\question{F046. Warum darf man zensierte 24 Stunden nicht als echte Recovery interpretieren?}
\textbf{Antwort:} Bei einem zensierten Lauf ist nur bekannt, dass Recovery bis zum Ende nicht beobachtet wurde. Der tatsächliche Zeitpunkt könnte später liegen oder nie eintreten. Der als Recovery-or-censor zusammengefasste Wert ist daher eine begrenzte Beobachtungskennzahl, nicht die mittlere physikalische Konvergenzzeit.

\textbf{Vertiefung:} Eine Survival-Auswertung könnte Zeit-bis-Ereignis mit Zensierung expliziter behandeln. Mit sechs Zellen bliebe ihre Unsicherheit groß. Für die Verteidigung die Zensurquote immer zusammen mit der Zeit nennen. Gerade ein Balken bei rund 24 h bedeutet häufig keine erfolgreiche Recovery nach 24 h.

\question{F047. Welche nominalen Zahlen müssen Sie sicher kennen?}
\textbf{Antwort:} Im Robustheitsbenchmark beträgt der zellweise gleichgewichtete MAE DM 0,0690, HDM 0,0412, HECM 0,0337 und DD 0,0258. Die RMSE-Werte sind 0,0740, 0,0442, 0,0388 und 0,0300. Das entspricht MAE von 6,90, 4,12, 3,37 und 2,58 Prozentpunkten.

\textbf{Grenze:} Es sind sechs Zellen und 16 standardisierte Fenster, nicht alle Samples aller Lebenszeiten gleich gewichtet. Die Hierarchie ist für diese Vertreter nominal klar, aber die Holm-korrigierten Signifikanztests liegen über 0,05. Deshalb praktische Effektgröße und Unsicherheit nennen, nicht eine pauschal statistisch gesicherte universelle Überlegenheit. \textbf{Beleg:} Abbildung 6.4.

\question{F048. Warum kann ein Bias den MAE verbessern?}
\textbf{Antwort:} MAE vergleicht mit einer Referenz, und die nominale Schätzung kann bereits einen gerichteten Fehler haben. Eine zusätzliche Störung kann diesen Fehler teilweise kompensieren. Das ist keine Aussage, dass ungenauere Sensoren grundsätzlich besser wären.

\textbf{Herleitung:} Hat die nominale Schätzung Fehler +0,04 und verschiebt die Störung sie um -0,02, sinkt der absolute Fehler von 0,04 auf 0,02. Bei +0,02 zusätzlicher Verschiebung steigt er auf 0,06. Deshalb werden Gain und Stromoffset mit beiden Vorzeichen getestet. \textbf{Nicht tun:} Ergebnisse oder Vorzeichen nach gewünschter Schlussfolgerung auswählen, ohne die adverse Richtungsdefinition offenzulegen.

\question{F049. Was bedeutet adverse-direction und ist das ein Worst-Case-Beweis?}
\textbf{Antwort:} Für jede Zelle und Betragsstufe wird die größere MAE-Änderung der beiden getesteten Vorzeichen verwendet. Anschließend werden Zellen gleich gewichtet. Das ist die ungünstigere der zwei untersuchten Richtungen, nicht der mathematische Worst Case aller möglichen Fehlerverläufe.

\textbf{Vertiefung:} Diese nichtlineare Auswahl unterscheidet sich vom Mittel beider Vorzeichen. Sie kann auch kleine positive Effekte betonen und muss als Auswahlregel vor der Aggregation bekannt sein. Ein monotoner Verlauf im betrachteten Raster ist ein Befund, kein allgemeiner Satz. Bei Clipping, Resets und nichtlinearen Modellen muss Monotonie nicht immer gelten. \textbf{Beleg:} Abbildung 6.5.

\question{F050. Warum ist additiver Stromoffset viel kritischer als Gain?}
\textbf{Antwort:} Offset wird auch bei nominal null Strom integriert. Gain ist proportional zum Strom und kann sich zwischen Laden und Entladen teilweise kompensieren. Bei 50 mA entstehen in 24 Stunden 1,2 Ah falsche Ladung, was bei 1,8 Ah etwa 66,7 Prozentpunkten ungebremster SOC-Verschiebung entspräche.

\textbf{Grenze:} Das ist keine Vorhersage des gemessenen MAE, weil Anker, Clipping, Lastwechsel und Korrektur wirken. Tatsächlich sind die adversen MAE-Zunahmen DD 0,0401, HECM 0,1763, DM 0,2467 und HDM 0,2727. Die Größenordnung erklärt, warum Nullpunktkalibrierung wichtig ist. \textbf{Beleg:} Kapitel 6.4.3.

\question{F051. Warum ist HDM beim Offset schlechter als DM, obwohl nominal besser?}
\textbf{Antwort:} HDM reduziert nominal Kapazitätsmismatch, besitzt aber weiterhin den Integrationskern. Bei gealterter Zelle kann der kleinere effektive Kapazitätsnenner denselben falschen Ladungsbetrag in eine größere SOC-Abweichung umsetzen. Damit kann bessere nominale Kapazitätsanpassung zugleich die Offsetempfindlichkeit erhöhen.

\textbf{Grenze:} Dieses Argument erklärt eine strukturelle Möglichkeit. Der exakte Zahlenunterschied enthält zusätzlich SOH-Verlauf, Anker, Clipping und Zellprofil. Die Arbeit isoliert nicht jeden Anteil. Die richtige Schlussfolgerung ist, dass nominal bessere Adaptation nicht automatisch höhere Robustheit gegen einen anderen Fehlermechanismus liefert.

\question{F052. Akkumuliert der Gainfehler über die gesamte Lebenszeit?}
\textbf{Antwort:} Nicht zwangsläufig monoton. Ohne Korrektur folgt der Integrationsfehler dem vorzeichenbehafteten Ladungsdurchsatz. Volladeanker verändern den Zustand, Lade- und Entladephasen ändern das Vorzeichen des Integrals, und Modellkorrekturen wirken zusätzlich. Die Lebenszeitanalyse zeigt eher wiederkehrenden Aufbau und teilweise Re-Ankerung als unbegrenzte Drift.

\textbf{Beleg:} Im kontinuierlichen High-Load-Replay steigen MAE um 0,0102 für DM, 0,0072 für HDM, 0,0037 für HECM und 0,0048 für DD. Hier hat HECM die kleinste Zusatzlast, obwohl DD im sechszelligen Fenstervergleich günstiger ist. \textbf{Grenze:} Eine Zelle über die volle Lebenszeit und mehrere standardisierte Testfenster beantworten verschiedene Fragen. Dieser Unterschied ist kein logischer Widerspruch.

\question{F053. Warum verursachen hohe lokale Ausschläge manchmal kaum globalen MAE?}
\textbf{Antwort:} Ein Mittel über viele Stunden verdünnt kurze Ereignisse. Ein Modell kann einen einzelnen gefährlichen Peak haben und trotzdem fast unveränderten mittleren Fehler. Deshalb verwendet die Arbeit ereignisausgerichtete Maxima, P95-Peaks und Excess Error zusätzlich zu globalen Kennzahlen.

\textbf{Herleitung:} Ein Fehler von 0,2 für einen einzigen Sample unter 86400 Samples trägt isoliert nur etwa 0,00000231 zum Tages-MAE bei. Ob der restliche Verlauf zusätzlich gestört wird, ist separat zu prüfen. \textbf{Beleg:} Abbildungen 6.12 und 6.13. DD zeigt eine lokale Spike-Empfindlichkeit, die in der globalen Heatmap fast unsichtbar ist.

\question{F054. Warum verstärken Ableitungsfeatures Rauschen und Spikes?}
\textbf{Antwort:} Eine Backward Difference zieht zwei Messungen voneinander ab und teilt durch den Zeitabstand. Unabhängiges Messrauschen addiert sich in der Varianz. Bei gleichem Δt hat die Differenz zweier unabhängiger Fehler mit Varianz σ² die Varianz 2σ². Kleine Zeitabstände vergrößern zusätzlich die Ableitungsamplitude.

\textbf{Vertiefung:} Ein einzelner Spannungsspike erzeugt typischerweise eine große positive und danach negative Ableitungsstörung. DD sieht damit sowohl absolute Spannung als auch dU/\allowbreak{}dt. Eine Filterung könnte helfen, verändert aber Verzögerung und Trainingsverteilung. \textbf{HYPOTHESE:} Dieser Pfad ist ein plausibler Grund für DD-Transienten, nicht durch kanalweise Ablation bewiesen.

\question{F055. Warum kann DM trotz fehlender Spannungsregelung auf Spannungsrauschen reagieren?}
\textbf{Antwort:} Spannung beeinflusst die Volladeerkennung. Wenn Rauschen den Zeitpunkt des bestätigten Ankers verändert, ändert sich der Qc-Reset und damit der spätere SOC-Verlauf. Spannung wird also nicht kontinuierlich zur Korrektur benutzt, ist aber nicht wirkungslos.

\textbf{Beleg:} Kapitel 6.4.4 berichtet kleine MAE-Zunahmen bei Spannungsrauschen auch für DM und HDM. \textbf{Grenze:} Ohne Ereignislog der Resetentscheidungen ist die Ankererklärung plausibel, aber nicht für jeden konkreten Knick direkt bewiesen. Eine sichere Diagrammantwort nennt den beobachteten Knick und den zu prüfenden Resetzeitpunkt, statt jede Bewegung nachträglich zu erklären.

\question{F056. Warum ist Jitter hier weniger kritisch als Sample Loss?}
\textbf{Antwort:} Der Jittertest erhält Messinformation und rekonstruiert zeitabhängige Größen mit den gestörten Zeitabständen. Beim Sample Loss wird Information durch gehaltene Werte ersetzt. Insbesondere fehlende Ladung kann Integration nicht nachträglich rekonstruieren.

\textbf{Grenze:} Das ist keine Entwarnung für beliebige asynchrone Sensoren. Kanalversatz, falsche Timestamps, Frequenzdrift oder aliasierte hochfrequente Signale sind andere Fehler. Die konkrete Jitterimplementierung muss zu der behaupteten realen Störung passen. \textbf{Beleg:} Abbildung 6.10, drei Stufen ±0,1, ±0,5 und ±0,9 s.

\question{F057. Was wird beim Burst Dropout eingefroren und warum 48 Stunden?}
\textbf{Antwort:} Der Versuch simuliert einen einstündig eingefrorenen Messstream. Der Start liegt im zulässigen Abschnitt mit größtem Betrag der integrierten Nettoladung. Es müssen mindestens zwölf Stunden davor und 24 Stunden danach verfügbar sein. Ein 48-Stunden-Baseline-Lauf ist nötig, um die Störung nicht gegen einen unterschiedlich langen nominalen Versuch zu vergleichen.

\textbf{Vertiefung:} Der letzte gültige Wert kann während des Ausfalls physikalisch falsch werden. Es geht nicht bloß um fehlende Bildpunkte. Abgeleitete Features müssen konsistent aus den eingefrorenen Messungen entstehen. \textbf{Beleg:} Abbildung 6.11. DD hat eine negative mittlere globale Änderung, deren Intervall null einschließt. Das belegt keine heilsame Wirkung von Ausfällen.

\question{F058. Warum verbessert perfekterer SOH nicht automatisch DD?}
\textbf{Antwort:} HDM benutzt SOH in einem expliziten Kapazitätsnenner, HECM zusätzlich zur Tabellenwahl. DD benutzt SOH innerhalb einer gelernten mehrdimensionalen Abbildung. Ein Austausch eines Features kann dessen Zusammenhang mit anderen Features oder mit dem Training verändern.

\textbf{Beleg:} Der Referenz-SOH-Austausch senkt den nominalen MAE für HDM um 0,0370 und für HECM um 0,0166. DD verschlechtert sich im Mittel um 0,0049. \textbf{HYPOTHESE:} Ko-Adaption und Featureverteilungsänderung können das erklären. Für einen Nachweis müsste man Trainingseingänge, Kalibrierung und Ablationen untersuchen. Es folgt nicht, dass bessere SOH-Messungen prinzipiell unerwünscht wären.

\question{F059. Was beweist die HECM-Sensitivitätsanalyse?}
\textbf{Antwort:} Sie prüft lokale Veränderungen der Widerstands- und Zeitkonstantentabellen um ±10 Prozent und OCV-Offsets um ±10 mV, jeweils einzeln und über alle betrachteten Störungssubfälle. Bewertet wird, wie sich die Störungspenalty relativ zur zugehörigen nominalen Variante verändert. Die größte absolute Interaktion beträgt 0,006116 MAE.

\textbf{Grenze:} Das ist keine Tabellenunabhängigkeit. Baseline-MAE und Recovery ändern sich. Bei Widerstand -10 Prozent steigt der mittlere Recovery-oder-Zensur-Wert von 1,20 auf 5,10 Stunden, weil eine Zelle nicht recovern kann. Es wurden keine gemeinsamen Parameterverteilungen, keine beliebigen Kennlinienformänderungen und keine globale Unsicherheitsanalyse untersucht. \textbf{Beleg:} Abbildung A.3 und Tabelle A.1.

\question{F060. Warum betrachtet man ΔΔMAE statt nur absoluten MAE?}
\textbf{Antwort:} Eine Lookup-Änderung kann bereits nominal einen anderen Fehler erzeugen. Die Störungswechselwirkung soll davon getrennt werden. Zuerst wird für dieselbe Lookup der gestörte minus nominale MAE gebildet. Davon wird die entsprechende Differenz unter der ursprünglichen Lookup abgezogen.

\textbf{Herleitung:} ΔΔ = (MAEgestört,perturbiert - MAEsauber,perturbiert) - (MAEgestört,nominal - MAEsauber,nominal). Kleine ΔΔ bedeutet ähnliche zusätzliche Störungskosten, nicht gleichen absoluten Fehler. \textbf{Kritischer Punkt:} In mehreren Subfällen schließen punktweise Intervalle null aus. Daher ist die Aussage gar keine Auswirkung sachlich zu stark, auch wenn der größte Offseteffekt überwiegend durch den Sensorfehler erklärt wird.

\topic{6. Statistik und wissenschaftliche Fairness}
\question{F061. Warum Cell-Macro statt alle Samples zusammenwerfen?}
\textbf{Antwort:} Zuerst werden Fenster und Wiederholungen innerhalb einer Zelle zusammengefasst. Danach bekommt jede der sechs Zellen dasselbe Gewicht. Andernfalls könnten lange Trajektorien, viele verfügbare Fenster oder viele Seeds die Gesamtaussage dominieren.

\textbf{Vertiefung:} Macro beantwortet die Frage nach dem durchschnittlichen Zellverhalten in diesem Holdout-Set. Ein samplegewichteter Wert würde eher einen zufällig gezogenen Zeitpunkt unter der vorhandenen Datendauer repräsentieren. Beide sind definierbar, aber nicht austauschbar. Die Auswahl der Gewichtung muss vor dem Ergebnis begründet werden. \textbf{Beleg:} Kapitel 6.2.6.

\question{F062. Was machen 10000 Bootstrap-Wiederholungen?}
\textbf{Antwort:} Aus den vorhandenen unabhängigen Einheiten werden wiederholt Stichproben mit Zurücklegen gezogen und die gewünschte Kennzahl neu berechnet. Hier werden Zellen und darin geschachtelte stochastische Wiederholungen berücksichtigt. Aus der resultierenden Verteilung wird ein Unsicherheitsintervall geschätzt.

\textbf{Grenze:} Zehntausend Wiederholungen stabilisieren die numerische Bootstrap-Schätzung, schaffen aber keine neuen Zellen und beheben keine fehlende Repräsentativität. Mit sechs Zellen bleiben extreme oder bislang unbesetzte Betriebsbedingungen schlecht erfasst. Bei gepaarten Modellen muss die Paarung im Resampling erhalten bleiben. \textbf{OFFEN:} Die genaue Behandlung mehrerer deterministischer SOH-Fenster sollte am Auswertungscode gezeigt werden können.

\question{F063. Warum ist der Signifikanztest bei sechs Zellen so grob?}
\textbf{Antwort:} Ein exakter Sign-Flip-Test mit sechs unabhängigen gepaarten Differenzen hat nur 2 hoch 6, also 64 Vorzeichenkombinationen. Für einen üblichen zweiseitigen extremen Test ist die kleinste erreichbare Wahrscheinlichkeit 2/\allowbreak{}64 = 0,03125, sofern keine degenerierten Bindungen vorliegen.

\textbf{Herleitung:} Bei sechs paarweisen Modellvergleichen beginnt Holm mit einem Vergleich gegen 0,05/\allowbreak{}6 = 0,00833. Schon der kleinste mögliche unadjustierte Wert ist größer. Unter dieser Konstellation kann deshalb kein Test die erste Holm-Stufe bestehen. Das erklärt fehlende korrigierte Signifikanz ohne den praktischen Effekt zu widerlegen. \textbf{Grenze:} Das Ergebnis hängt von der tatsächlichen Testfamilie und Testdefinition ab. \textbf{Beleg:} Kapitel 6.2.6 und 6.4.1.

\question{F064. Warum kann ein Konfidenzintervall null ausschließen, während p nicht signifikant ist?}
\textbf{Antwort:} Bootstrap-Intervalle, diskrete Sign-Flip-Tests und Holm-Korrektur benutzen unterschiedliche Verfahren und gegebenenfalls unterschiedliche Fehlerkontrolle. Ein punktweises unadjustiertes Intervall ist nicht dasselbe wie eine simultane Aussage über viele Tests. Daher ist die Konstellation nicht automatisch widersprüchlich.

\textbf{Vertiefung:} Aussagen sauber trennen: Effektgröße, punktweises Intervall, Testverfahren, Zahl der Vergleiche und korrigierter p-Wert. Nicht ein Verfahren auswählen, nur weil es die gewünschte Aussage liefert. Die Dissertation sollte durchgängig sagen, dass die kleine Zahl unabhängiger Zellen die inferenzielle Stärke begrenzt.

\question{F065. Was ist der Unterschied zwischen Fehlerbalken, Quartilen und Min-Max?}
\textbf{Antwort:} Quartile beschreiben die Verteilung beobachteter Werte, Min-Max deren Extremwerte und ein Konfidenzintervall die Unsicherheit einer geschätzten Kennzahl. Die Fehlerbalken im Robustheitskapitel sind überwiegend hierarchische 95-Prozent-Konfidenzintervalle, nicht automatisch Min-Max oder Interquartilsabstände.

\textbf{Vertiefung:} Boxplots im Embedded-Kapitel beschreiben dagegen zeitliche Fehlerverteilungen. Ein schmales Intervall um einen Mittelwert kann mit einer breiten Verteilung einzelner Fehler koexistieren. \textbf{Übung:} Vor jedem Diagramm zuerst sagen, was Punkt, Balken, Box und Strich jeweils darstellen. Die genaue Whisker-Regel eines Boxplots muss aus Skript oder Legende kommen, nicht aus Gewohnheit angenommen werden.

\question{F066. Sind gleiche Seeds für alle Modelle unfair?}
\textbf{Antwort:} Nein, gleiche Störungsrealisierungen sind für gepaarte Vergleiche sinnvoll. Alle Modelle sehen dann dieselben zufälligen Messfehler und dieselben Ausfallereignisse. Mehrere Seeds prüfen die Abhängigkeit vom konkreten Zufallsverlauf.

\textbf{Grenze:} Störungsseeds sind keine Trainingsseeds. Sie quantifizieren nicht die Variabilität neu trainierter Netze. Ebenso messen wiederholte MCU-Replays primär Ausführungsstreuung und nicht neue Zellgeneralisation. \textbf{Beleg:} Kapitel 6.3 nennt zehn Seeds für Gaußrauschen und fünf für zufällige Verluste, Jitter und Spikes, deterministische Fälle einmal.

\question{F067. Wie kann man 20 Szenarien, 22 Subfälle und 18 Heatmap-Zeilen erklären?}
\textbf{Antwort:} Es sind unterschiedliche Zählebenen. Das Protokoll enthält 20 Definitionen einschließlich Baseline und Initialisierung. Entfernt man diese beiden, bleiben 18 Messstörungs- und Integritätszeilen. Drei Gainstufen und eine Offsetstufe werden jeweils mit beiden Vorzeichen ausgeführt, wodurch vier zusätzliche Subfälle und damit 22 entstehen.

\textbf{Vertiefung:} Seeds, Zellfenster, Modelle und Diagnosezweige multiplizieren die tatsächliche Zahl der Auswertungen weiter. Die Arbeit nennt 7040 Modell-Fenster-Auswertungen für die Hauptkampagne. Ein prüfbarer Run-Manifest ist belastbarer als aus dem Gedächtnis eine vereinfachte Multiplikation zu konstruieren. \textbf{Beleg:} Abbildung A.1 und Kapitel 6.3.

\question{F068. Wie interpretieren Sie Radarplot und Robustheitsscore?}
\textbf{Antwort:} Die Skalen sind relative Zusammenfassungen innerhalb der betrachteten Kandidaten. Familienbalancierte Robustheit verhindert, dass eine Fehlerfamilie allein wegen vieler Severity-Stufen mehr Gewicht erhält. Andere Gewichtungen führen zu anderen Abständen und teilweise zu anderen Reihenfolgen.

\textbf{Grenze:} 0,967 ist keine Zuverlässigkeit von 96,7 Prozent und keine Fehlerfreiheitswahrscheinlichkeit. Min-Max-Normierung kann sich ändern, sobald ein weiterer Kandidat aufgenommen wird. Im getesteten Satz führt DD mehrere Aggregationen an, aber einzelne lokale Schwächen bleiben relevant. \textbf{Beleg:} Abbildung 6.16. Für Sicherheitsentscheidungen zuerst absolute Anforderungen und Szenariowerte prüfen.

\question{F069. Was wäre eine saubere Studie zur Altersabhängigkeit?}
\textbf{Antwort:} Für jede Zelle die gepaarte Störungspenalty in vergleichbaren SOH-Phasen betrachten, dazu Last, Temperatur und Resetabstände dokumentieren. Der Zielkontrast lautet Änderung der Penalty mit SOH innerhalb der Zelle, nicht bloß schlechterer MAE in späteren Daten.

\textbf{Grenze:} Alter und Betriebsbedingungen sind nicht unabhängig randomisiert. Fehlende aged-Zellen und veränderte Lastprofile können Scheineffekte erzeugen. Eine Regression mit Zellstruktur oder ein Mixed-Effects-Modell kann helfen, aber bei sechs Zellen und wenigen Fenstern wäre ein komplexes Modell instabil. Erst deskriptiv transparent berichten, dann gezielte neue Versuche planen.

\question{F070. Wie vermeiden Sie nachträgliches Schönreden eines Diagramms?}
\textbf{Antwort:} Vorzeichen, Fenster, Seeds, Aggregation, Einheiten und Metrik müssen nachvollziehbar sein. Ein unerwartet negativer ΔMAE wird nicht entfernt, sondern als mögliche Kompensation analysiert. Ein schönes Einzelzellbild wird nicht als Beleg für alle Zellen verkauft.

\textbf{Praktischer Ablauf:} Erst Roh- und Baselinekurven paaren, dann Ereigniszeitpunkt und Kausalität prüfen, dann lokale und globale Metriken unterscheiden, zuletzt interpretieren. Wenn sich die Story dadurch ändert, ändert man die Aussage und nicht die Daten. Explorative Analysen dürfen Teil einer kohärenten Arbeit sein, müssen aber nicht fälschlich als vorab registrierte Bestätigung bezeichnet werden.

\topic{7. Neuronale Grundlagen, Pruning und Quantization}
\question{F071. Wie erklären Sie LSTM und GRU an der Tafel?}
\textbf{Antwort:} Ein LSTM besitzt einen Cell State c und einen Hidden State h. Forget-, Input- und Output-Gates steuern Behalten, Schreiben und Ausgeben. Die Kandidateninformation wird typischerweise mit tanh berechnet. Eine GRU hat nur einen Recurrent State und verwendet Reset- und Update-Gate, um alte und neue Information zu mischen.

\textbf{Prüfpunkt:} Update-Gate-Konventionen unterscheiden sich. Im Grundlagenbild der Dissertation gewichtet z den Kandidaten, während beispielsweise PyTorch in seiner Dokumentation z den alten Zustand gewichten lässt. Das kann durch z gegen 1-z mathematisch äquivalent sein, wenn Parameter und Gleichungen konsistent sind. Auch reset-before und reset-after unterscheiden sich. Beim C-Export zählt die konkret trainierte Framework-Semantik, nicht nur der Name GRU. \textbf{Beleg:} Abbildungen 3.3 und 3.4, PyTorch-Quelle im Quellenabschnitt.

\question{F072. Warum helfen Gates gegen Vanishing Gradients, ohne Stabilität zu garantieren?}
\textbf{Antwort:} Der additive Cell-State-Pfad im LSTM erleichtert den Informationstransport über viele Schritte gegenüber einer ausschließlichen wiederholten nichtlinearen Transformation. Gates können steuern, wie viel Information erhalten bleibt. Dennoch können Activations sättigen, Gradienten verschwinden oder wachsen und lange Extrapolationen problematisch sein.

\textbf{Vertiefung:} Gradient Clipping beschränkt große Gradienten beim Training, beweist aber keine BIBO-Stabilität des trainierten Modells. Endlich lange fehlerfreie Replays liefern empirische Evidenz, keinen allgemeinen mathematischen Stabilitätsbeweis. Ein Sigmoid-Ausgang begrenzt SOC, nicht zwingend alle internen Zustände.

\question{F073. Was unterscheidet Training Chunk, Rolling Window und dauerhaften Zustand?}
\textbf{Antwort:} Ein Training Chunk begrenzt den Backpropagation-Horizont. Ein Rolling Window berechnet jede Ausgabe aus einem begrenzten Fenster mit neu initialisiertem Zustand. Dauerhaftes Streaming führt h und gegebenenfalls c von Sample zu Sample weiter. Diese Begriffe beschreiben unterschiedliche Ebenen und dürfen nicht gleichgesetzt werden.

\textbf{Beleg:} Kapitel 6 primär Rolling-GRU mit 2024 Samples, Hardware zusätzlich continuous und periodic reset. Kapitel 7 stateful LSTM-Ausführung mit Training über begrenzte Chunks. \textbf{OFFEN:} Ob Recurrent States im Training zwischen Chunks übernommen und nur vom Gradienten getrennt wurden oder tatsächlich zurückgesetzt wurden, muss die jeweilige Trainingsschleife belegen. Das beeinflusst die Interpretation des Trainings-Deployment-Mismatches.

\question{F074. Warum ist Structured Pruning für den Mikrocontroller attraktiv?}
\textbf{Antwort:} Ganze Hidden Channels und abhängige Matrixdimensionen werden entfernt. Es entstehen kleinere dichte Arrays, die bestehende dichte C-Schleifen unmittelbar verarbeiten können. Gewichte, recurrent state und ein Teil des Rechenaufwands schrumpfen gemeinsam.

\textbf{Grenze:} Unstrukturiertes Nullsetzen verkleinert dichte Arrays zunächst nicht und kann ohne Sparse-Kernel kaum Laufzeit sparen. Sparse-Indexierung verursacht eigenen Speicher- und Steuerungsaufwand. Structured Pruning ist hier eine hardwaregerechte Wahl, aber nicht allgemein genauer als unstrukturiertes. \textbf{Beleg:} Abbildungen 3.5 und 3.6 sowie 7.4.

\question{F075. Was genau wird beim LSTM-Pruning entfernt?}
\textbf{Antwort:} Entfernt wird ein ganzer Hidden Channel, also eine interne Dimension des LSTM. Dazu gehören jeweils eine Komponente des Hidden State h und des Cell State c sowie die zugehörigen Gewichte und Bias-Einträge. Auch die Verbindungen zu den übrigen Channels und zum nachfolgenden MLP müssen entfernt werden, damit das verkleinerte Netz zusammenpasst.

\textbf{Auswahl:} Für jeden Channel berechnet ihr einen L2-Score. Er fasst die Größe der zugehörigen Gewichte in den Zeilen der Input- und Recurrent-Gate-Matrizen zusammen. Channels mit kleinem Score werden als Kandidaten für das Pruning ausgewählt.

\textbf{Technisches Detail:} Beim Entfernen fallen zusätzlich Spalten der Recurrent-Matrizen und der MLP-Eingangsmatrix weg. Diese Spalten müssen mit entfernt werden, weil sie den gestrichenen Channel weiterverwenden würden. Sie werden in dem hier verwendeten Auswahl-Score jedoch nicht als eigener Beitrag mitgezählt.

\textbf{Grenze:} Kleine Gewichte bedeuten nicht automatisch, dass ein Channel für die Vorhersage unwichtig ist. Ein Channel kann Informationen speichern, die erst später gebraucht werden. Die Gewichtsnorm misst diesen Einfluss auf das Gedächtnis und spätere Vorhersagen nicht direkt.

\textbf{Beispiel:} Ein Channel könnte Information über eine frühere Belastung speichern, die eine spätere SOC-Schätzung unterstützt. Auch bei einem kleinen Gewichtsscore könnte sein Entfernen deshalb den Fehler erhöhen. Das ist eine Veranschaulichung, kein Nachweis für einen bestimmten Channel eures Netzes.

\textbf{Merksatz:} Wir wählen Channels anhand ihrer Gewichtsgröße aus. Das ist eine Näherung für ihre Wichtigkeit. Nach dem Pruning folgt kurzes Fine-Tuning; anschließend wird die Genauigkeit des verkleinerten Modells überprüft. Beleg: Kapitel 7.4, Abbildungen 7.4 und A.10.

\question{F076. Warum 30 Prozent und warum nicht viel mehr?}
\textbf{Antwort:} Die Studie verwendet einen moderaten Betriebspunkt, SOC 64 auf 45 und SOH 128 auf 90 Hidden Units. Das ist kein experimentell bewiesenes globales Optimum. Ein breiter Sweep müsste für jede Rate Auswahl, Fine-Tuning und unabhängige Bewertung wiederholen.

\textbf{Grenze:} Bei stärkerer Reduktion kann die zeitliche Repräsentation oder die lokale Dynamik leiden. Die ideale Rate hängt von Redundanz, Datensatz, Loss Function und Hardwareengpass ab. Ein einzelner günstiger Betriebspunkt reicht zur Demonstration einer Möglichkeit, nicht zur Behauptung, 30 Prozent seien universell optimal.

\clearpage
\question{F077. Warum spart 30 Prozent Breite etwa 45 Prozent MACs?}
\textbf{Kurzantwort:} Beim Pruning werden nicht nur Channels entfernt, sondern auch ihre Verbindungen untereinander. Die rekurrenten Matrizen verlieren dadurch Zeilen und Spalten. Deshalb sinkt dieser Rechenanteil stärker als die Hidden Size. Gemeint sind hier rund 45 Prozent, nicht 75 Prozent.

\textbf{Was ist ein MAC?} MAC steht für Multiply-Accumulate: Einen Wert mit einem Gewicht multiplizieren und das Produkt zu einer laufenden Summe addieren. Bei einer dichten Matrix-Vektor-Multiplikation benötigt jedes Matrixelement einen solchen Schritt.

\textbf{1. Zahlenbeispiel für den quadratischen Anteil:} Stell dir eine rekurrente Matrix mit 10 Channels vor. Jeder neue Channel kann Informationen aus allen 10 bisherigen Channels erhalten. Nach 30 Prozent Pruning bleiben 7 Channels.

\begin{center}\begin{tabular}{lrr}\toprule & Vorher & Nach Pruning \\ \midrule Hidden Channels & 10 & 7 \\ Matrixgröße & $10\times10$ & $7\times7$ \\ MACs für diese Matrix pro Zeitschritt & 100 & 49 \\ \bottomrule\end{tabular}\end{center}
Von 100 MACs bleiben 49 übrig. Für diese Matrix sparst du also 51 MACs beziehungsweise 51 Prozent. Das LSTM hat vier Gates; der gemeinsame Faktor vier ändert dieses Verhältnis nicht.

\textbf{2. Allgemeine Formel:} H bezeichnet die ursprüngliche Hidden Size. Bei genau 30 Prozent weniger Channels beträgt die neue Breite 0,7 H. Für den quadratischen Rechenanteil gilt:

\[ \frac{(0{,}7H)^2}{H^2}=0{,}49 \qquad\Rightarrow\qquad 1-0{,}49=0{,}51=51\,\%. \]
\textbf{3. Warum im gesamten Modell nur etwa 45 Prozent?} Nicht alle Berechnungen schrumpfen quadratisch. Bei den Input-Matrizen bleibt etwa die Zahl der Input Features gleich. Auch der nachfolgende MLP hat andere Skalierungsanteile. Außerdem werden ganzzahlige Channel-Zahlen verwendet: SOC 64 auf 45, SOH 128 auf 90. Die genaue Ersparnis ergibt sich daher aus der MAC-Zählung des gesamten Modells.

\begin{center}\begin{tabular}{lrrr}\toprule Modell & MACs vorher & MACs nachher & Einsparung \\ \midrule SOC & 22.080 & 12.124 & 45,1\,\% \\ SOH & 85.120 & 46.208 & 45,7\,\% \\ \bottomrule\end{tabular}\end{center}
\textbf{Rechnung für das SOC-Modell:}

\[ \frac{22\,080-12\,124}{22\,080}\cdot100\,\%=45{,}1\,\%. \]
\textbf{Merksatz für die Prüfung:} Weniger Channels bedeuten auch weniger Verbindungen zwischen den Channels. Deshalb sparen 30 Prozent weniger Breite bei unseren Modellen etwa 45 Prozent MACs.

\textbf{Grenze und Beleg:} Weniger MACs bedeuten nicht automatisch denselben prozentualen Laufzeitgewinn. Bias, Activations, Speicherbewegungen und Schleifen sind in diesen MAC-Zahlen nicht enthalten. Beleg: Kapitel 7 und die MAC-Zählung der Arbeit.

\clearpage
\question{F078. Warum kann ein gepruntes Modell genauer werden?}
\textbf{Antwort:} Pruning und Fine-Tuning verändern die Funktion. Eine solche Änderung kann redundante oder ungünstige Abhängigkeiten reduzieren und auf einem bestimmten Testverlauf näher an das Ziel führen. In Kapitel 7 ist der SOC-MAE für Pruned niedriger als für Base.

\textbf{Grenze:} Eine Regularisierungswirkung ist eine mögliche Erklärung, keine isoliert bewiesene Ursache. Unterschiedliches Fine-Tuning-Budget kann ebenfalls wirken. Zur Attribution braucht man eine weitertrainierte ungeprunte Kontrollgruppe, mehrere Initialisierungen und mehrere Zellen. SOH verschlechtert sich bei Kompression im primären Vergleich, was eine allgemeine Verbesserung widerlegt.

\question{F079. Wie funktioniert die verwendete INT8-Quantization?}
\textbf{Antwort:} Für jede Matrixzeile wird ein Maßstab aus ihrem maximalen Absolutgewicht geteilt durch 127 berechnet. Die Gewichte werden durch diesen Maßstab geteilt, auf ganze Zahlen gerundet und im symmetrischen Bereich -127 bis 127 gespeichert. Im Kernel wird der effektive Gewichtswert mit dem zeilenspezifischen FP32-Maßstab rekonstruiert.

\textbf{Herleitung:} Bei endlichen Werten und der beschriebenen Max-Skalierung beträgt der Rundungsfehler eines Gewichts höchstens eine halbe Schrittweite. Das ist keine Schranke für den finalen SOC-Fehler, da viele Operationen und rekurrente Rückkopplungen folgen. \textbf{Beleg:} Kapitel 7.5. Ein all-zero-row-Schutz verhindert einen Maßstab von null.

\question{F080. Warum werden -128 und volle Integer-Arithmetik nicht genutzt?}
\textbf{Antwort:} Der genau symmetrische Bereich -127 bis 127 besitzt gleich große positive und negative Endpunkte und stellt null exakt dar. Der zusätzliche negative INT8-Code bleibt ungenutzt. Vollständige Integer-Arithmetik wäre eine andere Implementierung mit quantisierten Activations, Akkumulatoren, Requantization und geeigneten Kernen.

\textbf{Beleg:} In dieser Arbeit sind nur die recurrent matrices Wih und Whh INT8 gespeichert. Bias, Zeilenskalen, Zustände, Activations, MLP und Rechenpfad bleiben FP32. Das Modell ist daher weight-only mixed precision, nicht vollständig INT8. Diese Präzisionsgrenze muss bei jedem Speicher- und Laufzeitargument genannt werden.

\clearpage
\question{F081. Warum spart Quantization nicht 75 Prozent des gesamten Flash?}
\textbf{Kurzantwort:} Unser Modell besteht aus einem LSTM und einem nachgeschalteten MLP. Nur die LSTM-Matrizen Input-to-Hidden (Wih) und Hidden-to-Hidden (Whh) werden als INT8 gespeichert. Deren einzelne Gewichte brauchen ein Byte statt vier Byte. Die MLP-Gewichte bleiben FP32.

\textbf{Was bleibt zusätzlich im Flash?} Programmcode, Bibliotheken, MLP-Gewichte, Bias-Werte und die FP32-Skalierungsfaktoren je quantisierter Matrixzeile. Deshalb gelten 75 Prozent Einsparung nur für die quantisierten Matrixelemente. Für RAM lässt sich diese Rechnung nicht übernehmen: Hidden State, Cell State und Arbeitsdaten bleiben in dieser Implementierung FP32.

\textbf{Warum wurde der MLP nicht ebenfalls quantisiert?} Eine sachliche Begründung für den gewählten Umfang ist, dass der größte Matrixgewichtsblock im LSTM liegt. Der MLP-Output-Head bleibt als unveränderte Ausgabestufe in FP32. Das ist eine nachvollziehbare Eingrenzung der Untersuchung, aber kein Nachweis, dass MLP-Quantization die Genauigkeit verschlechtern würde oder diese Auswahl optimal ist.

\begin{center}\small\begin{tabular}{lrrr}\toprule Modell & LSTM-Matrixgewichte & MLP-Matrixgewichte & LSTM-Anteil\\\midrule SOC & 17\,920 & 4\,160 & 81,2\,\%\\ SOH & 68\,608 & 16\,512 & 80,6\,\%\\\bottomrule\end{tabular}\end{center}
\textbf{Rechnung:} LSTM: 4H(D + H), MLP: HM + M, jeweils ohne Bias und zusätzliche Skalierungswerte. Mit D = 6, H = M = 64 beim SOC sowie H = M = 128 beim SOH liegen rund 80 Prozent der Matrixgewichte im LSTM. Eine zusätzliche MLP-Quantization wäre ein weiterer Optimierungsschritt und müsste separat auf Genauigkeit, Speicher und Laufzeit geprüft werden.

\textbf{Gemessener Flash:} SOC 105,32 auf 52,48 KB; SOH 335 auf 138 KB. \textbf{Offen:} Die KB/\allowbreak{}KiB-Konvention des ursprünglichen Messskripts ist für exakte Bytevergleiche zu klären. kB bedeutet 1.000 Byte, KiB bedeutet 1.024 Byte.

\textbf{Für die Verteidigung:} Wir untersuchen gezielt Weight-only Quantization der LSTM-Matrizen. Sie betreffen den größten Gewichtsblock. Der MLP-Head bleibt FP32; deshalb erwarten wir keine 75 Prozent Einsparung für die gesamte Firmware.

\clearpage
\question{F082. Warum ist Ihr quantisiertes Netz langsamer?}
\textbf{Kurzantwort:} Die Gewichte sind platzsparend als INT8 gespeichert, die Rechnung läuft aber weiterhin in FP32. Der Kernel konvertiert die Gewichte und berücksichtigt die gespeicherten Skalen während der Berechnung. Die Modell-MAC-Zahl bleibt gleich, zusätzliche Arbeit kommt hinzu. Die Skalen werden dabei verwendet, nicht bei jedem Schritt neu bestimmt.

\textbf{Beleg und Grenze:} Die SOC-Kernel Time steigt von 1,40 auf 6,99 ms, die SOH-Kernel Time von 22,73 auf 29,21 ms. Laut Arbeit ist die Compileroptimierung ausgeschaltet. Wie viel Konversion, Skalierung, Speicherzugriffe oder Schleifen jeweils beitragen, wurde nicht einzeln gemessen. Auch die Messgrenzen sind zu beachten (F083). Diese Ergebnisse belegen keine allgemeine Verlangsamung durch Quantization.

\textbf{Was ist ein Integer Kernel?} Ein Kernel ist eine Rechenroutine, etwa für ein Dot Product. Ein Integer Kernel verarbeitet dabei Ganzzahlen. Ein FP32-Kernel arbeitet dagegen mit Gleitkommazahlen. Der frühere Ausdruck Implementierungsraum meint lediglich eine andere technische Umsetzung der Berechnung.

\textbf{Warum nicht jede Zwischenzahl INT8 sein kann:} INT8 reicht von -128 bis 127. Zwei Eingabewerte können hineinpassen, ihr Produkt aber nicht:

\[100\cdot100=10\,000,\qquad 10\,000>127.\]
\textbf{Accumulator:} Produkte und ihre Summe benötigen einen größeren Datentyp, typischerweise INT32. Am Ende passt der Kernel die Skala an, rundet und begrenzt auf den vorgesehenen Ausgangsbereich (Requantization). Größere Integer-Zwischenwerte machen daraus keine FP32-Rechnung.

\textbf{Was übernimmt die Bibliothek?} CMSIS-NN stellt optimierte Rechenroutinen für Arm-Cortex-M-Prozessoren bereit. Passende Routinen übernehmen unter anderem Integer-Multiplikation, Accumulation und Requantization. Das exportierte Modell muss jedoch zum unterstützten Operator, Zahlenformat und Skalierungsschema passen. Nur INT8-Gewichte zu speichern stellt die übrige Verarbeitung nicht automatisch um.

\textbf{Besonderheit LSTM:} Der Cell State trägt Informationen über viele Zeitschritte weiter. Zustände können deshalb größere Integer-Datentypen benötigen. Auch Sigmoid und Tanh brauchen passende Integer- oder Fixed-Point-Verfahren. Ob eine Bibliothek die konkrete LSTM-Variante unterstützt, muss geprüft werden. Diese Arbeit hat eine solche vollständige Integer-Umsetzung nicht benchmarked.

\clearpage
{\large\bfseries\color{accent}F082 -- Vertiefung: Mikrocontroller, NPU und Quantization}\par
\textbf{Kann unser STM32 Integer rechnen?} Ja. Der STM32H753ZI besitzt einen Cortex-M7 mit Integer-/\allowbreak{}DSP-Befehlen und einer Floating-Point Unit. Er ist nicht auf FP32 beschränkt, besitzt aber keine dedizierte NPU. Welche Variante schneller ist, hängt vom Kernel, Compiler, Speicherzugriff und den unterstützten Befehlen ab. Aus unseren Zeiten folgt nicht, dass Integer-Inferenz auf dieser Hardware langsamer sein muss.

\textbf{Was ist eine NPU?} Neural Processing Unit bezeichnet spezialisierte Hardware für neuronale Netze. Sie führt viele MACs parallel aus und nutzt Gewichte und Zwischenwerte möglichst mehrfach, um Speichertransfers zu sparen. Bei gleichem Modell kann die mathematische MAC-Zahl gleich bleiben, obwohl Laufzeit und Energiebedarf sinken. Weniger MACs durch Pruning und schnellere Ausführung derselben MACs durch Hardware sind verschiedene Effekte.

\textbf{Was hat das mit Quantization zu tun?} Quantization stellt Zahlen mit geringerer Genauigkeit dar. INT8-Werte brauchen weniger Speicher und können kompakte, parallele Recheneinheiten ermöglichen. Eine NPU ist dagegen die Hardware, die unterstützte Netzoperationen ausführt. Nicht alle Netze sind quantisiert und nicht jede NPU ist auf INT8 beschränkt; die unterstützten Formate sind gerätespezifisch.

\textbf{Konkrete Beispiele:} STM32N6 besitzt den Neural-ART-Beschleuniger von ST für unterstützte quantisierte Netze. Arm Ethos-U55 ist ein weiterer Beschleuniger mit INT8-/\allowbreak{}INT16-Unterstützung. Solche Hardware ist eine mögliche weitere Untersuchungsplattform, keine garantierte Beschleunigung unseres LSTM. Operator-Unterstützung, recurrent states und Übergaben zwischen CPU und NPU müssen zum Modell passen.

\textbf{Wer bereitet das Netz vor?} Die Softwarewerkzeuge exportieren beziehungsweise konvertieren das Modell und ordnen unterstützte Operationen dem Beschleuniger zu. Andere Operationen können auf der CPU bleiben, sofern die Toolchain dies unterstützt. Die NPU selbst macht aus unserem FP32-LSTM nicht automatisch ein korrekt quantisiertes Modell.

\textbf{Für die Verteidigung:} Unsere Variante spart Gewichtsspeicher, rechnet jedoch weiter in FP32. Eine optimierte Integer-Implementierung oder eine NPU wäre eine andere Ausführung, deren Genauigkeit und Laufzeit separat zu prüfen wären.

\textbf{Merksatz:} NPU = spezialisierte Rechenhardware. Quantization = kompaktere Zahlendarstellung. CMSIS-NN = Softwarebibliothek mit optimierten Rechenroutinen.

{\small\textbf{Quellen zur Einordnung (24.09.2026):} \href{https://github.com/ARM-software/CMSIS-NN}{Arm CMSIS-NN}; \href{https://www.st.com/en/microcontrollers-microprocessors/stm32h753zi.html}{ST: STM32H753ZI}; \href{https://www.st.com/en/development-tools/stm32n6-ai.html}{ST: STM32N6-AI}; \href{https://support.arm.com/compute-ip/ethos-u55}{Arm: Ethos-U55}. Eigene Messwerte und Architektur: Dissertation, Kapitel 7.}\par
\clearpage
\question{F083. Warum unterscheiden sich SOC- und SOH-Laufzeitfaktoren so stark?}
\textbf{Antwort:} MAC-Zahl ist nur ein statischer Teil der Kosten. Activations, Indexierung, Schleifen, Speicherzugriffe und Funktionsgrenzen haben andere Skalierungen. Außerdem sind die Zeitmessgrenzen nicht vollständig gleich, etwa bezüglich SOH-Skalierung im Quantized-Pfad.

\textbf{Grenze:} Es gibt in der Arbeit keine operation-level-Zerlegung, mit der der Faktor 4,99 bei SOC gegenüber 1,29 bei SOH eindeutig erklärt werden könnte. Die ehrliche Antwort benennt gemessenen Unterschied und plausible Beiträge, ohne einen speziellen Cacheeffekt als bewiesene Ursache auszugeben. Ein nächster Versuch würde identische Messgrenzen, Optimierungsflags und Kernelprofiling verwenden.

\topic{8. Embedded-Messung, Filterung und Systemgrenzen}
\question{F084. Warum darf man die GRU- und LSTM-Hardwarezahlen nicht direkt vergleichen?}
\textbf{Antwort:} Es sind unterschiedliche Modelle, Eingangsvektoren, Messgrenzen und Firmwarepfade. Kapitel 6 misst isolierte SOC-Kerne mit extern bereitgestelltem kausalem SOH. Kapitel 7 misst eigenständige stateful SOC- und SOH-LSTM-Modelle mit vorbereiteten Featurevektoren.

\textbf{Beleg:} Continuous-DD in Kapitel 6 benötigt etwa 0,425 ms und 4,1 KiB Peak-RAM. Das SOC-Base-LSTM in Kapitel 7 benötigt 1,40 ms Kernel Time und 4,93 KB RAM. Daraus lässt sich kein sauberer Architektursieger GRU gegen LSTM ableiten. Dafür müssten Features, Genauigkeitsziel, Toolchain und Rechenmodus kontrolliert werden.

\question{F085. Was zeigen 724 ms gegenüber 0,425 ms?}
\textbf{Antwort:} Beim Rolling-Verfahren werden für jede Ausgabe 2024 Schritte mit neuem Zustand berechnet. Continuous verarbeitet nur den neuen Schritt. Der Quotient liegt bei rund 1704, nicht exakt 2024, weil unterschiedliche feste Kosten und Ausführungsdetails beitragen.

\textbf{Grenze:} Es ist kein zusätzlicher Pruninggewinn und keine algebraisch identische Beschleunigung desselben zeitlichen Algorithmus. Das Gedächtnis der kontinuierlichen Ausführung kann länger zurückreichen. Die ähnlichen nominalen Fehler der sechs Replays validieren den untersuchten Betrieb, aber nicht automatisch die identische Robustheit unter allen Störungen. \textbf{Beleg:} Abbildung 6.20.

\question{F086. Warum sind periodische Resets trotz ähnlicher Laufzeit schlecht?}
\textbf{Antwort:} Ein Reset verwirft die akkumulierte Hidden Representation. Das Netz muss Kontext neu aufbauen und kann währenddessen große Fehler erzeugen. Das ist nicht dasselbe wie das rollende Neuberechnen eines vollständigen historischen Fensters.

\textbf{Beleg:} Der Mittelwert der zellweisen maximalen Dataset-SOC-Fehler liegt bei periodischem Reset um 29 Prozentpunkte, einzelne Zellen erreichen etwa 39. Continuous und Rolling haben hier viel günstigere Verläufe. \textbf{Nachfrage:} Was bei Stromausfall tun? State Checkpointing, definierte Kaltstartprozedur, verifizierte Anker und ein Fallback sind mögliche Erweiterungen, aber ihre sichere Funktion ist nicht durch diese Studie bereits nachgewiesen.

\clearpage
\question{F087. Wie wurden Flash und RAM bestimmt?}
\textbf{Kurzantwort wie auf Folie 36:} Flash wird aus dem fertig gelinkten Programm bestimmt. RAM umfasst Variablen, Puffer und den maximal beobachteten Stack. Die folgenden Bilanzen verwenden dieselben Begriffe wie die Präsentation.

\textbf{Flash-Bilanz der Folie:}

\[ F_{\mathrm{Firmware}}=F_{\mathrm{Gewichte}}+F_{\mathrm{Code}}+F_{\mathrm{Konstanten}}. \]
Gewichte sind technisch ebenfalls Konstanten; hier werden sie als eigene Gruppe gezählt. Der letzte Term meint deshalb die übrigen Konstanten und weiteren Flash-Inhalte, einschließlich gespeicherter Startwerte für initialisierte RAM-Variablen. Gemessen werden die belegten Flash-Bereiche des Firmware-Images. Code und Konstanten dürfen nicht doppelt gezählt werden; manche Size-Werkzeuge fassen sie bereits unter text zusammen.

\textbf{RAM-Bilanz der Folie für die Embedded-Auswertung:}

\[ R_{\mathrm{gesamt}}=R_{\mathrm{Variablen}}+R_{\mathrm{Puffer}}+R_{\mathrm{Stack,max}}. \]
\textbf{Was bedeuten die drei Teile?}

\noindent\textbf{\(R_{\mathrm{Variablen}}\):} RAM für die übrigen statisch gespeicherten Werte, zum Beispiel Zähler oder Statusvariablen. Die separat gezählten Puffer sind hier ausgenommen.\par\smallskip
\noindent\textbf{\(R_{\mathrm{Puffer}}\):} RAM für statisch angelegte Datenfelder, zum Beispiel Eingaben, Hidden State, Cell State und Zwischenergebnisse. Puffer sind im Code ebenfalls Variablen, werden in dieser Bilanz aber als eigene Gruppe gezählt.\par\smallskip
\noindent\textbf{\(R_{\mathrm{Stack,max}}\):} Größter während des Tests beobachteter Stack-Bedarf. Der Stack enthält unter anderem lokale Variablen und Informationen zu laufenden Funktionsaufrufen; seine Belegung kann sich während der Ausführung ändern.\par\smallskip
\textbf{Verbindung zur technischen Messung in Kapitel 7:} Die statischen Variablen und Puffer liegen in den Speicherbereichen .data und .bss. .data enthält initialisierte, .bss mit null initialisierte statische Daten. Daher gilt für die beschriebene Auswertung:

\[ R_{\mathrm{Variablen}}+R_{\mathrm{Puffer}}=|\mathrm{.data}|+|\mathrm{.bss}|. \]
Puffer sind darin bereits enthalten und werden nicht nochmals addiert. Die statischen Größen werden aus dem Firmware-Build bestimmt, der Stack-Peak während der Ausführung beobachtet.

\textbf{Rechenbeispiel, erfunden:} 2 KiB Variablen + 6 KiB Puffer + 3 KiB maximaler Stack ergeben 11 KiB RAM. Die ersten beiden Teile ergeben zusammen 8 KiB statischen RAM. Ein KiB entspricht 1.024 Byte.

\textbf{Abgrenzung zu Kapitel 6:} Dort wird zusätzlich dynamisch angeforderter Speicher (Heap) berücksichtigt:

\[ R_{\mathrm{reported,Kap.6}}=R_{\mathrm{static}}+H_{\max,\mathrm{obs}}+S_{\max,\mathrm{obs}}. \]
\textbf{Grenzen:} Beobachtete Maxima sind kein Beweis für den maximal möglichen Bedarf aller Ausführungspfade. Einzelne Peaks müssen nicht gleichzeitig auftreten. Unbenutzte Reserven zählen nicht als tatsächliche Belegung; bereits gezählte Puffer dürfen nicht doppelt eingehen. Quellen: Präsentation, Folie 36; Kapitel 6.2.4 (S. 90) und 7.6.2 (S. 137). Die Folienbilanzen sind keine nummerierten Gleichungen der Dissertation.

\clearpage
\textbf{Ergänzung zu F087: Einheiten für RAM und Datenübertragung}

Die Tabelle unterscheidet zwei Dinge: Bit gegenüber Byte sowie dezimale gegenüber binären Vielfachen. Für RAM sind sowohl kB als auch KiB möglich; entscheidend ist, dass die Einheit zur verwendeten Umrechnung passt.

\begin{center}\renewcommand{\arraystretch}{1.5}\begin{tabular}{|p{2.1cm}|p{2.0cm}|p{3.2cm}|p{7.0cm}|}\hline \textbf{Einheit}&\textbf{Abkürzung}&\textbf{Bedeutung}&\textbf{Typische Verwendung}\\\hline Bit&bit (auch b)&Eine 0 oder 1&Datenübertragung\\\hline Byte&B&8 Bit&Speicherbedarf, Dateigrößen\\\hline Kilobit&kbit (auch kb)&1.000 Bit&Übertragungsraten, etwa kbit/s\\\hline Kilobyte&kB&1.000 Byte&Dateigrößen und dezimale Speicherangaben\\\hline Kibibyte&KiB&1.024 Byte&RAM und Speicherangaben in Zweierpotenzen\\\hline\end{tabular}\end{center}
\textbf{Merksatz:} Großes B bedeutet Byte. kB und KiB sind beide Byte-Einheiten, keine Bit-Einheiten. Die Angabe kbit/\allowbreak{}s bezeichnet dagegen eine Datenmenge pro Sekunde, also eine Übertragungsrate.

\textbf{Warum KiB bei RAM?} Speicher wird über binäre Adressen angesprochen und ist häufig in Größen organisiert, die Zweierpotenzen entsprechen. KiB passt zu diesen Größen und liefert dafür glatte Zahlen. Es ist aber keine Pflicht: kB wäre ebenfalls zulässig. Die Einheit ändert nicht den tatsächlichen Speicherbedarf.

\textbf{Umrechnungsbeispiel:} 4,1 KiB entsprechen 4,1984 kB, gerundet 4,2 kB. Hier wird ausschließlich dieselbe Datenmenge in einer anderen Einheit angegeben.

\clearpage
\question{F088. Warum ändern Zellen den Flash kaum, aber die Laufzeit möglicherweise schon?}
\textbf{Antwort:} Bei identischer Firmware bleiben Code, Gewichte und statische Puffer unabhängig von den Eingangssequenzen gleich. Laufzeit kann dagegen durch datenabhängige Zweige, Activation Functions, Cachezustand oder Interrupts variieren. Wiederholungen prüfen diese Variation.

\textbf{Grenze:} Laufzeitverteilungen sind keine verschiedenen Modellgrößen. Umgekehrt kann eine andere Eingangsfolge einen höheren Stack- oder Heap-Peak auslösen, wenn datenabhängige Pfade existieren. Deswegen bleibt der Messumfang wichtig. In deterministischen dichten Kernen ist geringe Streuung durchaus plausibel und kein Beweis für fehlende Zellvielfalt.

\question{F089. Was misst der DWT-Cycle-Counter, was die Host Latency?}
\textbf{Antwort:} Der Cycle-Counter zählt CPU-Zyklen zwischen instrumentierten Grenzen auf dem Mikrocontroller. Teilt man durch den tatsächlichen Prozessortakt, erhält man die verstrichene Zeit dieses Bereichs. Host Latency enthält zusätzlich UART-Transfer, Protokoll, Scheduling und Rücktransport.

\textbf{Vertiefung:} Interrupts innerhalb der Messgrenzen können mitzählen. Messoverhead, Taktkonfiguration, Cache/\allowbreak{}Warm-up und Zählerüberlauf müssen dokumentiert werden. Bei 480 MHz entsprechen 204000 Zyklen 0,425 ms. \textbf{Beleg:} Kapitel 6.2.4 und 7.6.2. Eine hohe Host Latency widerlegt nicht automatisch einen schnellen Kern.

\question{F090. Haben Sie Echtzeitfähigkeit bewiesen?}
\textbf{Antwort:} Die untersuchten isolierten Ausführungen bleiben innerhalb der 1-s-Abtastperiode. Das zeigt Machbarkeit unter dem Messaufbau. Ein formaler Worst-Case-Execution-Time-Nachweis oder die volle Terminierbarkeit aller parallelen BMS-Aufgaben ist damit nicht erbracht.

\textbf{Vertiefung:} Für ein Gesamtsystem müssen Sensorerfassung, Featurebildung, Kommunikation, SOH-Service, Schutzaufgaben, Interrupts und Fehlerbehandlung hinzugerechnet werden. Besonders Rolling-DD belegt mit 724 ms bereits etwa 72,4 Prozent der Periode. Continuous-DD hat mehr Reserve. \textbf{Nicht sagen:} Ein schneller Mittelwert allein garantiere jede Deadline im Feld.

\question{F091. Warum keine unabhängige Energieaussage?}
\textbf{Antwort:} Eest wird aus angenommener konstanter Boardleistung 0,5 W mal gemessener Kernel Time berechnet. Der Wert besitzt damit exakt dieselbe relative Rangfolge wie die Zeit. Es wurde für die Varianten keine unabhängige Leistungsaufnahme gemessen.

\textbf{Herleitung:} 0,5 W mal 0,8 ms sind 0,4 mJ. Bei 6,99 ms sind es 3,495 mJ. \textbf{Grenze:} Wirkliche Energie hängt von Aktivität, Speicher, Peripherie, Spannungsreglern und Schlafphasen ab. Für eine reale Messung wären synchronisierte Strom-/\allowbreak{}Spannungserfassung, definierte Systemgrenzen und Wiederholungen erforderlich. Der Proxy ist nützlich, darf aber nicht als zusätzlicher unabhängiger Erfolg verkauft werden.

\question{F092. Welche SOH-Filterung wird verwendet und warum ist sie kritisch?}
\textbf{Antwort:} Zuerst wird die Rohvorhersage auf den anfänglichen SOH ausgerichtet. Eine erste Stufe begrenzt Änderungen und mittelt mit α = 0,02. Die zweite Stufe verwendet α = 0,000001 und begrenzt zusätzlich Abwärtsänderungen. Bei 1 Hz hat der lineare EMA-Anteil dieser zweiten Stufe eine Zeitkonstante von etwa 11,57 Tagen und eine 90-Prozent-Antwortzeit von 26,65 Tagen.

\textbf{Grenze:} Eine ruhige Kurve kann stark verzögert sein. Kleine Fehler gegenüber einer langsam interpolierten Referenz beweisen nicht, dass das Rohnetz den SOH schnell erkennt. Rate-Limiter sind nichtlinear und besitzen nicht einfach dieselbe Grenzfrequenz wie ein EMA. \textbf{Beleg:} Kapitel 7.2, Abbildung A.6.

\question{F093. Ist die initiale SOH-Ausrichtung kausal?}
\textbf{Antwort:} Die Formel verwendet den initialen Referenzwert y0. Das ist zeitlich kausal, wenn dieser Startwert vor Einsatz aus einem tatsächlichen Kapazitätstest oder verifizierten Inbetriebnahmewert bekannt ist. Es ist aber eine zusätzliche Informationsannahme, kein kostenlos verfügbarer Online-Messkanal.

\textbf{Grenze:} Bei unbekanntem gebrauchten Akku oder Neustart ohne Historie wäre y0 nicht selbstverständlich bekannt. Die Sensitivität gegenüber falschem Start-SOH müsste separat untersucht werden. Eine genaue Antwort sagt deshalb initial kalibriert statt vollständig blind geschätzt. Außerdem sollte der Nenner der Ausrichtung auf problematisch kleine Anfangsausgaben geprüft werden.

\clearpage
\question{F094. Warum unterscheiden sich SOH-Zahlen zwischen Dashboard und Filteranalyse?}
\textbf{Kurzantwort für die Verteidigung:} Die Zahlen stammen aus zwei getrennten Auswertungen. Die ursprünglichen Benchmarkwerte lassen sich aus den gespeicherten Vorhersagen nachrechnen und gehören zur ersten Filterstufe. Für die spätere Filteranalyse wurden die Vorhersagen neu berechnet und anschließend die Filterstufen verglichen. Diese Wiederholungsrechnung reproduziert den ursprünglichen Verlauf nicht exakt. Deshalb dürfen beide Zahlenreihen nicht als eine gemeinsame Vorher-nachher-Kette interpretiert werden.

\textbf{Welche Werte gehören zusammen?} Alle Angaben sind SOH-MAE in Prozentpunkten; gerundet auf zwei Nachkommastellen.

\begin{center}\renewcommand{\arraystretch}{1.3}\begin{tabular}{|p{8.5cm}|r|r|r|}\hline \textbf{Auswertung und Verarbeitungsstand}&\textbf{Base}&\textbf{Pruned}&\textbf{Quantized}\\\hline Ursprünglicher Benchmark: gespeichert nach Stufe 1&0,85&1,46&1,41\\\hline Wiederholungsrechnung: vor zeitlicher Filterung, bereits am Startwert ausgerichtet&1,85&1,98&2,11\\\hline Wiederholungsrechnung: nach Stufe 1&1,68&1,64&1,93\\\hline Wiederholungsrechnung: nach Stufe 1 und 2&1,48&1,69&1,03\\\hline\end{tabular}\end{center}
\textbf{Was ist geklärt?} Die ursprünglichen Werte 0,85 /\allowbreak{} 1,46 /\allowbreak{} 1,41 sind nicht als Ergebnisse beider Filterstufen zu behandeln: Sie stammen aus den nach der ersten Stufe gespeicherten Benchmarkverläufen. Der Vergleich dieser Werte mit der letzten Tabellenzeile vermischt also sowohl unterschiedliche Läufe als auch unterschiedliche Filterstände.

\textbf{Was bleibt ungeklärt?} Auch beim Vergleich derselben Filterstufe stimmen die Läufe nicht überein: Beim Base-Modell stehen nach Stufe 1 ursprünglich 0,85 gegenüber später 1,68 Prozentpunkten. Die zusätzliche zweite Stufe erklärt diese Abweichung nicht. Unterschiede in numerischer Verarbeitung und Datenaufbereitung sind mögliche Ursachen; ihr konkreter Beitrag wurde nicht isoliert nachgewiesen. Die Abweichung darf deshalb nicht einfach als Rundungsfehler oder als nachgewiesener Filtereffekt erklärt werden.

\textbf{Was zeigt die Filteranalyse trotzdem?} Innerhalb der Wiederholungsrechnung wird der Einfluss der Filter auf dieselben neu berechneten Vorhersagen verglichen. Beim Base-Modell ergibt sich die Kette 1,85 zu 1,68 zu 1,48. Außerdem kann die Filterung die Rangfolge der Varianten verändern. Daraus lässt sich aber nicht der ursprüngliche Benchmarkwert 0,85 ableiten.

\textbf{Bei kritischer Nachfrage:} Die gespeicherten Benchmarkwerte sind nachrechenbar. Die spätere Rechnung ist keine exakte Reproduktion dieses Laufs. Diese Trennung hätte klarer gekennzeichnet werden müssen; die Ursache der verbleibenden Abweichung ist noch offen. Die Zahlen werden nicht nachträglich so ersetzt, dass sie scheinbar übereinstimmen.

\clearpage
\question{F095. Was beweist die Langhorizontanalyse zur Quantization?}
\textbf{Antwort:} Die Trajektorie wird in zehn aufeinanderfolgende Teile geteilt, ohne recurrent state zurückzusetzen. Beim SOC steigen sowohl Base- als auch Quantized-Zielfehler. Gleichzeitig sinkt ihre mittlere gegenseitige Abweichung von 0,332 auf 0,271 Prozentpunkte. Das spricht gegen wachsende quantization-bedingte Abweichung in diesem Replay.

\textbf{Grenze:} Das beweist keine generelle rekurrente Stabilität über beliebige Laufzeiten. Gemeinsame Zielabweichung kann durch schwierigere späte Betriebsbedingungen entstehen. Die Pruned-zu-Base-Differenz wächst, obwohl Pruned insgesamt kleineres MAE hat. Nähe zum Base ist daher nicht identisch mit Genauigkeit. \textbf{Beleg:} Abbildung A.5.

\question{F096. Warum ist Report-Loss nicht dasselbe wie Input-Dropout?}
\textbf{Antwort:} Im Embedded-Report-Loss-Test läuft die Inferenz im Hintergrund weiter. Nur der ausgegebene Bericht fehlt, und die letzte gültige Ausgabe wird gehalten. Im Messstream-Dropout aus Kapitel 6 fehlen dem Estimator dagegen neue Eingangsinformationen. Die internen Zustände können sich dadurch falsch entwickeln.

\textbf{Beleg:} Kapitel 7.6.1 und 7.10. Bei langsamem SOH verändert gehaltene Ausgabe wenig, bei dynamischem SOC kann ein langer Report Loss relevante Abweichungen erzeugen. \textbf{Nicht sagen:} Kaum Fehler bei 10 Prozent Report Loss belege Robustheit des rekurrenten Modells gegen 10 Prozent fehlende Sensorwerte.

\question{F097. Was zeigt der Bitflip-Versuch und was nicht?}
\textbf{Antwort:} In einer einzelnen FP32-Eingangszahl wird das höchstwertige Fraction-Bit 22 verändert. Vorzeichen und Exponent bleiben gleich. Die gestörte und saubere Modellkopie starten aus demselben Zustand und werden anschließend mit denselben 60 sauberen Samples weitergeführt. Damit wird die Nachwirkung eines begrenzten Eingangsbufferfehlers geprüft.

\textbf{Grenze:} Es werden keine Gewichts-, Zustands-, Exponenten- oder Hardware-Strahlungsfehler untersucht. Die Recovery-Schwelle ist ereignisrelativ, maximal aus zehn Prozent des Peaks und einem sehr kleinen Mindestwert gebildet. Daher ist sie nicht identisch mit der Zwei-Prozentpunkte-Recovery aus Kapitel 6. \textbf{Beleg:} Abbildung A.8 und Tabelle A.7.

\question{F098. Was kann der Utility-Score und was bedeuten 1771 Gewichtungen?}
\textbf{Antwort:} U ist eine gewichtete Summe von Verhältnissen zum Base für MAE, Flash, RAM und Energieproxy. Niedriger ist günstiger, Base ist eins. Die 1771 Kombinationen entstehen aus vier nichtnegativen Gewichten in Schritten von 0,05 mit Summe eins.

\textbf{Herleitung:} 20 Gittereinheiten werden auf vier Gewichte verteilt. Die Anzahl ist der Binomialkoeffizient 23 über 3 = 1771. \textbf{Grenze:} 94,64 Prozent gewonnene Gewichtungen sind keine Erfolgswahrscheinlichkeit in realen Anwendungen. Das Ergebnis hängt vom Gitter, den Kriterien und der Base-Normierung ab. Der Energieproxy bringt nur Timinginformation ein. \textbf{Beleg:} Abbildung A.7.

\question{F099. Warum dual-core und warum ist das noch keine vollständige Isolation?}
\textbf{Antwort:} Das Plattformkonzept legt klassische Überwachung und Schutzaufgaben auf den Cortex-M4 und die schwereren Schätzer auf den Cortex-M7. Dadurch lassen sich Algorithmen austauschen, ohne jede BMS-Funktion neu zu strukturieren. Gemeinsame Busse, Speicher und Kommunikation können trotzdem Interferenzen erzeugen.

\textbf{Grenze:} Eine Architekturzeichnung beweist weder Freedom from Interference noch eine geprüfte Safety-Partition. Benötigt werden definierte Schnittstellen, Zeitbudgets, Watchdogs, Datenkonsistenz, Fehlerpfade und Systemtests. \textbf{Beleg:} Kapitel 4.3.3. Die Custom-Plattform verwendet STM32H755, die isolierten Benchmarks STM32H753. Diese Hardwareebenen getrennt benennen.

\question{F100. Kann das System direkt ein Feld-BMS ersetzen?}
\textbf{Antwort:} Die Arbeit zeigt eine Laborplattform und reproduzierbare Embedded-Schätzung. Sie ersetzt keine Produktqualifikation. Feldbetrieb bringt weitere Zellstreuung, Temperaturgradienten, EMV, Sensoralterung, Versorgungsausfälle und Packkopplung mit sich.

\textbf{Vertiefung:} Zell-SOC ist nicht automatisch Pack-SOC. In einem Serienpack teilen Zellen den Strom, besitzen aber unterschiedliche Kapazität und Grenzen. Die schwächste Zelle kann nutzbare Energie begrenzen. Ein BMS braucht unabhängige Schutzfunktionen auch dann, wenn die ML-Schätzung plausibel wirkt. \textbf{Beleg:} Kapitel 4.3 und 8.5. Konkrete Normkonformität wird hier weder behauptet noch geprüft.

\topic{9. Kritische Punkte vor der Verteidigung}
Diese Liste enthält keine automatischen Änderungen an der Dissertation. Sie unterscheidet beobachtete Darstellungsprobleme von offenen Implementierungsfragen. Ziel ist eine ehrliche Antwort, nicht das Einüben einer unbelegten Rechtfertigung.

\question{P01. Abbildung 4.4: Split-Zuordnung ist im vorhandenen PDF nicht erkennbar}
\textbf{Beobachtung:} Die Legende auf PDF-Seite 85 enthält Lade-C-Rate, Entlade-C-Rate und DOD. Zellnamen und explizite Training-/\allowbreak{}Validation-/\allowbreak{}Testzuordnung fehlen in der sichtbaren Grafik. Die Bildunterschrift behauptet dagegen eine Farbunterscheidung dieser drei Splits.

\textbf{Antwort auf eine Rückfrage:} Die verlässliche Zuordnung steht in Tabelle 4.3. Aus dieser Grafik allein kann ich den Split nicht eindeutig zeigen. \textbf{To-do:} Vor endgültiger Verteidigungsversion Bild und Caption abgleichen. Eine Backup-Folie mit Zellen und Splits erstellen. Hier wurde das Original nicht verändert.

\question{P02. Kapazitätsnenner zwischen Label und HDM}
\textbf{Beobachtung:} Relatives SOH wird auf Cref,0 normiert, HDM verwendet in der gedruckten Gleichung Cnom mal SOH. \textbf{To-do:} Den tatsächlich verwendeten SOH-Scaler, die Trainingstargets und die Kapazitätskonstante im Runner gemeinsam zeigen. Falls es eine bewusste unterschiedliche Normierung gibt, Umrechnung dokumentieren. Nicht ohne Implementierungsprüfung einen Fehler behaupten, aber auch nicht Konsistenz voraussetzen.

\question{P03. EFC-Konvention und Faktor zwei}
\textbf{Beobachtung:} Kapitel 6 schreibt EFC = gesamte absolute Ladung geteilt durch Cnom. Bei vollständig mitgezähltem Laden und Entladen ergibt ein kompletter Zyklus damit zwei. Häufig wird ein vollständiger Lade-Entlade-Durchlauf als ein EFC über Division durch 2C definiert.

\textbf{Antwort:} Die verwendete Durchsatzkoordinate kann als Feature konsistent sein, benötigt aber eine explizite Konvention. \textbf{To-do:} Prüfen, ob nur eine Richtung oder tatsächlich beide Richtungen gezählt werden. Datensatz-EFC, Feature-EFC und Achsen der NMC-Grafiken nicht ungeprüft gleichsetzen.

\question{P04. Dimensionen und Vorzeichen der Integrationsgleichungen}
\textbf{Beobachtung:} Wenn Zeit in Sekunden und Kapazität in Ah angegeben werden, ist ein Faktor 3600 erforderlich. In der HECM-Übergangsgleichung steht ηΔt/\allowbreak{}Cb ohne expliziten Umrechnungsfaktor. Die allgemeine SOC-Integralgleichung hat eine andere Vorzeichenkonvention als das ladungspositive Qc.

\textbf{Antwort:} Beides kann korrekt sein, wenn Cb in Coulomb beziehungsweise die Zeit in Stunden und die Stromkonvention passend definiert sind. \textbf{To-do:} Einheiten und Codekoordinaten offenlegen. Die Vorbereitung liefert keine Bestätigung der Firmwareeinheiten.

\question{P05. MLP-Komplexität und Hyperparameterwahl}
\textbf{Beobachtung:} Strukturell einfach ist nicht gleich klein. Das finale MLP ist breit und tief. Die Lag-Matrix zeigt exemplarische Testzellen, weshalb die Trennung von Auswahl und Demonstration nachvollziehbar bleiben muss. \textbf{To-do:} model.summary beziehungsweise Parameterzahl, Suchprotokoll und unberührte Testentscheidung bereithalten. Den vierten NMC-Testfall ausdrücklich nennen können.

\question{P06. Filterprovenienz und Anfangskalibrierung}
\textbf{Beobachtung:} Haupt-SOH-Auswertung und Filteranalyse haben verschiedene Werte und Rangfolgen. Initialalignment verwendet den Anfangsreferenzwert. \textbf{To-do:} Für beide Exporte Zell-ID, Samplebereich, Filterparameter, initiale Zustände, Masken und MAE-Berechnung nebeneinander dokumentieren. Nicht aus dem glatten Bild auf ungefilterte Onlinegenauigkeit schließen.

\question{P07. Robustheit einer Semantik ist nicht automatisch Robustheit der anderen}
\textbf{Beobachtung:} Der große Robustheitsvergleich verwendet Rolling-GRU. Die schnelle Hardwareempfehlung nutzt Continuous-State. Ähnliche nominale Ergebnisse über sechs Zellen schließen Unterschiede unter Störungen nicht aus. \textbf{To-do:} Als Grenze nennen. Eine vollständige Übertragung würde die relevanten Störungs- und Recoveryfälle im Continuous-Modus wiederholen.

\question{P08. Hardwaremessgrenzen und Compilerflags}
\textbf{Beobachtung:} Kapitel 7 misst handgeschriebene, nicht optimiert kompilierte Referenzkerne. Bei SOH liegt die Skalierung nicht für alle Varianten innerhalb derselben Zeitgrenze. \textbf{To-do:} Compiler, Flags, Takt, Warm-up, Interruptbedingungen, gemessene Funktion und Hostprotokoll auf einer Backup-Folie festhalten. Die Ergebnisse nicht als maximale Leistungsfähigkeit des STM32 oder von INT8 ausgeben.

\question{P09. Tatsächliche Systemintegration versus Konzept}
\textbf{Beobachtung:} Die Plattformbeschreibung umfasst Schutz, Sensing und Dual-Core-Trennung. Die berichteten Kernbenchmarks schließen zahlreiche Systemfunktionen aus. \textbf{To-do:} Genau trennen, was auf der eigenen Platine tatsächlich getestet wurde, was auf dem Nucleo gemessen wurde und was eine Architekturabsicht ist. Dafür Messprotokolle statt nur Blockdiagramme bereithalten.

\question{P10. Keine Tabellenunabhängigkeit, keine Sicherheit aus einem Score}
\textbf{Beobachtung:} Die HECM-Analyse zeigt lokale Robustheitsstabilität in vielen Fällen, aber eine relevante Recovery-Änderung. Radar- und Utility-Scores hängen von Normierung und Gewichtung ab. \textbf{To-do:} Die Formulierungen unabhängig, sicherstes Modell und universell robust vermeiden, wenn sie nicht durch ein definiertes absolutes Sicherheitskriterium gedeckt sind.

\topic{10. Rechenübungen mit Lösungen}
\question{R01. Stromoffset und Zeit}
\textbf{Aufgabe:} Eine 1,8-Ah-Zelle wird mit einem konstanten Messoffset von +50 mA integriert. Wie groß ist der rein rechnerische SOC-Fehler nach einer, sechs und 24 Stunden ohne Anker oder Clipping?

\textbf{Lösung:} 0,05 A mal t /\allowbreak{} 1,8 Ah ergibt 0,02778, 0,16667 und 0,66667. Das sind 2,78, 16,67 und 66,67 Prozentpunkte. Das Vorzeichen hängt von der Stromkonvention ab. Diese Werte sind Endabweichungen, nicht automatisch zeitliche MAE. Bei linear von null wachsendem Fehler wäre der mittlere absolute Fehler über den Zeitraum halb so groß.

\question{R02. Gainfehler bei Lade-Entlade-Wechsel}
\textbf{Aufgabe:} Eine Zelle erfährt zuerst +1 Ah und anschließend -1 Ah Nettoladung. Der Strom hat +3 Prozent Gainfehler. Was passiert mit dem idealisierten Integrationsfehler?

\textbf{Lösung:} Nach der ersten Phase sind +0,03 Ah Zusatzfehler vorhanden. Nach beiden Phasen ist die Summe null. Zwischenzeitlich bestand trotzdem ein Fehler. Bei 1,8 Ah entspricht der Peak 1,67 Prozentpunkten. Reale Asymmetrien, Resetbedingungen, Effizienz und Kapazitätsänderungen verhindern eine automatische vollständige Kompensation.

\question{R03. MAE und RMSE}
\textbf{Aufgabe:} Normierte Fehler sind 0,01, -0,01 und 0,10. Berechne MAE, RMSE und Bias.

\textbf{Lösung:} MAE = 0,12/\allowbreak{}3 = 0,04. MSE = 0,0102/\allowbreak{}3 = 0,0034. RMSE = 0,05831. Bias = 0,10/\allowbreak{}3 = 0,03333. In Prozentpunkten sind dies 4,00, 5,83 und 3,33. RMSE betont den großen Ausreißer stärker. Der positive Bias zeigt die gerichtete Abweichung, obwohl ein negativer Einzelfehler vorkommt.

\question{R04. Makro- und Mikrogewichtung}
\textbf{Aufgabe:} Zelle A hat MAE 0,02 über 100 Samples, Zelle B MAE 0,08 über 900 Samples. Was ergibt sich?

\textbf{Lösung:} Cell-Macro = (0,02 + 0,08)/\allowbreak{}2 = 0,05. Samplegewichtetes Mittel = (100 mal 0,02 + 900 mal 0,08)/\allowbreak{}1000 = 0,074. Keines ist allein durch Mathematik richtiger. Die Forschungsfrage bestimmt, ob jede Zelle oder jeder Zeitpunkt gleich gewichtet werden soll.

\question{R05. EMA-Zeitkonstante}
\textbf{Aufgabe:} Warum dauert α = 10\textasciicircum{}-6 bei 1 Hz so lange?

\textbf{Lösung:} Nach n Schritten ist der verbleibende Anteil eines Sprungs (1-α)\textasciicircum{}n. Die Zeitkonstante ist -1/\allowbreak{}ln(1-α), näherungsweise eine Million Sekunden beziehungsweise 11,57 Tage. Für 90 Prozent Antwort setzt man den Rest auf 0,1. Daraus folgen ln(0,1)/\allowbreak{}ln(1-α), etwa 2,303 Millionen Sekunden oder 26,65 Tage. Nichtlineare Rate-Limiter können das reale Verhalten zusätzlich verändern.

\question{R06. Quantization einer Zeile}
\textbf{Aufgabe:} Die Gewichtzeile lautet -0,8, 0,1, 0,4. Bestimme Skala, Codes und rekonstruierten mittleren Wert.

\textbf{Lösung:} s = 0,8/\allowbreak{}127 = 0,0062992. Gerundet ergeben sich -127, 16 und 64, wobei exakte Halbwerte von der verwendeten Rundungsregel abhängen können. Für 0,1 erhält man 16s = 0,100787. Der Fehler 0,000787 liegt unter s/\allowbreak{}2 = 0,003150. Die Rekonstruktion von 0,4 kann an der Halbwertgrenze je nach Regel geringfügig anders ausfallen. Framework- und C-Rundung müssen zusammenpassen.

\question{R07. SOC-MACs vor und nach Pruning}
\textbf{Aufgabe:} Nutze D = 6, M = 64 und N = 4H(D+H) + HM + M für H = 64 und H = 45.

\textbf{Lösung:} Base: 4 mal 64 mal 70 + 64 mal 64 + 64 = 22080. Pruned: 4 mal 45 mal 51 + 45 mal 64 + 64 = 12124. Die Reduktion ist 9956/\allowbreak{}22080 = 45,09 Prozent. Es werden 19 von 64 recurrent channels entfernt, also 29,69 Prozent. Der MLP-Kopf behält M = 64.

\question{R08. Speicher des rollenden Eingabefensters}
\textbf{Aufgabe:} Wie viel Speicher benötigen 2024 Samples mit acht FP32-Features allein als Rohpuffer?

\textbf{Lösung:} 2024 mal 8 mal 4 = 64768 Byte = 63,25 KiB. Das erklärt einen großen Teil des berichteten Rolling-Peaks von 67,3 KiB. Restlicher Speicher entfällt auf Zustände, temporäre Daten und Stack. Continuous braucht diesen vollständigen Rohfensterpuffer nicht für dieselbe rekurrente Schrittausführung.

\question{R09. Zensur verändert den Mittelwert}
\textbf{Aufgabe:} Fünf Zellen recovern je nach 1,2 h. Eine bleibt bis 24 h außerhalb des Bandes. Wie sieht ein vereinfachtes Recovery-or-censor-Mittel aus?

\textbf{Lösung:} (5 mal 1,2 + 24)/\allowbreak{}6 = 5 h. Ein einzelner zensierter Fall kann damit einen ungefähr vierfachen Mittelwert erzeugen, obwohl die anderen fünf fast unverändert sind. Das erklärt qualitativ die Empfindlichkeit des HECM-Resistance-Falls. Die Rechnung ist ein Lehrbeispiel, nicht die genaue Rekonstruktion der Fensteraggregation des Experiments.

\question{R10. Energieproxy und Auslastung}
\textbf{Aufgabe:} Rechne 0,5 W bei 1,40 ms und 0,80 ms. Wie viel Prozent einer 1-s-Periode benötigt 724 ms?

\textbf{Lösung:} 0,70 mJ und 0,40 mJ. Der relative Proxy-Rückgang beträgt 42,86 Prozent und ist identisch zum Zeitrückgang. 724 ms entsprechen 72,4 Prozent der Periode. Bei zusätzlicher BMS-Last kann die verbleibende Reserve knapp werden, während ein Median keine Worst-Case-Garantie bietet.

\question{R11. Gewichtsraum des Utility-Scores}
\textbf{Aufgabe:} Warum sind es 1771 Gewichtskombinationen und nicht 21 hoch 4?

\textbf{Lösung:} Jedes Gewicht kann zwar 21 Werte von 0 bis 1 annehmen, aber die vier Gewichte müssen sich zu eins addieren. Gesucht sind nichtnegative ganzzahlige Lösungen von a+b+c+d = 20. Stars-and-bars liefert (23 mal 22 mal 21)/\allowbreak{}(3 mal 2 mal 1) = 1771. Gleiches Rastergewicht ist keine empirische Häufigkeit realer Nutzerpräferenzen.

\question{R12. Was muss man aus 0,006116 gegenüber 0,1763 schließen?}
\textbf{Aufgabe:} Setze maximale Lookup-Störungsinteraktion und nominale HECM-Offsetpenalty ins Verhältnis.

\textbf{Lösung:} Der Quotient beträgt ungefähr 0,0347, also 3,47 Prozent. Das ordnet die lokale Interaktion gegenüber dem großen Offseteffekt ein. Es ist keine globale Sensitivitätskonstante. Die Maxima und signed Subfälle müssen in ihrer Definition beachtet werden, und Recovery kann trotzdem stark reagieren. Ein kleines Verhältnis rechtfertigt deshalb nicht die Aussage, die Lookup habe keine Auswirkung.

\topic{11. Lern- und Rechercheplan}
\question{Priorität A: Ohne Nachschlagen erklären können}
\noindent\textbullet\enspace Die drei Forschungsfragen, je ein Ergebnis und je eine Grenze. Dabei MLP/\allowbreak{}NMC, GRU/\allowbreak{}LFP und LSTM/\allowbreak{}LFP sauber trennen.

\noindent\textbullet\enspace SOC/\allowbreak{}SOH-Definitionen, Vorzeichen, Ah-zu-Coulomb-Umrechnung, Offset- und Gainfehler einschließlich Rechenbeispiel.

\noindent\textbullet\enspace Unterschied von nominalem MAE, gepaarter Störungspenalty und gepaarter Recovery. Zwei Prozentpunkte sind nicht zwei Prozent relativer Fehler.

\noindent\textbullet\enspace Cell-Macro, sechs unabhängige Zellen, Bootstrap, Zensierung und warum 10000 Wiederholungen keine 10000 unabhängigen Experimente sind.

\noindent\textbullet\enspace Weight-only INT8, nicht Full-Integer. Gemessene Laufzeit, gelinkter Flash, beobachteter RAM-Peak und Energieproxy klar auseinanderhalten.

\noindent\textbullet\enspace Die vollständige SOH-Filterkette mit initialer Information und 26,65 Tagen EMA-Antwortzeit.

\question{Priorität B: Mit einer Skizze herleiten können}
\noindent\textbullet\enspace 2RC-Modell, EKF-Prädiktion, Spannungsresiduum und lokale Beobachtbarkeit auf dem LFP-Plateau.

\noindent\textbullet\enspace LSTM- und GRU-Gates einschließlich Framework-Konventionen. Zeige, wo der Zustand gespeichert wird.

\noindent\textbullet\enspace Gate-konsistentes Pruning, Matrixdimensionen und quadratische MAC-Skalierung.

\noindent\textbullet\enspace Zeilenweise Quantization, Halbschritt-Fehlergrenze und warum sie keine Ausgangsfehlerschranke ist.

\noindent\textbullet\enspace Gepaarten Kontrast, Lookup-Interaktion und unterschiedliche Unsicherheitsdarstellungen.

\question{Priorität C: Vor der Verteidigung anhand von Unterlagen klären}
\noindent\textbullet\enspace Exakte NMC-Testzellen, Split- und Skalierungslogik, Historienbins und Rolle der MAE-Matrix bei der Auswahl.

\noindent\textbullet\enspace Kapitel-7-Testzelle, genauer Samplebereich und vollständiger Trainingssplit. Nicht automatisch die Kapitel-6-Zellen einsetzen.

\noindent\textbullet\enspace Unterschiedliche SOH-Filterauswertungen, Anfangskalibrierung und verwendete Referenzdenominator.

\noindent\textbullet\enspace EFC-Konvention, HECM-Einheiten, Full-Charge-Schwellen und Reset-Logs.

\noindent\textbullet\enspace Hardwaremessgrenzen, Compilerflags, Takt, Speicherdateien und tatsächlicher Integrationsstand der Custom-Platine.

\question{Vorschlag für zehn Lerntage}
\textbf{Tag 1:} F001 bis F015. Eine dreiminütige Gesamterklärung aufnehmen. SOC-Offset, LFP-Plateau und EKF ohne Unterlagen zeichnen.

\textbf{Tag 2:} F016 bis F025. Die Zell-/\allowbreak{}Split-Tabelle auswendig erklären, Labels von Onlinefeatures trennen. P01 bis P04 mit Originalunterlagen abgleichen.

\textbf{Tag 3:} F026 bis F036. Kapitel 5 anhand aller sieben Abbildungen erklären. Parameterzahl und Lag-Zeiten nachrechnen. Hyperparameterauswahl rekonstruieren.

\textbf{Tag 4:} F037 bis F046. Vier Modelle, Eingänge und Zustandseingriffe erklären. Ein künstliches Recovery-Beispiel mit Rückfall zeichnen.

\textbf{Tag 5:} F047 bis F060. Für jede Störungsfamilie Eingriff, sichtbaren Befund, plausiblen Mechanismus und Grenze in vier Sätzen notieren.

\textbf{Tag 6:} F061 bis F070. Statistik an den Rechenübungen trainieren. Alle Fehlerbalken im Robustheitskapitel korrekt benennen.

\textbf{Tag 7:} F071 bis F083. LSTM/\allowbreak{}GRU zeichnen, Pruning-Matrizen ausschneiden, eine Quantization-Zeile von Hand berechnen.

\textbf{Tag 8:} F084 bis F100. Hardwaremesskette und Filterung erklären. Hauptzahlen auf eine einzige Backup-Seite bringen.

\textbf{Tag 9:} Den vollständigen Abbildungsatlas durchgehen. Pro Abbildung 30 Sekunden Aussage, 30 Sekunden Methodik und 30 Sekunden Einschränkung. Die kritischen Punkte mit offenen Nachweisen nicht überspringen.

\textbf{Tag 10:} Zwei Probeprüfungen zu je 30 bis 45 Minuten. Eine konzentriert sich auf Batterien und Referenzen, die andere auf Statistik und Embedded-Umsetzung. Ungeklärte Antworten als offene Punkte notieren statt improvisieren.

\question{Simulation einer kritischen Prüfung}
\textbf{Prüfpfad A:} Sie nennen niedrigen SOH-MAE. Wie wurde SOH gemessen? Warum steigt er? Nutzt die Interpolation Zukunftsinformation? Welche Informationen bekommt das Modell tatsächlich? Warum ist die Filterantwort länger als viele betriebliche Ereignisse? Welche Aussage bleibt dann vom Netz selbst?

\textbf{Antwortkern:} Kapazität als operative Referenz erklären, offline Label von Onlineinput trennen, Rohnetz und Filter getrennt bewerten, initiale Information offenlegen und keine schnelle Kapazitätserkennung aus glatten Kurven ableiten.

\textbf{Prüfpfad B:} DD gewinnt. Ist der Vergleich fair? Qc enthält doch schon SOC. Wurden alle Features gestört? Warum andere Initialzustände? Warum kein signifikanter Test? Warum nehmen Sie anschließend eine andere Inferenzsemantik auf der Hardware?

\textbf{Antwortkern:} Gleiche kausale Informationsbasis und rekonstruierte Features, aber unterschiedliche Hypothesenräume erläutern. Repräsentantenvergleich statt Universalaussage. Paarung und kleine Zellzahl erklären. Continuous als separat nominal validierten Deployment-Modus behandeln, dessen volle Störungsübertragung nicht bereits bewiesen ist.

\textbf{Prüfpfad C:} Quantization soll effizient sein. Warum ist sie langsamer? Was ist überhaupt quantisiert? Wo sind Energie und RAM gemessen? Wie viel schneller würde optimierter Code sein?

\textbf{Antwortkern:} Weight-only-Speicherung von Recurrent-Matrizen, FP32-Rekonstruktion, unveränderte Topologie und nicht optimierte Referenzkerne erklären. Messgrenzen und Proxy nennen. Keine nicht gemessene Beschleunigung durch einen optimierten Kernel versprechen.

\question{Wie mit einer unbekannten Frage umgehen?}
\textbf{Antwortmuster:} Zuerst den gesicherten Teil nennen. Dann präzise abgrenzen, welcher zusätzliche Nachweis fehlt. Abschließend einen Versuch oder eine Analyse vorschlagen, die die Frage beantworten würde. Beispiel: Die Tabelle zeigt Recovery-Sensitivität bei -10 Prozent Widerstand. Welcher RC-Zweig den Effekt dominiert, habe ich nicht isoliert. Dazu würde ich Ri, R1 und R2 einzeln perturbieren und Innovation, Kalman-Gain und SOC-Korrektur ereignisweise auswerten.

Nicht verwenden: Das muss am Rauschen liegen. Der Plot sieht gut aus. Mehr Daten machen es automatisch besser. Das Framework wird schon stimmen. Gerade transparente Grenzen zeigen wissenschaftliches Verständnis.

\topic{12. Quellen und konkrete Rechercheaufträge}
\textbf{Primärquelle der Ergebnisse:} main.pdf und die zugehörigen TeX-Dateien im Dissertation-Verzeichnis, Stand nach Git-Aktualisierung auf 2d8fdc9. Der Abbildungsatlas enthält eine direkt aus dem PDF erzeugte Seitenzuordnung. figure\_\allowbreak{}inventory.json speichert zusätzlich den SHA-256-Hash der gelesenen PDF. Die folgenden externen Quellen unterstützen Hintergrundlernen, nicht die erneute Verifikation der eigenen Messdaten.

\question{Q1. Rekurrente Gleichungen und Export}
PyTorch GRU-Dokumentation (https:/\allowbreak{}/\allowbreak{}docs.pytorch.org/\allowbreak{}docs/\allowbreak{}main/\allowbreak{}generated/\allowbreak{}torch.nn.GRU.html). Nachschlagen: Update-Gate-Konvention und abweichende Position des Reset-Gates gegenüber ursprünglichen Formulierungen. Lernziel: Eine Framework-GRU nicht nur anhand einer allgemeinen Blockskizze nachimplementieren. Diese offizielle Dokumentation wurde für die Vorbereitung geprüft.

Die LSTM-Grundlagen und Framework-Referenzen stehen außerdem in der Dissertation, insbesondere in Kapitel 3.4 und der Bibliografie unter den LSTM-/\allowbreak{}PyTorch-Referenzen. Lernauftrag: Gate-Reihenfolge und getrennte Biasvektoren im tatsächlich exportierten Modell identifizieren. Bei konkreter Softwareversion deren Dokumentation verwenden, nicht unbesehen eine aktuelle API.

\question{Q2. Stromintegration und Unsicherheit}
Movassagh et al., A Critical Look at Coulomb Counting (https:/\allowbreak{}/\allowbreak{}arxiv.org/\allowbreak{}abs/\allowbreak{}2101.05435). Die Originalarbeit untersucht unter anderem Strommessung, Integrationsnäherung, Kapazitätsunsicherheit und Zeitbasis. Lernauftrag: Eigene Definition von Offset, Gain, Clock-Fehler und Kapazitätsfehler danebenstellen. Die konkrete Fehlerformel dieses Lernskripts selbst herleiten und ihre Voraussetzungen nennen.

\question{Q3. Optimierte Embedded-Kernel}
Arm CMSIS-NN, offizielles Repository (https:/\allowbreak{}/\allowbreak{}github.com/\allowbreak{}ARM-software/\allowbreak{}CMSIS-NN) und offizielle Fully-Connected-API (https:/\allowbreak{}/\allowbreak{}arm-software.github.io/\allowbreak{}CMSIS-NN/\allowbreak{}latest/\allowbreak{}group\_\allowbreak{}\_\allowbreak{}FC.html). Lernauftrag: Gewichts-, Aktivierungs-, Bias- und Akkumulator-Datentypen vergleichen. Herausarbeiten, weshalb ein INT8-Gewichtsarray in einem FP32-C-Kernel noch keine solche Integer-Ausführung ist. Die Quellen wurden zur Einordnung geprüft. Keine daraus abgeleitete Beschleunigungszahl für die eigene Firmware behaupten.

\question{Q4. Literatur bereits in der Dissertation}
\textbf{Plett, Battery Management Systems:} Zustandsdefinition, Coulomb Counting, OCV, Observer und Parametrierung wiederholen. \textbf{Hyperband, Li et al.:} Budgetzuteilung und Early-Stopping-Verfahren lernen. \textbf{Structured Sparsity, Wen et al., sowie recurrent pruning nach Narang/\allowbreak{}Lobacheva:} Strukturauswahl von Deploymentrepräsentation unterscheiden. \textbf{Quantization nach Jacob et al. und Krishnamoorthi:} Post-Training-Verfahren, Quantization-Aware Training und Requantization auseinanderhalten. Die genauen bibliografischen Angaben stehen in bib/\allowbreak{}Dissertation.bib. Diese Werke sind gezielte Leseaufträge, keine Behauptung, sämtliche Volltexte seien für diese Vorbereitung erneut geprüft worden.

\question{Q5. Eigene Belegmappe}
Für jede Ergebnisfamilie einen Belegordner anlegen: Konfiguration, Zellliste, Modellhash, Scaler, Skriptversion, Laufmanifest, Ergebnistabelle und Abbildung. Für Hardware zusätzlich Toolchain, Flags, Linker-Map, Firmwarehash und Zeitmessgrenzen. Für die beiden SOH-Auswertungen außerdem Filterzustand und Labeldefinition. Diese Mappe beantwortet die Rückfrage Woher genau stammt diese Zahl wesentlich besser als weitere schöne Diagramme.

\topic{13. Abbildungsatlas: Lesen, erklären, hinterfragen}
Dieser Atlas folgt den 67 tatsächlich nummerierten Abbildungen der angegebenen PDF. Originalseiten werden ausschließlich als separate Kopien in das Lernskript eingebettet. Die Abbildungen der Dissertation werden nicht verändert. Enthält eine Seite zwei Bilder, ist jeweils die bezeichnete Nummer gemeint. Die Seitenausschnitte zeigen bewusst auch die Caption, damit Unterschiede zwischen Grafik und Text erkennbar bleiben.

\textbf{Antwortschema für jede Abbildung:} Was ist auf den Achsen? Was sind Beobachtungseinheit und Zeitraum? Was wurde verändert und was konstant gehalten? Was zeigt das Ergebnis? Welche Erklärung ist belegt und welche nur plausibel? Was kann diese Grafik nicht beantworten? Bei einer schematischen Abbildung entfällt die empirische Interpretation, ihre Bausteine und Annahmen müssen stattdessen erklärt werden.

\textbf{Nummerierungsfalle:} Beispielsweise ist Figure\_\allowbreak{}09 im Ergebnisordner hier Abbildung 6.9. Die ADC-Grafik heißt im Quellenordner Figure\_\allowbreak{}16, aber im Dissertation-PDF Abbildung 6.14. Im Gespräch immer die Nummer der tatsächlich gezeigten Folie beziehungsweise PDF verwenden.

\topic{Abbildung 2.1: Anforderungsrahmen}
\begin{center}\includegraphics[width=\linewidth,height=0.60\textheight,keepaspectratio]{figures/thesis-page-036.pdf}\end{center}
Originalseite der Dissertation: PDF-Seite 36. Die folgende Interpretation bezieht sich auf die bezeichnete Abbildung.

\textbf{Frage:} Warum steht das BMS im Zentrum und weshalb sind die drei Bereiche gleichberechtigt?

\textbf{Antwort:} Der Rahmen verknüpft Schätzqualität, Hardwaregrenzen und die Umsetzungs-/\allowbreak{}Wartungskette. Die Darstellung ist eine konzeptionelle Ordnung, keine gemessene Gleichgewichtung. Ein fehlender Sicherheits- oder Zeitnachweis lässt sich nicht durch bessere mittlere Accuracy ausgleichen. \textbf{Vertiefung:} Zeige für jeden Bereich eine konkrete Ergebnisabbildung der Dissertation. \textbf{Grenze:} Das Schaubild weist noch keine Vollständigkeit aller Anforderungen eines Serien-BMS nach. Siehe F005 und F006.

\topic{Abbildung 3.1: Schätzerfamilien}
\begin{center}\includegraphics[width=\linewidth,height=0.60\textheight,keepaspectratio]{figures/thesis-page-042.pdf}\end{center}
Originalseite der Dissertation: PDF-Seite 42. Die folgende Interpretation bezieht sich auf die bezeichnete Abbildung.

\textbf{Frage:} Warum sind hybride Modelle keine völlig disjunkte vierte Physik?

\textbf{Antwort:} Hybrid bezeichnet Kombinationen von Integrations-, Modell- und Lernkomponenten. HDM verändert den Kapazitätspfad des Coulomb Counting, HECM verbindet einen Observer mit geschätztem SOH, DD verwendet ebenfalls physikalisch motivierte Features. Die Familien sind konzeptionelle Orientierung. \textbf{Grenze:} Aus den Kästen darf man nicht ableiten, dass alle Verfahren innerhalb einer Familie denselben Informationsgehalt oder dieselbe Güte hätten. Siehe F004 und F037.

\topic{Abbildung 3.2: MLP}
\begin{center}\includegraphics[width=\linewidth,height=0.60\textheight,keepaspectratio]{figures/thesis-page-046.pdf}\end{center}
Originalseite der Dissertation: PDF-Seite 46. Die folgende Interpretation bezieht sich auf die bezeichnete Abbildung.

\textbf{Frage:} Was repräsentiert eine Verbindung und wie berechnet sich die nächste Schicht?

\textbf{Antwort:} Jede Verbindung ist ein Gewicht. Ein Neuron summiert gewichtete Eingänge plus Bias und wendet eine Activation Function an. Die ganze Schicht berechnet a = f(Wx+b). Training verändert W und b, nicht die aktuelle Eingangsmessung. \textbf{Vertiefung:} Für n Eingänge und m Ausgänge gibt es nm Gewichte plus m Biasparameter. \textbf{Grenze:} Die gezeichnete kleine Knotenzahl ist schematisch und keine Größenangabe des finalen SOH-Modells. Siehe F030.

\topic{Abbildung 3.3: LSTM-Zelle}
\begin{center}\includegraphics[width=\linewidth,height=0.60\textheight,keepaspectratio]{figures/thesis-page-048.pdf}\end{center}
Originalseite der Dissertation: PDF-Seite 48. Die folgende Interpretation bezieht sich auf die bezeichnete Abbildung.

\textbf{Frage:} Warum gibt es zwei horizontale Zustandswege?

\textbf{Antwort:} c speichert den Cell State, h die nach außen gegebene rekurrente Repräsentation. Forget-Gate und Input-Gate bestimmen alte und neue Information im Cell State, Output-Gate steuert die sichtbare Ausgabe. Die obere additive Verbindung erleichtert lange Abhängigkeiten. \textbf{Grenze:} Ein offenes Gate garantiert weder physikalische Interpretierbarkeit noch beliebig langes korrektes Gedächtnis. Auf Nachfrage die Dimensionen sämtlicher Gates angeben: jeweils H Komponenten. Siehe F071 und F072.

\topic{Abbildung 3.4: GRU-Zelle}
\begin{center}\includegraphics[width=\linewidth,height=0.60\textheight,keepaspectratio]{figures/thesis-page-050.pdf}\end{center}
Originalseite der Dissertation: PDF-Seite 50. Die folgende Interpretation bezieht sich auf die bezeichnete Abbildung.

\textbf{Frage:} Welcher Zweig wird hier mit z gewichtet und ist das überall gleich?

\textbf{Antwort:} Die Dissertation verwendet die Kandidatengewicht-Konvention. Andere Implementierungen gewichten mit z den bisherigen Zustand. Das ist nur bei konsistenter Umdefinition äquivalent. Reset-before und reset-after beeinflussen ebenfalls die konkrete Rechnung. \textbf{Grenze:} Nicht allein anhand des Diagramms einen Framework-Export bestätigen. Das endgültige C-Modell muss gegen exakt dessen Gateordnung, Biasbehandlung und Activations geprüft werden. Siehe F071.

\topic{Abbildung 3.5: Pruning-Granularität}
\begin{center}\includegraphics[width=\linewidth,height=0.60\textheight,keepaspectratio]{figures/thesis-page-055.pdf}\end{center}
Originalseite der Dissertation: PDF-Seite 55. Die folgende Interpretation bezieht sich auf die bezeichnete Abbildung.

\textbf{Frage:} Warum sparen die roten Nullen im mittleren Bild nicht zwingend Laufzeit?

\textbf{Antwort:} Unstrukturierte Nullen behalten die ursprünglichen Matrixdimensionen. Eine dichte Schleife multipliziert sie weiterhin, sofern kein spezieller Sparse-Pfad verwendet wird. Rechts entfernt Structured Pruning ganze Zeilen und passende Spalten, sodass wirklich kleinere dichte Tensoren entstehen. \textbf{Grenze:} Die Zeichnung illustriert Speicherrepräsentation, keine gemessenen Zeitverhältnisse. Sparsity-Anteil und effektiver Speed-up sind getrennte Größen. Siehe F074.

\topic{Abbildung 3.6: Abhängige Strukturen}
\begin{center}\includegraphics[width=\linewidth,height=0.60\textheight,keepaspectratio]{figures/thesis-page-057.pdf}\end{center}
Originalseite der Dissertation: PDF-Seite 57. Die folgende Interpretation bezieht sich auf die bezeichnete Abbildung.

\textbf{Frage:} Warum ist unabhängig beschnittenes Pruning vor einer Addition ungültig?

\textbf{Antwort:} Addierte Zweige müssen dieselben Dimensionen und die beabsichtigte semantische Zuordnung besitzen. Werden verschiedene Kanäle entfernt, können Dimensionen oder Bedeutungen nicht mehr zusammenpassen. Gemeinsame Auswahlregeln erhalten die Struktur. \textbf{Übertragung:} Im LSTM betrifft derselbe Hidden Channel mehrere Gates, recurrent columns, Zustände und MLP-Eingänge. Alle müssen zusammen angepasst werden. Siehe F075.

\topic{Abbildung 3.7: Quantization Points und Gitter}
\begin{center}\includegraphics[width=\linewidth,height=0.60\textheight,keepaspectratio]{figures/thesis-page-065.pdf}\end{center}
Originalseite der Dissertation: PDF-Seite 65. Die folgende Interpretation bezieht sich auf die bezeichnete Abbildung.

\clearpage
\textbf{Frage:} Was zeigt Abbildung 3.7, und welche Teile haben wir tatsächlich quantisiert?

\textbf{Kurzantwort:} Die Abbildung zeigt allgemeine Möglichkeiten der Quantization. Unsere Umsetzung verwendet Weight-only Quantization: Die beiden LSTM-Gewichtsmatrizen Input-to-Hidden und Hidden-to-Hidden werden als INT8 gespeichert. Beim Rechnen werden die Gewichtswerte in FP32 umgewandelt und ihre Scales berücksichtigt. Die MLP-Gewichte, Biases, States und Activations bleiben FP32. Es wird also nicht das gesamte Netz auf INT8 umgestellt.

\textbf{Welche Gewichte genau?} Sowohl W\_\allowbreak{}ih (Input-to-Hidden) als auch W\_\allowbreak{}hh (Hidden-to-Hidden), jeweils für alle vier LSTM-Transformationen. Der frühere Ausdruck „recurrent weights“ war hier missverständlich: Gemeint sind die LSTM-Gewichtsmatrizen insgesamt, nicht nur W\_\allowbreak{}hh. „Zeilenweise symmetrisch“ bedeutet: Jede Matrixzeile erhält einen eigenen Scale; verwendet werden die Integer-Codes von -127 bis +127 mit Zero Point 0.

\textbf{Warum nicht alle Gewichte?} Die untersuchte Umsetzung konzentriert sich auf den großen LSTM-Gewichtsblock. Der MLP-Head bleibt FP32. Auch seine Gewichte zu quantisieren wäre ein zusätzlicher Optimierungsschritt, dessen Genauigkeit und Laufzeit separat geprüft werden müssten. Siehe F079 bis F082.

\textbf{Was unterscheidet Weights, Activations und Akkumulatoren?} Weights sind gespeicherte Modellparameter. Activations sind Zwischenwerte, die beim Verarbeiten der Eingaben entstehen. Ein Akkumulator sammelt die Produkte einer gewichteten Summe. Diese Summe kann einen größeren Wertebereich benötigen: Schon 100 mal 100 ergibt 10 000 und passt nicht in INT8. Integer Kernels verwenden deshalb häufig INT32 für die Summen. In unserer Umsetzung erfolgt die Rechnung weiterhin in FP32. INT8-Gewichte bedeuten also nicht automatisch INT8-Zwischenwerte oder INT8-Summen.

\textbf{Was ist der Zero Point?} Er ist die gespeicherte Ganzzahl z, die den realen Wert null repräsentiert. Im oberen, symmetrischen Beispiel gilt z = 0: Integer-Code 0 bedeutet realer Wert 0. Bei einer asymmetrischen Abbildung könnte beispielsweise Code 100 den realen Wert 0 darstellen. Es geht nicht bloß um positive und negative Werte; diese gibt es auch bei symmetrischer Quantization.

\textbf{Wie liest man die drei Gitter rechts?} Oben stehen die Integer-Codes q, darunter die zugehörigen realen Werte x. Das obere Gitter ist symmetrisch um null. Das mittlere nutzt UINT8 von 0 bis 255 für nichtnegative Werte, ebenfalls mit Zero Point 0. Das untere nutzt einen verschobenen Zero Point für einen asymmetrischen Wertebereich. „Gitter“ meint die diskreten darstellbaren Zahlenstufen; Werte dazwischen werden gerundet.

\textbf{Was bedeutet Full Integer?} Die Netzoperationen werden mit geeigneten Integer- beziehungsweise Fixed-Point-Verfahren ausgeführt, statt die Gewichte zum Rechnen wieder nach FP32 zu konvertieren. Dabei müssen nicht alle Zwischenwerte INT8 sein: INT32-Summen oder breitere States sind möglich. Dafür braucht man passende Kernels und Scales auch für die Zwischenwerte. Das ist eine andere Umsetzung als unsere Weight-only Quantization. Abbildung 3.7 beschreibt die Grundlagen und belegt keine Full-Integer-Ausführung unserer Firmware.

\topic{Abbildung 4.1: NMC-Alterung über Zyklen und Zeit}
\begin{center}\includegraphics[width=\linewidth,height=0.60\textheight,keepaspectratio]{figures/thesis-page-082.pdf}\end{center}
Originalseite der Dissertation: PDF-Seite 82. Die folgende Interpretation bezieht sich auf die bezeichnete Abbildung.

\textbf{Frage:} Warum verändern sich scheinbare Abstände zwischen den oberen und unteren Kurven?

\textbf{Antwort:} Die obere Achse normiert auf Durchsatzzyklen, die untere auf Kalenderzeit. Zellen mit anderen Raten und DOD erreichen vergleichbaren Durchsatz in anderer Zeit. Daher erscheinen Alterungsunterschiede je nach Koordinate anders. Zwei Zellen pro Szenario zeigen zudem Replikatstreuung. \textbf{Grenze:} Aus den Kurven allein keine separate Aktivierungsenergie oder eindeutige Dominanz eines Alterungsmechanismus schätzen. Die Temperaturwirkung ist sichtbar, andere Faktoren wirken gekoppelt. Siehe F009 und F010.

\topic{Abbildung 4.2: LFP-DoE-Würfel}
\begin{center}\includegraphics[width=\linewidth,height=0.60\textheight,keepaspectratio]{figures/thesis-page-084.pdf}\end{center}
Originalseite der Dissertation: PDF-Seite 84. Die folgende Interpretation bezieht sich auf die bezeichnete Abbildung.

\clearpage
{\large\bfseries\color{accent}Abbildung 4.2 -- DoE kurz erklärt}\par
\textbf{Frage:} Warum verwendet man ein Design of Experiments (DoE), und was zeigt der Würfel?

\textbf{Zweck:} DoE ist statistische Versuchsplanung. Lade-C-Rate, Entlade-C-Rate und DoD werden gezielt kombiniert, um ihre Einflüsse auf die Alterung mit begrenzter Versuchskapazität zu untersuchen. Ein Designpunkt ist eine Kombination dieser drei Einstellungen.

\textbf{Full Factorial Design:} Bei drei Faktoren mit je zwei Stufen (niedrig/\allowbreak{}hoch) werden alle 2 x 2 x 2 = 8 Kombinationen getestet. Sie bilden die acht grünen Würfelecken. Man kann damit Haupteffekte untersuchen, etwa den durchschnittlichen Einfluss einer höheren Lade-C-Rate, und Wechselwirkungen: Wirkt schnelles Laden bei hoher DoD anders als bei niedriger DoD?

\textbf{Central Composite Design (CCD):} Die gezeichnete Anordnung ergänzt die acht Eckpunkte um einen roten Mittelpunkt und sechs blaue Axialpunkte. An jedem Axialpunkt wird ein Faktor weiter nach oben oder unten variiert, während die anderen auf mittlerer Einstellung bleiben. Diese zusätzlichen Punkte ermöglichen eine quadratische Modellierung und damit die Untersuchung von Krümmung, etwa eines überproportional steigenden Kapazitätsverlusts bei hoher Belastung.

\textbf{Nutzen für unsere Modelle:} Die Kampagne liefert Messdaten aus systematisch unterschiedlichen Belastungs- und Alterungsverläufen für Training und Bewertung der SOC-/\allowbreak{}SOH-Schätzer. Wiederholungen derselben Einstellung würden zusätzlich die Zellstreuung erfassen. Ob Einflüsse tatsächlich nachgewiesen sind, ergibt erst die statistische Auswertung, nicht der Würfel allein.

\textbf{Einordnung der Unterlagen:} Die Grafik zeigt 15 unterschiedliche Positionen; der Begleittext nennt dagegen acht Betriebspunkte mit Wiederholungen. Die konkrete Zellzuordnung muss am Versuchsplan abgeglichen werden. Die frühere pauschale Erklärung, die 15 Zellen seien lediglich Wiederholungen von acht Punkten, ist damit nicht belegt. Siehe F025.

\topic{Abbildung 4.3: Struktur eines LFP-Verlaufs}
\begin{center}\includegraphics[width=\linewidth,height=0.60\textheight,keepaspectratio]{figures/thesis-page-084.pdf}\end{center}
Originalseite der Dissertation: PDF-Seite 84. Die folgende Interpretation bezieht sich auf die bezeichnete Abbildung.

\textbf{Frage:} Warum wechseln regelmäßige Zyklen, Diagnoseblöcke und dynamische Abschnitte?

\textbf{Antwort:} Die Kampagne kombiniert beschleunigte zyklische Alterung mit Check-ups und anwendungsbezogenen Lastprofilen. Check-ups liefern vergleichbare Referenzinformation, dynamische Abschnitte beanspruchen die Schätzer anders als konstante Ströme. \textbf{Grenze:} Ein dargestelltes PV-Heimspeicherprofil repräsentiert nicht jede Microgrid-Nutzung. Außerdem muss erläutert werden, welche Diagnosephasen im jeweiligen Modellinput enthalten sind und welche für Labels verwendet werden. Siehe F020.

\topic{Abbildung 4.4: LFP-SOH über die Kampagne}
\begin{center}\includegraphics[width=\linewidth,height=0.60\textheight,keepaspectratio]{figures/thesis-page-085.pdf}\end{center}
Originalseite der Dissertation: PDF-Seite 85. Die folgende Interpretation bezieht sich auf die bezeichnete Abbildung.

\textbf{Frage:} Warum haben einige Kurven starke Einbrüche und lokale Erholung, und welche ist C29?

\textbf{Antwort:} Die Kurven verbinden diskrete Kapazitätsanker über die Versuchszeit. Unterschiedliche Lastbedingungen und Eigenerwärmung gehen mit unterschiedlichen Alterungsverläufen einher. Lokale Anstiege sind gemessene Kapazitätsvariation, nicht automatisch irreversible Regeneration. \textbf{Konkreter offener Punkt:} Im vorhandenen PDF nennt die Legende Betriebsbedingungen, aber keine Zellnamen oder eindeutig ausgewiesenen Splits. C29 darf deshalb nicht allein aus einer vermuteten Farbe identifiziert werden. Für diese Zuordnung Tabelle 4.3 benutzen. Siehe P01.

\topic{Abbildung 4.5: Holdout-Abdeckung}
\begin{center}\includegraphics[width=\linewidth,height=0.60\textheight,keepaspectratio]{figures/thesis-page-087.pdf}\end{center}
Originalseite der Dissertation: PDF-Seite 87. Die folgende Interpretation bezieht sich auf die bezeichnete Abbildung.

\textbf{Frage:} Weshalb bleibt eine Zelle ausschließlich fresh, und kann High eine Statistikgruppe sein?

\textbf{Antwort:} C27 deckt im vorhandenen Zeitbereich nur SOH 0,930 bis 1 ab. Es gibt keine beobachteten mid-life-/\allowbreak{}aged-Fenster dieser Zelle. High besteht allein aus C29 und ist deskriptiv. \textbf{Vertiefung:} Temperaturbereich, P95-C-Rate und SOH sind verschiedene Coverage-Größen. Ein breiter Balken bedeutet nicht viele unabhängige Wiederholungen. \textbf{Grenze:} Keine belastbare zwischenzellige Varianz für High und keine rein kausale Lastklassenwirkung. Siehe F022 und F023.

\topic{Abbildung 4.6: BMS-Platine}
\begin{center}\includegraphics[width=\linewidth,height=0.60\textheight,keepaspectratio]{figures/thesis-page-091.pdf}\end{center}
Originalseite der Dissertation: PDF-Seite 91. Die folgende Interpretation bezieht sich auf die bezeichnete Abbildung.

\textbf{Frage:} Welche sichtbaren Hardwarebereiche sind für Ihre Beiträge entscheidend?

\textbf{Antwort:} Der Text ordnet Sensing/\allowbreak{}Balancing, Strommess- und Schutzpfad sowie Controller/\allowbreak{}Kommunikation als Funktionsdomänen ein. Die physische Trennung von Leistungs- und Signaldomäne unterstützt Messqualität und EMV. \textbf{Grenze:} Bauteilwerte, Kalibriergenauigkeit, reale Leiterplattenbestückung und Teststatus nicht aus dem Rendering erraten. Schaltplan, Stückliste und Messprotokoll bereithalten. Eine Platinenabbildung ist kein Funktions- oder Sicherheitsnachweis. Siehe F099 und P09.

\topic{Abbildung 4.7: BMS-Systemblock}
\begin{center}\includegraphics[width=\linewidth,height=0.60\textheight,keepaspectratio]{figures/thesis-page-092.pdf}\end{center}
Originalseite der Dissertation: PDF-Seite 92. Die folgende Interpretation bezieht sich auf die bezeichnete Abbildung.

\textbf{Frage:} Wo verlaufen Messung, Schätzung und Schutz, und was passiert bei Ausfall der KI?

\textbf{Antwort:} Messfront-end und Stromsensor liefern Daten an die Controllerdomäne. Die konzeptionelle M4/\allowbreak{}M7-Trennung hält konventionelle BMS-Aufgaben von austauschbaren Schätzern getrennt. Eine reale Ausfallantwort benötigt aber Watchdog, sichere Grenzprüfung, Plausibilität und definierte Fallbacks. \textbf{Grenze:} Das Blockbild allein belegt nicht, dass alle diese Maßnahmen fertig implementiert und getestet sind. Die isolierten H753-Benchmarks nicht als vollständigen H755-Systemtest ausgeben. Siehe F099.

\topic{Abbildung 5.1: Korrelationsmatrix}
\begin{center}\includegraphics[width=\linewidth,height=0.60\textheight,keepaspectratio]{figures/thesis-page-099.pdf}\end{center}
Originalseite der Dissertation: PDF-Seite 99. Die folgende Interpretation bezieht sich auf die bezeichnete Abbildung.

\textbf{Frage:} Warum sind negative und positive Zusammenhänge beide rot?

\textbf{Antwort:} Farbe stellt hier den Betrag dar. Die gedruckten Zahlen enthalten Vorzeichen und exakten Koeffizienten. Es ist eine symmetrische Matrix, daher genügt eine Dreieckshälfte. \textbf{Vertiefung:} Durchsatz und SOH haben gemeinsame zeitliche Entwicklung. Eine große Pearson-Korrelation beweist weder eine Ursache noch robuste Generalisierung auf andere Lastprofile. \textbf{Prüfübung:} Eine konkrete Zelle der Matrix mit Vorzeichen, Größenordnung und möglicher Confounding-Variable erklären. Siehe F028.

\topic{Abbildung 5.2: Lag-Konstruktion}
\begin{center}\includegraphics[width=\linewidth,height=0.60\textheight,keepaspectratio]{figures/thesis-page-100.pdf}\end{center}
Originalseite der Dissertation: PDF-Seite 100. Die folgende Interpretation bezieht sich auf die bezeichnete Abbildung.

\clearpage
{\large\bfseries\color{accent}Abbildung 5.2 -- Resampling und Aggregation}\par
\textbf{Frage:} Wie werden die Messdaten für das MLP zeitlich aufbereitet, und welche Information bleibt erhalten?

\textbf{Die beiden Zeitraster auseinanderhalten:} Laut Dissertation werden die zunächst je Prozessschritt unterschiedlich abgetasteten NMC-Messungen auf ein Sekundenraster gebracht. Für die MLP-Historie folgt eine gröbere Zeitauflösung. Der Startpunkt der Untersuchung ist 30 Minuten mit 12 Lags; als beste untersuchte Kombination nennt die Arbeit 10 Minuten mit 16 Lags, also 160 Minuten Historie. Das sind zwei verschiedene Verarbeitungsschritte, keine widersprüchlichen Abtastraten.

\textbf{Resampling: Welches Zeitraster?} Es legt die neuen Zeitintervalle fest. Im gefundenen NMC-Vorverarbeitungsskript wird ein Sekundenraster erzeugt und interpoliert. Resampling erzeugt dabei keine neuen unabhängigen Messungen. Im späteren Lag-Code wird auf die übergebene gröbere Intervalllänge resampled; für die in der Arbeit genannte 10-Minuten-Konfiguration sind das 600 Sekunden pro Intervall.

\textbf{Aggregation: Wie wird ein Intervall zusammengefasst?} Der gefundene Lag-Code verwendet den arithmetischen Mittelwert pro Zeitintervall und Messspalte, anschließend lineare Interpolation fehlender Werte. Auf einem lückenlosen Sekundenraster enthält ein 10-Minuten-Intervall beispielsweise 600 Werte, aus denen ein Mittelwert wird. Es wird also nicht lediglich jeder 600. Messwert ausgewählt. Beispiel: 300 Sekunden mit 3,6 V und 300 Sekunden mit 3,8 V ergeben 3,7 V als Intervallmittel. Derselbe Mittelwert kann aus unterschiedlichen zeitlichen Verläufen entstehen.

\textbf{Was bleibt erhalten, was geht verloren?} Erhalten bleiben mittlere Signalniveaus, die Reihenfolge der Zeitintervalle, ihr fester Abstand und Veränderungen zwischen ihnen. Die genaue Reihenfolge innerhalb eines Intervalls, kurze Spitzen und schnelle Wechsel sind aus dem Mittelwert nicht rekonstruierbar. Interpolation füllt Lücken durch angenommene Zwischenwerte; sie stellt den unbekannten Verlauf nicht wieder her. Das MLP erhält damit eine gröbere zeitliche Darstellung, nicht den ursprünglichen Sekundenverlauf.

\textbf{Besonderheit der kumulierten Ströme:} Im gefundenen Lag-Code wird zuerst der Strom pro Intervall gemittelt. Danach werden positive und negative Intervallmittel getrennt aufsummiert und mit der Intervalllänge in Ah umgerechnet. Diese Summen werden anschließend ebenfalls verzögert als Features abgelegt. Sie tragen Nutzungsinformation von vor dem Lag-Fenster mit. Lade- und Entladeanteile innerhalb desselben Intervalls können sich allerdings schon beim Mitteln aufheben; ihre getrennte Integration vor dem Mitteln wäre nicht dieselbe Rechnung.

\textbf{Warum diese Aufbereitung?} Statt Tausenden Sekundenwerten bekommt das MLP eine kompakte Historie. Das reduziert Eingangsgröße und glättet schnelle Schwankungen. Zu grobe Intervalle verlieren relevante Dynamik; zu feine Intervalle decken bei gleicher Lag-Zahl nur wenig Zeit ab. Deshalb wurden in der Arbeit mehrere Zeitauflösungen und Lag-Längen verglichen.

\clearpage
{\large\bfseries\color{accent}Abbildung 5.2 -- Lag-Fenster und Überlappung}\par
\textbf{Frage:} Wie erhält ein MLP ohne Recurrent State eine Historie, wie überlappen die Fenster, und warum?

\textbf{Konstruktion im gefundenen Code:} Für jede Featuregröße werden der aktuelle Intervallwert und verschobene Kopien angelegt: t, t-1, t-2 bis t-L. Die Schleife läuft einschließlich L; L = 16 bedeutet hier daher 16 vergangene Werte plus den aktuellen, insgesamt 17 Werte je Feature. Die Kapazität des aktuellen Intervalls dient als Zielwert. Im gezeigten Trainingsaufruf werden Kapazitätsspalten aus den Inputs entfernt. Die Arbeit beschreibt das Ziel als aktuellen SOH, nicht als Zukunftsprognose.

\textbf{Konkrete Rechnung mit der in der Arbeit genannten Einstellung:} Bei 10 Minuten Abstand und L = 16 liegen zwischen ältestem und aktuellem Zeitindex 160 Minuten. Für Spannung, Temperatur sowie positive und negative kumulierte Ladung wären es nach dieser Code-Konvention 4 x 17 = 68 Eingabewerte. Das ist eine Ableitung aus der gefundenen Lag-Klasse, keine Bestätigung der Eingangsgröße des final trainierten Modells. 16 Werte einschließlich des aktuellen hätten dagegen nur 15 Abstände, also 150 Minuten zwischen den Zeitindizes. Bei Intervallmitteln muss zusätzlich zwischen Zeitindex-Abstand und der gesamten Breite der enthaltenen Intervalle unterschieden werden.

\textbf{Überlappung bei einem Schritt Vorschub:} Jede nächste Zeile wird ein neues Beispiel; ein zusätzlicher größerer Stride ist in dieser Lag-Klasse nicht eingebaut. Bei 10-Minuten-Auflösung würde das nächste Fenster daher zehn Minuten später beginnen und enden:

\[\begin{aligned}A &: [0,10,20,\ldots,150,160]\ \text{min},\\B &: [10,20,30,\ldots,160,170]\ \text{min}.\end{aligned}\]
\textbf{Was wird wiederverwendet?} Beide Fenster enthalten die 16 Zeitindizes von Minute 10 bis Minute 160. Ein alter Eintrag fällt heraus, ein neuer kommt hinzu. Nach dieser Code-Konvention überlappen somit 16 von 17 Einträgen, etwa 94 Prozent. Die Zeitangaben sind schematische Rasterindizes; aus der linken Beschriftung eines gemittelten Intervalls darf nicht abgeleitet werden, dass sein Mittelwert schon zu Intervallbeginn verfügbar war.

\textbf{Warum überlappen lassen?} Nach jedem neuen Intervall kann eine Schätzung mit aktualisierter Historie entstehen, ohne auf einen komplett neuen, nicht überlappenden Datenblock zu warten. Im Training werden unterschiedliche Endzeitpunkte und Ausschnitte eines Verlaufs nutzbar. Die zeitliche Reihenfolge steckt in der festen Reihenfolge der Input-Spalten. Das MLP selbst behält zwischen zwei Aufrufen keinen Hidden State.

\textbf{Grenze bei Training und Validation:} Überlappende Fenster sind stark abhängig und keine neuen unabhängigen Versuche. Im gefundenen Projekt werden die Lag-Zeilen vor dem Training gemischt; der gezeigte Aufruf reserviert anschließend 20 Prozent für Validation. Dadurch können stark ähnliche Fenster auf beiden Seiten liegen. Ein separater Test auf vollständig zurückgehaltenen Zellen ist davon zu unterscheiden. Eine rein zeitliche Trennung müsste zusätzlich die Fensterüberlappung an der Grenze berücksichtigen.

\textbf{Quellenstatus:} Dissertation, Kapitel 4 (Sekundenraster) und 5 (Zeitauflösung/\allowbreak{}Lags); NMC-Projekt mg\_\allowbreak{}farm\_\allowbreak{}main\_\allowbreak{}project, Klassen LagDataset und LagDataCollection sowie main.py; Vorverarbeitung resample\_\allowbreak{}BX\_\allowbreak{}SOH\_\allowbreak{}Combined\_\allowbreak{}to\_\allowbreak{}second\_\allowbreak{}v2.py. Die gefundene main.py nutzt einen älteren Aufruf mit 30 Minuten und unterschiedlichen Lag-Längen. Der finale 10-Minuten-/\allowbreak{}16-Lag-Trainingslauf ist damit nicht vollständig reproduziert. Die Code-Details sind deshalb ausdrücklich als Befund dieses Projektstands gekennzeichnet. Siehe F026 und F029.

\topic{Abbildung 5.3: Trainings- und Evaluationsworkflow}
\begin{center}\includegraphics[width=\linewidth,height=0.60\textheight,keepaspectratio]{figures/thesis-page-102.pdf}\end{center}
Originalseite der Dissertation: PDF-Seite 102. Die folgende Interpretation bezieht sich auf die bezeichnete Abbildung.

\textbf{Frage:} An welcher Stelle darf das Testset Einfluss nehmen?

\textbf{Antwort:} Auf die finale Auswertung, nicht auf Scalerfitting, Hyperparameterwahl oder frühes Stoppen. Datenaufbereitung mit gelernten Parametern muss innerhalb der Entwicklungsgrenze bleiben. \textbf{Kritische Nachfrage:} Ist das im gezeichneten Ablauf oder in der Implementierung garantiert? Das Schema beschreibt Absicht, die tatsächliche Trennung muss der Code belegen. \textbf{Offen:} Zufällige überlappende Validation-Sequenzen und historischer Gebrauch der Lag-MAE-Matrix prüfen. Siehe F031 und F032.

\topic{Abbildung 5.4: SOH-MLP}
\begin{center}\includegraphics[width=\linewidth,height=0.60\textheight,keepaspectratio]{figures/thesis-page-104.pdf}\end{center}
Originalseite der Dissertation: PDF-Seite 104. Die folgende Interpretation bezieht sich auf die bezeichnete Abbildung.

\textbf{Frage:} Warum sieht das Modell klein aus, während die Tabelle viele Schichten zeigt?

\textbf{Antwort:} Das Bild erklärt Feature Groups und Vollvernetzung schematisch. Die tatsächlich gewählten Breiten stehen in Tabelle 5.1. Die Darstellung ist kein maßstabsgetreues Netz mit der richtigen Neuronenzahl. \textbf{Grenze:} Einfach meint feedforward und nachvollziehbare Featurebildung, nicht zwingend wenige Parameter. Den Unterschied zur Eingangssuche mit acht Neuronen erklären. Die bedingte Parameterrechnung in F030 bereithalten.

\topic{Abbildung 5.5: Historien-MAE-Matrix}
\begin{center}\includegraphics[width=\linewidth,height=0.60\textheight,keepaspectratio]{figures/thesis-page-108.pdf}\end{center}
Originalseite der Dissertation: PDF-Seite 108. Die folgende Interpretation bezieht sich auf die bezeichnete Abbildung.

\textbf{Frage:} Warum liegt das günstige Gebiet zwischen den Extremen?

\textbf{Antwort:} Zu grobe Intervalle können relevante Lade-/\allowbreak{}Entladedynamik verwischen. Zu feine Intervalle liefern bei kurzer Sequenz wenig Gesamtzeit oder bei langer Sequenz hohe Redundanz. Das beobachtete Raster bevorzugt 16 Schritte mit zehn Minuten Aggregation. \textbf{Grenze:} Es ist ein diskretes datenabhängiges Optimum. Unterschiede zwischen Zellen und fehlende Trainingswiederholungen erlauben keinen universellen Optimalitätsbeweis. Prüfen, ob das Bild nachträgliche Sensitivität oder Auswahlgrundlage war. Siehe F029 und F031.

\topic{Abbildung 5.6: NMC-SOH-Trajektorien}
\begin{center}\includegraphics[width=\linewidth,height=0.60\textheight,keepaspectratio]{figures/thesis-page-109.pdf}\end{center}
Originalseite der Dissertation: PDF-Seite 109. Die folgende Interpretation bezieht sich auf die bezeichnete Abbildung.

\textbf{Frage:} Warum weicht Zelle 9 stärker ab, obwohl die Form teilweise getroffen wird?

\textbf{Antwort:} Formtreue allein reicht nicht. Ein systematisches vertikales Verschieben erzeugt trotz ähnlichem Trend erhebliche Residuen. Die U-förmige Referenz ist atypisch und schlechter in Absolutwerten reproduziert. \textbf{Grenze:} Der Plot enthält nachträgliche Rolling-Glättung. Die sichtbare Ruhe ist deshalb nicht automatisch rohe Netzqualität. Fensterbreite und kausale beziehungsweise zentrierte Filterdefinition vor der Prüfung nachschlagen. Siehe F033 und F034.

\topic{Abbildung 5.7: NMC-Scatter}
\begin{center}\includegraphics[width=\linewidth,height=0.60\textheight,keepaspectratio]{figures/thesis-page-109.pdf}\end{center}
Originalseite der Dissertation: PDF-Seite 109. Die folgende Interpretation bezieht sich auf die bezeichnete Abbildung.

\textbf{Frage:} Was bedeutet Abstand von der Diagonalen und warum geht Zeitinformation verloren?

\textbf{Antwort:} Jeder Punkt vergleicht Vorhersage und Referenz. Die Identitätslinie bedeutet perfekte Übereinstimmung. Systematische Abstände zeigen Bias oder Kalibrierfehler. Viele Punkte auf engem Bereich können durch lange Aufenthaltsdauer dort entstehen. \textbf{Grenze:} Scatter zeigt nicht, ob Fehler zusammenhängend über Stunden auftreten oder rasch wechseln. Deshalb zusammen mit der Trajektorie lesen. MSE-Einheiten sind Quadrat-Prozentpunkte, MAE-Einheiten Prozentpunkte. Siehe F034.

\topic{Abbildung 6.1: Gemeinsame Benchmarkkette}
\begin{center}\includegraphics[width=\linewidth,height=0.60\textheight,keepaspectratio]{figures/thesis-page-116.pdf}\end{center}
Originalseite der Dissertation: PDF-Seite 116. Die folgende Interpretation bezieht sich auf die bezeichnete Abbildung.

\textbf{Frage:} Wo muss eine Störung eingespeist werden, damit der Vergleich fair bleibt?

\textbf{Antwort:} Sie muss die tatsächlichen Rohkanäle und alle daraus kausal gebildeten Features beeinflussen. Ein gestörter Strom bei ungestörtem Qc würde DD einen künstlich sauberen Nebenkanal lassen. Die Kette soll diese Inkonsistenz verhindern. \textbf{Vertiefung:} Labels bleiben unverändert, da die Zielbatterie nicht durch einen Sensorfehler physikalisch anders betrieben wird. \textbf{Grenze:} Ein realer Regelkreis, der auf falsche Schätzungen reagiert, wäre eine andere Versuchsart. Siehe F039 und F040.

\topic{Abbildung 6.2: Störungstaxonomie}
\begin{center}\includegraphics[width=\linewidth,height=0.60\textheight,keepaspectratio]{figures/thesis-page-121.pdf}\end{center}
Originalseite der Dissertation: PDF-Seite 121. Die folgende Interpretation bezieht sich auf die bezeichnete Abbildung.

\textbf{Frage:} Warum gehören Gain, Offset, Rauschen, Verfügbarkeit und Initialzustand nicht in dieselbe Fehlerbeschreibung?

\textbf{Antwort:} Sie ändern unterschiedliche Mechanismen: Skalierung, Nullpunkt, zufällige Fluktuation, Informationsverlust oder internen Startzustand. Dadurch können sich gleiche nominale Modelle unter den Tests anders ordnen. \textbf{Grenze:} Die Taxonomie ist eine relevante Auswahl, kein vollständiger kartesischer Produktraum sämtlicher Sensorfehler. Zusammengesetzte Fehler und zeitabhängige Drift sind nicht umfassend abgedeckt. Die konkrete Stärke steht in der Szenariotabelle. Siehe F050 und F056.

\topic{Abbildung 6.3: Fensterprotokoll}
\begin{center}\includegraphics[width=\linewidth,height=0.60\textheight,keepaspectratio]{figures/thesis-page-126.pdf}\end{center}
Originalseite der Dissertation: PDF-Seite 126. Die folgende Interpretation bezieht sich auf die bezeichnete Abbildung.

\textbf{Frage:} Warum ein sauberes Vorfenster, gemeinsamer Scorebeginn und verlängertes Dropoutfenster?

\textbf{Antwort:} Das SOH-Netz benötigt kausalen Kontext, die GRU ein gültiges Eingabefenster. Erst danach wird gleichzeitige Leistung verglichen. Dropout braucht ausreichend Nachbeobachtung und eine gleich lange Baseline. \textbf{Grenze:} Das ist keine Kaltstartvalidierung ohne Vorgeschichte. Der typische Fensterselektor erfasst nicht jeden seltenen ungünstigen Betriebspunkt. Auf Nachfrage die sechs Selection Features und die Behandlung fehlender Zell-SOH-Kombinationen nennen. Siehe F023, F024 und F042.

\topic{Abbildung 6.4: Nominale Accuracy}
\begin{center}\includegraphics[width=\linewidth,height=0.60\textheight,keepaspectratio]{figures/thesis-page-128.pdf}\end{center}
Originalseite der Dissertation: PDF-Seite 128. Die folgende Interpretation bezieht sich auf die bezeichnete Abbildung.

\textbf{Frage:} Ist DD hier wirklich am besten und was bedeuten die Punkte?

\textbf{Antwort:} DD hat das kleinste zellgewichtete MAE und RMSE im dargestellten Protokoll. Die Punkte zeigen Zell-/\allowbreak{}Fensterwerte, die Balken den Makrowert und die Fehlerstriche ein hierarchisches 95-Prozent-Intervall. DD-MAE 0,0258 bedeutet 2,58 Prozentpunkte. \textbf{Grenze:} Nicht jede Zelle oder jeder Zeitpunkt muss dieselbe Rangfolge haben. Die korrigierten exakten Tests sind wegen kleiner Zellzahl nicht signifikant. Siehe F047, F063 und F065.

\topic{Abbildung 6.5: Current Gain}
\begin{center}\includegraphics[width=\linewidth,height=0.60\textheight,keepaspectratio]{figures/thesis-page-129.pdf}\end{center}
Originalseite der Dissertation: PDF-Seite 129. Die folgende Interpretation bezieht sich auf die bezeichnete Abbildung.

\textbf{Frage:} Warum steigt das obere Diagramm monoton, während eine Einzelkurve gelegentlich näher an der Referenz liegen kann?

\textbf{Antwort:} Oben wird je Betrag und Zelle die ungünstigere von zwei Gainrichtungen zusammengefasst. Unten steht eine konkrete Realisierung und ein Ausschnitt einer Zelle. Eine lokale Kompensation ist damit vereinbar. \textbf{Grenze:} Die obere Kurve ist keine Messung einer zufälligen positiven Drift und kein allgemeiner Monotoniebeweis. Die +3-Prozent-Eingangskurve illustriert eine Richtung, nicht jede adverse Auswahl. Siehe F048 und F049.

\topic{Abbildung 6.6: Gain über Lebenszeit und Ladeanker}
\begin{center}\includegraphics[width=\linewidth,height=0.60\textheight,keepaspectratio]{figures/thesis-page-130.pdf}\end{center}
Originalseite der Dissertation: PDF-Seite 130. Die folgende Interpretation bezieht sich auf die bezeichnete Abbildung.

\textbf{Frage:} Akkumuliert der Fehler jetzt oder wird er zurückgesetzt?

\textbf{Antwort:} Zwischen Ladeankern kann sich der Beitrag aufbauen. An einem Anker ändert sich der Integrationszustand, aber der Sensorfehler bleibt und wirkt danach erneut. Die Panels zeigen zeitlich schwankende Zusatzlast, normierte Zwischenankerentwicklung und Gesamt-MAE. \textbf{Wichtige Aussage:} HECM hat in diesem einzelnen High-Load-Lebenszeitlauf die kleinste Zusatzlast. Das widerspricht nicht dem sechszelligen Fensterresultat zugunsten DD. \textbf{Grenze:} Unterschiedliche Datengrundlage und Gewichtung. Siehe F052.

\topic{Abbildung 6.7: Lokales Stromrauschen}
\begin{center}\includegraphics[width=\linewidth,height=0.60\textheight,keepaspectratio]{figures/thesis-page-132.pdf}\end{center}
Originalseite der Dissertation: PDF-Seite 132. Die folgende Interpretation bezieht sich auf die bezeichnete Abbildung.

\textbf{Frage:} Warum schwankt DD lokal sichtbar, während die globale Fehleränderung klein bleibt?

\textbf{Antwort:} Rohstrom und Ableitung wirken direkt auf DD. Die integrierenden Verfahren glätten wechselnde Fehler stärker im sichtbaren SOC, obwohl jeder verrauschte Sample eingeht. Nullgemittelte lokale Änderungen können im Tagesmittel wenig ausmachen. \textbf{Grenze:} Ein Knick ist nicht automatisch ein Fehler im Modell. Resetlogs, Lastwechsel und SOH-Stundenupdates müssen unterschieden werden. Der gezeigte Ausschnitt wird als reset-frei beschrieben. Siehe F054 und F055.

\topic{Abbildung 6.8: Rauschen über sechs Zellen}
\begin{center}\includegraphics[width=\linewidth,height=0.60\textheight,keepaspectratio]{figures/thesis-page-133.pdf}\end{center}
Originalseite der Dissertation: PDF-Seite 133. Die folgende Interpretation bezieht sich auf die bezeichnete Abbildung.

\textbf{Frage:} Darf man bei einem höheren Mittelbalken sicher größere Empfindlichkeit behaupten?

\textbf{Antwort:} Ein höherer Punktwert beschreibt die beobachtete Richtung, ein breites Intervall aber erhebliche Unsicherheit. Insbesondere hoher Stromnoise zeigt HECM mit größerer mittlerer Penalty, deren Intervall null einschließt. \textbf{Grenze:} Nicht aus dem Balken allein eine bewiesene Populationseigenschaft machen. Die Seeds messen Störungsvariation, die Zellen bleiben die unabhängigen Einheiten. Siehe F061 bis F066.

\topic{Abbildung 6.9: Initialisierungs-Recovery}
\begin{center}\includegraphics[width=\linewidth,height=0.60\textheight,keepaspectratio]{figures/thesis-page-134.pdf}\end{center}
Originalseite der Dissertation: PDF-Seite 134. Die folgende Interpretation bezieht sich auf die bezeichnete Abbildung.

\textbf{Frage:} Warum bleibt DM nahe zehn Prozentpunkten, während DD unter die Schwelle fällt und wieder steigt?

\textbf{Antwort:} DM/\allowbreak{}HDM behalten die falsche Ladungsbilanz bis zu einer wirksamen Verankerung. HECM korrigiert über Spannung. DD kann die Qc-Abhängigkeit zeitweise kompensieren, aber bei anderem Kontext wieder stärker auf das gestörte Feature reagieren. Panel b zeigt gepaarte Ausgangsdifferenz, nicht Dataset-MAE. Panel c aggregiert mehrere Fenster und Zellen, nicht nur die Beispielkurve. \textbf{Grenze:} Erste Rückkehr ist nicht persistent. Siehe F041 bis F046.

\topic{Abbildung 6.10: Signalverlust und Jitter}
\begin{center}\includegraphics[width=\linewidth,height=0.60\textheight,keepaspectratio]{figures/thesis-page-136.pdf}\end{center}
Originalseite der Dissertation: PDF-Seite 136. Die folgende Interpretation bezieht sich auf die bezeichnete Abbildung.

\textbf{Frage:} Weshalb ist 2-Prozent-Sample Loss gravierender als großer Zeitjitter?

\textbf{Antwort:} Gefrorene Werte verlieren reale Signaländerung, während der Jittertest erhaltene Werte zeitbewusst verarbeitet. Die integrativen Modelle verlieren insbesondere Ladungsinformation. \textbf{Grenze:} Das hängt von der exakten Interventionssemantik ab. Asynchrone Spannung/\allowbreak{}Strom-Paare oder falsche Zeitstempel könnten viel kritischer sein und sind nicht automatisch durch diese Balken abgedeckt. Fehlerstriche sind Konfidenzintervalle, keine Worst-Case-Grenzen. Siehe F056.

\topic{Abbildung 6.11: Burst Dropout}
\begin{center}\includegraphics[width=\linewidth,height=0.60\textheight,keepaspectratio]{figures/thesis-page-136.pdf}\end{center}
Originalseite der Dissertation: PDF-Seite 136. Die folgende Interpretation bezieht sich auf die bezeichnete Abbildung.

\textbf{Frage:} Warum ist DD nach dem Ausfall relativ unauffällig, HECM aber länger verschoben?

\textbf{Antwort:} Im Rolling-DD verlassen eingefrorene Rohwerte später das Fenster. Neue Multikanalinformation kann den Einfluss des falschen Qc reduzieren. HECM muss eine verschobene Integrations-/\allowbreak{}Polarisationshistorie mit begrenzter LFP-Spannungsinformation korrigieren. Das sind plausible Mechanismen. \textbf{Grenze:} Der negative globale DD-ΔMAE mit nullüberlappendem Intervall ist kein Nachweis einer Verbesserung. Panel b misst Abweichung vom sauberen Modell, nicht absoluten Referenzfehler. Siehe F057.

\topic{Abbildung 6.12: Spike-Penalty nach Zelle und SOH}
\begin{center}\includegraphics[width=\linewidth,height=0.60\textheight,keepaspectratio]{figures/thesis-page-137.pdf}\end{center}
Originalseite der Dissertation: PDF-Seite 137. Die folgende Interpretation bezieht sich auf die bezeichnete Abbildung.

\textbf{Frage:} Warum reagiert DD hier stärker als in der globalen Heatmap?

\textbf{Antwort:} Diese Grafik konzentriert sich auf Spike-Samples und löst nach Zelle und SOH auf. Die globale Heatmap mittelt über lange Fenster, in denen fast alle Samples ungestört sind. Beide Aussagen können gleichzeitig richtig sein. \textbf{Grenze:} Farbintensität nicht ohne jeweilige Skala vergleichen. Eine stärker gefärbte Zelle beweist keinen kausalen Alterungseinfluss, da Last und SOH gekoppelt sein können. Siehe F053 und F069.

\topic{Abbildung 6.13: Ereignisausgerichteter Spikeverlauf}
\begin{center}\includegraphics[width=\linewidth,height=0.60\textheight,keepaspectratio]{figures/thesis-page-138.pdf}\end{center}
Originalseite der Dissertation: PDF-Seite 138. Die folgende Interpretation bezieht sich auf die bezeichnete Abbildung.

\textbf{Frage:} Was ist Excess Absolute Error und kann er negativ sein?

\textbf{Antwort:} Er beschreibt zusätzlichen absoluten Zielfehler gegenüber dem gepaarten sauberen Lauf. Er kann negativ sein, wenn die Störung eine bestehende Abweichung kompensiert. Davon verschieden ist der immer nichtnegative Betrag der Differenz zwischen gestörter und sauberer Ausgabe. \textbf{Prüfpunkt:} Ein kausaler Spike darf keine vorherige Wirkung verursachen. Bei vorgezogenen Unterschieden Zeitmarken, Rolling-Alignment, Vorereignisse und Mittelung prüfen. \textbf{Grenze:} Der plausible dU/\allowbreak{}dt-Pfad braucht Ablation zur isolierten Attribution. Siehe F053 und F054.

\topic{Abbildung 6.14: ADC-Quantization}
\begin{center}\includegraphics[width=\linewidth,height=0.60\textheight,keepaspectratio]{figures/thesis-page-139.pdf}\end{center}
Originalseite der Dissertation: PDF-Seite 139. Die folgende Interpretation bezieht sich auf die bezeichnete Abbildung.

\textbf{Frage:} Warum sind Stufen im Sensorsignal sichtbar, aber nicht gleichermaßen im SOC?

\textbf{Antwort:} Integration, Observer und zeitliche Lernabbildung übertragen Eingangsrundung unterschiedlich. Schritte von 0,01 A, 5 mV und 0,5 °C werden gemeinsam im gewählten Szenario geprüft. Alle Makrointervalle überlappen hier null. \textbf{Grenze:} Das ist Sensor Quantization, nicht Weight Quantization. Es bewertet weder beliebige ADC-Bitzahlen noch Front-end-Kalibrierung vollständig. Full-scale-Bereich und effektive Auflösung wären nötig, um Bits aus Schrittweiten abzuleiten. Siehe F079 für den anderen Quantization-Begriff.

\topic{Abbildung 6.15: Cross-Scenario-Heatmap}
\begin{center}\includegraphics[width=\linewidth,height=0.60\textheight,keepaspectratio]{figures/thesis-page-141.pdf}\end{center}
Originalseite der Dissertation: PDF-Seite 141. Die folgende Interpretation bezieht sich auf die bezeichnete Abbildung.

\textbf{Frage:} Was bedeuten positive, negative und sehr kleine Werte?

\textbf{Antwort:} Positiv bedeutet höheren MAE relativ zur jeweiligen Modellbaseline, negativ geringeren. Die Heatmap zeigt Zusatzfehler, keine absoluten nominalen Fehler. Current Offset dominiert mehrere Modelle. \textbf{Grenze:} Kleine globale Spikewerte verbergen lokale Peaks. Gain/\allowbreak{}Offset verwenden adverse Richtungen, andere Zeilen andere definierte Interventionen. Die Heatmap ist daher ein kompakter Index mit methodisch unterschiedlichen Störungsformen, kein universeller Sicherheitsmaßstab. Siehe F048 und F067.

\topic{Abbildung 6.16: Entscheidungssynthese}
\begin{center}\includegraphics[width=\linewidth,height=0.60\textheight,keepaspectratio]{figures/thesis-page-142.pdf}\end{center}
Originalseite der Dissertation: PDF-Seite 142. Die folgende Interpretation bezieht sich auf die bezeichnete Abbildung.

\textbf{Frage:} Warum gewinnt nicht überall derselbe Vertreter und was bedeutet große Dreiecksfläche?

\textbf{Antwort:} Die Achsen fassen verschiedene relative Kriterien zusammen. DD ist bei nominaler Accuracy und mehreren Robustheitsdefinitionen stark, HECM bei Recovery. Die rechte Seite ändert Prioritätsgewichte. \textbf{Grenze:} Flächen sind nicht linear in jedem Einzelwert und die Achsenskalierung ist relativ. Ein Score nahe eins ist keine Zuverlässigkeitswahrscheinlichkeit. Für eine reale Entscheidung muss die konkrete Störung und ein absolutes Anforderungslimit betrachtet werden. Siehe F068.

\topic{Abbildung 6.17: Hardware-Software-Äquivalenz}
\begin{center}\includegraphics[width=\linewidth,height=0.60\textheight,keepaspectratio]{figures/thesis-page-143.pdf}\end{center}
Originalseite der Dissertation: PDF-Seite 143. Die folgende Interpretation bezieht sich auf die bezeichnete Abbildung.

\textbf{Frage:} Ist der sehr kleine Fehler hier die Batterie-SOC-Genauigkeit?

\textbf{Antwort:} Nein. Verglichen werden zwei Implementierungen desselben Schätzers. Kleine MCU-zu-Software-Abweichung zeigt die numerische Konsistenz des Exports. Beide können gegenüber Dataset-SOC einen deutlich größeren Fehler besitzen. \textbf{Vertiefung:} Logarithmische Achse und Prozentpunkteinheit nennen. \textbf{Grenze:} Gleiche Ausgaben beweisen nicht die Richtigkeit der gemeinsamen Referenz oder des Modells. Siehe F084 und F089.

\topic{Abbildung 6.18: Inferenzlatenzen}
\begin{center}\includegraphics[width=\linewidth,height=0.60\textheight,keepaspectratio]{figures/thesis-page-144.pdf}\end{center}
Originalseite der Dissertation: PDF-Seite 144. Die folgende Interpretation bezieht sich auf die bezeichnete Abbildung.

\textbf{Frage:} Warum sind DM/\allowbreak{}HDM nahezu deterministisch und DD breiter verteilt?

\textbf{Antwort:} DM/\allowbreak{}HDM führen sehr wenige Operationen aus. DD berechnet zahlreiche Matrix- und Activation Operations, HECM Lookup/\allowbreak{}Observer-Schritte. Wiederholte Zellreplays zeigen beobachtete Laufzeitstreuung. \textbf{Grenze:} Die hier gezeigte DD-Ausführung ist continuous, nicht die teure Rolling-Variante des primären Robustheitsbenchmarks. Der Maximalwert ist ein beobachteter Extremwert, kein formaler WCET-Beweis. Siehe F088 bis F090.

\topic{Abbildung 6.19: Flash und RAM}
\begin{center}\includegraphics[width=\linewidth,height=0.60\textheight,keepaspectratio]{figures/thesis-page-145.pdf}\end{center}
Originalseite der Dissertation: PDF-Seite 145. Die folgende Interpretation bezieht sich auf die bezeichnete Abbildung.

\textbf{Frage:} Warum braucht HECM viel Flash, aber wenig RAM?

\textbf{Antwort:} Seine voridentifizierten Parameterflächen liegen als persistente Tabellen im Flash. Die aktuellen Observerzustände bleiben klein. DD speichert Gewichte im Flash und Zustände/\allowbreak{}Arbeitsdaten im RAM. Das Ergebnis betrifft die konkrete Tabellendichte und Firmware, nicht jede ECM-Implementierung. \textbf{Grenze:} Das gemeinsame SOH-LSTM ist nicht Teil dieser MCU-Messung. Ein vollständiger SOC/\allowbreak{}SOH-Systembedarf wäre größer und müsste separat gemessen werden. Siehe F087.

\topic{Abbildung 6.20: DD-Ausführungsmodi}
\begin{center}\includegraphics[width=\linewidth,height=0.60\textheight,keepaspectratio]{figures/thesis-page-146.pdf}\end{center}
Originalseite der Dissertation: PDF-Seite 146. Die folgende Interpretation bezieht sich auf die bezeichnete Abbildung.

\textbf{Frage:} Wie können gleiche Gewichte gleichzeitig sehr unterschiedliche Fehler und Laufzeiten erzeugen?

\textbf{Antwort:} Rolling baut aus dem letzten vollständigen Fenster neu auf. Continuous trägt den Zustand fort. Periodic Reset löscht Kontext, ohne das ganze historische Fenster neu einzuspielen. Unterschiedliche zeitliche Information verändert die Funktion trotz gleicher Gewichte. \textbf{Beleg:} Rolling etwa 724 ms und 67,3 KiB, Continuous etwa 0,425 ms und 4,1 KiB. \textbf{Grenze:} Ähnliche nominale Accuracy ist keine vollständige Störungsäquivalenz. Siehe F085 und F086.

\topic{Abbildung 7.1: Stateful LSTM plus Kopf}
\begin{center}\includegraphics[width=\linewidth,height=0.60\textheight,keepaspectratio]{figures/thesis-page-157.pdf}\end{center}
Originalseite der Dissertation: PDF-Seite 157. Die folgende Interpretation bezieht sich auf die bezeichnete Abbildung.

\textbf{Frage:} Was läuft einmal pro Sample und was bleibt über Aufrufe erhalten?

\textbf{Antwort:} Das LSTM verarbeitet den aktuellen Featurevektor. h und c bleiben erhalten, der MLP-Kopf bildet h auf SOC oder SOH ab. Für SOC ist H = 64, für SOH H = 128. \textbf{Grenze:} Das Diagramm ist nicht das stündliche SOH-Netz mit 20 aggregierten Eingängen aus Kapitel 6. Hier werden beide Aufgaben auf dem 1-Hz-Stream behandelt. Siehe F073 und F084.

\topic{Abbildung 7.2: Aufgabenabhängige Architekturen}
\begin{center}\includegraphics[width=\linewidth,height=0.60\textheight,keepaspectratio]{figures/thesis-page-158.pdf}\end{center}
Originalseite der Dissertation: PDF-Seite 158. Die folgende Interpretation bezieht sich auf die bezeichnete Abbildung.

\textbf{Frage:} Warum unterschiedliche Features und verschiedene Ausgangsaktivierungen?

\textbf{Antwort:} SOC erhält unter anderem Qc und Ableitungen für kurzfristige Dynamik. SOH erhält Zeit und Durchsatz für langfristige Alterungsinformation. SOC wird durch Sigmoid auf 0 bis 1 begrenzt, SOH bleibt vor Nachverarbeitung linear. \textbf{Grenze:} Zeit kann ein starker Datensatzproxy sein. Eine Modellleistung mit Zeitfeature beweist keine robuste Extrapolation auf andere Alterungsraten. Ein SOH-Wert über eins kann bei nominaler Referenz physikalisch möglich sein, muss aber geprüft werden. Siehe F019.

\topic{Abbildung 7.3: Deploymentpipeline}
\begin{center}\includegraphics[width=\linewidth,height=0.60\textheight,keepaspectratio]{figures/thesis-page-160.pdf}\end{center}
Originalseite der Dissertation: PDF-Seite 160. Die folgende Interpretation bezieht sich auf die bezeichnete Abbildung.

\textbf{Frage:} Wie kontrollieren Sie, dass Kompressionsvarianten fair verglichen werden?

\textbf{Antwort:} Gleicher Base-Ausgangspunkt, identischer Featurestream, dokumentierte Gewichte/\allowbreak{}Scaler und dieselbe Zustandsführung. Danach getrennte Kompression, C-Export, numerischer Vergleich und Hardwaremessung. \textbf{Grenze:} Unterschiedliches Fine-Tuning ist selbst eine Intervention und braucht für kausale Attribution eine passende Kontrolle. Das Schema allein ersetzt keine Versions- und Datenprovenienz. Siehe F078 und Quellenauftrag Q5.

\topic{Abbildung 7.4: Numerisches Pruning-Beispiel}
\begin{center}\includegraphics[width=\linewidth,height=0.60\textheight,keepaspectratio]{figures/thesis-page-162.pdf}\end{center}
Originalseite der Dissertation: PDF-Seite 162. Die folgende Interpretation bezieht sich auf die bezeichnete Abbildung.

\textbf{Frage:} Warum verschwindet h2 nicht nur in einer Gatezeile?

\textbf{Antwort:} Alle vier Gates, der recurrent input, Bias, Zustände und MLP-Eingänge beziehen sich auf denselben Kanal. Entfernen muss diese Abhängigkeiten gemeinsam berücksichtigen. Im Beispiel hat h2 den kleinsten aggregierten L2-Score. \textbf{Grenze:} Die kleinen Matrixzahlen sind illustrativ, keine echten trainierten Gewichte. Der Score misst Gewichtsmagnitude, nicht direkt Sensitivität des SOH/\allowbreak{}SOC-Ausgangs. Siehe F075.

\topic{Abbildung 7.5: Quantization und Export}
\begin{center}\includegraphics[width=\linewidth,height=0.60\textheight,keepaspectratio]{figures/thesis-page-164.pdf}\end{center}
Originalseite der Dissertation: PDF-Seite 164. Die folgende Interpretation bezieht sich auf die bezeichnete Abbildung.

\textbf{Frage:} Was bleibt grün/\allowbreak{}rot beziehungsweise als Ganzzahl oder Float gespeichert?

\textbf{Antwort:} Die Recurrent-Matrizen werden zeilenweise in Codes und Skalen aufgeteilt. Die effektiven Werte entstehen im Rechenpfad wieder in Float. Die Netzstruktur bleibt erhalten. \textbf{Grenze:} Ein gespeichertes INT8-Muster sagt nichts darüber, ob Multiplikation und Akkumulation als Integer laufen. Die Ein- und Ausgabe und recurrent states bleiben hier FP32. Die Detailgrenze steht zusätzlich in A.11. Siehe F079 bis F082.

\topic{Abbildung 7.6: SOC über die gesamte Trajektorie}
\begin{center}\includegraphics[width=\linewidth,height=0.60\textheight,keepaspectratio]{figures/thesis-page-169.pdf}\end{center}
Originalseite der Dissertation: PDF-Seite 169. Die folgende Interpretation bezieht sich auf die bezeichnete Abbildung.

\textbf{Frage:} Warum sieht man oben fast nur eine dichte Fläche und unten trotzdem systematische Fehler?

\textbf{Antwort:} Viele Ladezyklen werden über lange Zeit zusammengedrängt. Kleine Unterschiede überlagern sich visuell. Das Fehlerpanel macht sie sichtbar. Base/\allowbreak{}Pruned/\allowbreak{}Quantized haben global MAE 2,68/\allowbreak{}2,34/\allowbreak{}2,79 Prozentpunkte. \textbf{Grenze:} Visuelle Überdeckung ist kein numerischer Gleichheitsnachweis. Es ist die repräsentative Einzelzelltrajektorie der Embedded-Studie, nicht automatisch das sechszellige Robustheitsmittel. Siehe F084 und F095.

\topic{Abbildung 7.7: Gefilterter SOH}
\begin{center}\includegraphics[width=\linewidth,height=0.60\textheight,keepaspectratio]{figures/thesis-page-170.pdf}\end{center}
Originalseite der Dissertation: PDF-Seite 170. Die folgende Interpretation bezieht sich auf die bezeichnete Abbildung.

\textbf{Frage:} Wie viel der ruhigen Schätzung stammt vom Netz und wie viel vom Filter?

\textbf{Antwort:} Die gezeigte Ausgabe durchläuft Anfangsausrichtung, zwei EMAs und Rate-Limiter. Deshalb ist die Linienruhe eine Eigenschaft der gesamten Pipeline. Das primäre Ergebnis bevorzugt Base gegenüber den komprimierten Varianten. \textbf{Grenze:} Die zweite EMA hat eine sehr lange Antwortzeit. Die Diagrammwerte dürfen nicht als ungefilterter SOH-MAE zitiert werden. Andere Werte in A.6 benötigen klare Sequenzprovenienz. Siehe F092 bis F094.

\topic{Abbildung 7.8: SOC-Fehlerverteilung}
\begin{center}\includegraphics[width=\linewidth,height=0.60\textheight,keepaspectratio]{figures/thesis-page-171.pdf}\end{center}
Originalseite der Dissertation: PDF-Seite 171. Die folgende Interpretation bezieht sich auf die bezeichnete Abbildung.

\textbf{Frage:} Warum sind Boxplot und Histogramm beide nötig?

\textbf{Antwort:} Die Boxplots zeigen Verteilung absoluter Fehler, Histogramme zusätzlich deren Vorzeichenstruktur. Pruned hat den kleinsten Median, aber seine Tails und lokalen Abschnitte müssen separat betrachtet werden. \textbf{Grenze:} Überlappende Histogramme können zeitlich völlig verschiedene Fehler enthalten. Die Beobachtungen sind korrelierte Samples, keine Millionen unabhängigen Versuche. Boxplot-Whisker nicht automatisch als Min-Max interpretieren. Siehe F065 und F078.

\topic{Abbildung 7.9: SOH-Fehlerverteilung}
\begin{center}\includegraphics[width=\linewidth,height=0.60\textheight,keepaspectratio]{figures/thesis-page-172.pdf}\end{center}
Originalseite der Dissertation: PDF-Seite 172. Die folgende Interpretation bezieht sich auf die bezeichnete Abbildung.

\textbf{Frage:} Warum ist die Form unregelmäßiger als beim SOC?

\textbf{Antwort:} SOH verändert sich langsam, wird aus spärlichen Kapazitätsankern interpoliert und stark gefiltert. Bestimmte Alterungsabschnitte können lange einen ähnlichen Bias tragen und Häufungen erzeugen. \textbf{Grenze:} Verteilungsform identifiziert nicht eindeutig einen elektrochemischen Alterungsmechanismus. Ein größeres Absolutfehlerband beim komprimierten Modell beweist keine rekurrente Instabilität. Die Zeitlokalisation liefert A.5. Siehe F094 und F095.

\topic{Abbildung 7.10: Dynamischer SOC-Ausschnitt}
\begin{center}\includegraphics[width=\linewidth,height=0.60\textheight,keepaspectratio]{figures/thesis-page-172.pdf}\end{center}
Originalseite der Dissertation: PDF-Seite 172. Die folgende Interpretation bezieht sich auf die bezeichnete Abbildung.

\textbf{Frage:} Warum ist Pruned trotz besserem Gesamt-MAE lokal teilweise weiter von Base entfernt?

\textbf{Antwort:} Pruning verändert die Funktion. Global günstigeres Verhalten muss nicht in jedem dynamischen Abschnitt günstiger sein. Außerdem ist die Base-Ausgabe keine Ground Truth. Der untere Residuenplot ist für die tatsächliche Genauigkeit wichtiger als allein die Überdeckung der Modellkurven. \textbf{Grenze:} Der gewählte 30k-bis-40k-Sekunden-Ausschnitt ist eine Illustration, kein repräsentatives Mittel aller Transienten. Siehe F078 und F095.

\topic{Abbildung 7.11: Check-up-Ausschnitt}
\begin{center}\includegraphics[width=\linewidth,height=0.60\textheight,keepaspectratio]{figures/thesis-page-173.pdf}\end{center}
Originalseite der Dissertation: PDF-Seite 173. Die folgende Interpretation bezieht sich auf die bezeichnete Abbildung.

\textbf{Frage:} Was zeigt der längere Abschnitt, das im kurzen Pulsfenster fehlt?

\textbf{Antwort:} Längere Ladungsintegration, Betriebsphasenwechsel und Diagnoseintervalle zeigen andere Dynamik. Die Varianten bleiben im gezeigten Abschnitt grundsätzlich konsistent, mit sichtbaren lokalen Fehlerunterschieden. \textbf{Grenze:} Ein stabiler Ausschnitt beweist keine allgemeine Langzeitstabilität. Die Datenauswahl der Diagnosephasen und die Maskierung von Kapazitätstests müssen nachvollziehbar sein. Siehe F020 und F095.

\topic{Abbildung 7.12: Ressourcenlandschaft}
\begin{center}\includegraphics[width=\linewidth,height=0.60\textheight,keepaspectratio]{figures/thesis-page-177.pdf}\end{center}
Originalseite der Dissertation: PDF-Seite 177. Die folgende Interpretation bezieht sich auf die bezeichnete Abbildung.

\textbf{Frage:} Warum stimmen theoretischer Gewichtsspeicher und gelinkter Flash nicht überein?

\textbf{Antwort:} Die Theorie zählt Parameterrepräsentationen. Die Firmware enthält zusätzlich Code, Skalierung, konstante Tabellen und Laufzeitunterstützung. RAM ist wiederum ein anderer Speicherbereich mit Zuständen, Puffern und Stack. \textbf{Grenze:} Weight-only-Quantization lässt viele FP32-Daten unverändert. Die genauen RAM-Unterschiede nicht allein mit einem Viertel pro Gewicht erklären. Für vollständige Attribution Linker-Sektionen und Stackmessung zeigen. Siehe F081 und F087.

\topic{Abbildung 7.13: Hostlatenzverteilungen}
\begin{center}\includegraphics[width=\linewidth,height=0.60\textheight,keepaspectratio]{figures/thesis-page-178.pdf}\end{center}
Originalseite der Dissertation: PDF-Seite 178. Die folgende Interpretation bezieht sich auf die bezeichnete Abbildung.

\textbf{Frage:} Warum steht Quantized weiter rechts und weshalb sind SOH-Abstände größer?

\textbf{Antwort:} Der Mixed-Precision-Kernel erhöht hier den Aufwand, und SOH besitzt den größeren recurrent core. Host Latency umfasst aber auch UART und Rücktransport. \textbf{Grenze:} Ein fester Kommunikationsanteil kann Kernelunterschiede relativ verkleinern. Es sind keine unabhängigen Energieverteilungen. Bei einem Vergleich mit Kapitel 6 zuerst Messgrenzen und Modellsemantik klären. Siehe F082, F083 und F089.

\topic{Abbildung A.1: Vollständige Testmatrix}
\begin{center}\includegraphics[width=\linewidth,height=0.60\textheight,keepaspectratio]{figures/thesis-page-216.pdf}\end{center}
Originalseite der Dissertation: PDF-Seite 216. Die folgende Interpretation bezieht sich auf die bezeichnete Abbildung.

\textbf{Frage:} Können Sie jede Zeile einem Eingriff und einer Wiederholungszahl zuordnen?

\textbf{Antwort:} Die Matrix trennt manipulierte Kanäle/\allowbreak{}Zustände von wiederholten Seeds. Baseline und Initialisierung erklären zusammen mit 18 Störungszeilen die 20 Definitionen. Die beiden Vorzeichen von drei Gainstufen und einem Offset erweitern die Messstörungsfälle auf 22 Subfälle. \textbf{Grenze:} Nicht aus der Zahl der farbigen Punkte die unabhängige Stichprobengröße ableiten. Die unabhängigen Einheiten sind weiterhin sechs Zellen. Siehe F066 und F067.

\topic{Abbildung A.2: DD-Latenzmodi}
\begin{center}\includegraphics[width=\linewidth,height=0.60\textheight,keepaspectratio]{figures/thesis-page-217.pdf}\end{center}
Originalseite der Dissertation: PDF-Seite 217. Die folgende Interpretation bezieht sich auf die bezeichnete Abbildung.

\textbf{Frage:} Warum liegen kontinuierlich und periodisch reset nahezu zusammen, obwohl ihre Genauigkeit stark verschieden ist?

\textbf{Antwort:} Beide führen pro Update nur einen rekurrenten Schritt aus. Der gelegentliche Zustandsreset spart keine relevanten Matrixoperationen und kostet wenig. Er verändert aber die verfügbare Historie. Rolling wiederholt dagegen das gesamte Fenster. \textbf{Grenze:} Laufzeitähnlichkeit bedeutet keine funktionale Gleichwertigkeit. Dieses Anhangsbild ergänzt die Fehler- und Speicherpanels in Abbildung 6.20. Siehe F085 und F086.

\topic{Abbildung A.3: Lookup-Sensitivität}
\begin{center}\includegraphics[width=\linewidth,height=0.60\textheight,keepaspectratio]{figures/thesis-page-218.pdf}\end{center}
Originalseite der Dissertation: PDF-Seite 218. Die folgende Interpretation bezieht sich auf die bezeichnete Abbildung.

\textbf{Frage:} Kann man daraus sagen, die HECM-Tabelle sei unwichtig?

\textbf{Antwort:} Nein. Links sind lokale Änderungen der Störungspenalty relativ zur jeweiligen sauberen Variante, rechts Recovery unter derselben Lookupvariation. Viele Interaktionen sind klein gegenüber dem großen Offseteffekt, aber Widerstand -10 Prozent zeigt einen zensierten Zellfall und ein deutlich höheres Zeitmittel. \textbf{Grenze:} OAT-Variationen prüfen keine kombinierten Unsicherheiten oder falschen Kennlinienformen. Tabelle A.1 zusätzlich nennen, nicht nur farbige Felder. Siehe F059 und F060.

\topic{Abbildung A.4: Pruning-Diagnostik}
\begin{center}\includegraphics[width=\linewidth,height=0.60\textheight,keepaspectratio]{figures/thesis-page-222.pdf}\end{center}
Originalseite der Dissertation: PDF-Seite 222. Die folgende Interpretation bezieht sich auf die bezeichnete Abbildung.

\textbf{Frage:} Warum schneiden Sie links die kleinsten 19 Scores ab und was beweist die Gewichtsverteilung rechts?

\textbf{Antwort:} Das entspricht der implementierten deterministischen Saliency-Regel für 64 auf 45 SOC-Kanäle. Rechts wird eine zentrale Gewichtungsverteilung vor und nach Auswahl verglichen. \textbf{Grenze:} Eine ähnliche Verteilungsform bedeutet weder ähnliche Funktion noch bewiesene Regularisierung. Rekurrente Struktur und Vorzeichen-/\allowbreak{}Kanalzusammenhänge gehen im Histogramm verloren. Für Kausalität braucht es zusätzliche Kontrollen. Siehe F075 und F078.

\topic{Abbildung A.5: Zehn Zeitsegmente}
\begin{center}\includegraphics[width=\linewidth,height=0.60\textheight,keepaspectratio]{figures/thesis-page-223.pdf}\end{center}
Originalseite der Dissertation: PDF-Seite 223. Die folgende Interpretation bezieht sich auf die bezeichnete Abbildung.

\textbf{Frage:} Ist der steigende Fehler rechts im Leben ein Quantization-Drift-Problem?

\textbf{Antwort:} Nicht zwingend. Sowohl Base als auch Quantized werden gegenüber dem Ziel schlechter, während ihre gegenseitige Abweichung nicht entsprechend wächst. Das spricht in diesem Replay gegen akkumulierte Quantization-Abweichung als Hauptursache. \textbf{Grenze:} Die Segmente sind gleich viele Samples, nicht unabhängige Wiederholungen. Veränderungen von SOH, Last und Labelqualität können gemeinsam wirken. Siehe F095.

\topic{Abbildung A.6: SOH-Filterstufen}
\begin{center}\includegraphics[width=\linewidth,height=0.60\textheight,keepaspectratio]{figures/thesis-page-224.pdf}\end{center}
Originalseite der Dissertation: PDF-Seite 224. Die folgende Interpretation bezieht sich auf die bezeichnete Abbildung.

\textbf{Frage:} Warum wechselt das Modellranking nach dem Filter?

\textbf{Antwort:} Filterung wirkt auf Richtung, Frequenz und zeitliche Struktur der Residuen, nicht nur auf deren Amplitude. Ein Modell mit mehr hochfrequentem Anteil kann stärker profitieren als eines mit langsamem Bias. Die lange EMA verändert den Zeitbezug zusätzlich. \textbf{Grenze:} Die hier genannten MAE-Werte weichen von der Haupt-SOH-Auswertung ab. Die präzise Daten-/\allowbreak{}Zustandsprovenienz vor einer eindeutigen Erklärung klären. Siehe F092 bis F094 und P06.

\topic{Abbildung A.7: Utility-Gewichtungen}
\begin{center}\includegraphics[width=\linewidth,height=0.60\textheight,keepaspectratio]{figures/thesis-page-225.pdf}\end{center}
Originalseite der Dissertation: PDF-Seite 225. Die folgende Interpretation bezieht sich auf die bezeichnete Abbildung.

\textbf{Frage:} Bedeutet viel Rot, dass Pruning immer gewinnt?

\textbf{Antwort:} Es gewinnt für einen großen Teil der untersuchten diskreten Prioritätskombinationen, insbesondere beim SOC. Beim SOH können Accuracy-Priorität Base und Flash-Priorität Quantized attraktiv machen. \textbf{Grenze:} Die Häufigkeit über ein künstliches Gitter ist keine Nutzer- oder Einsatzwahrscheinlichkeit. Harte Grenzen gehören vor den gewichteten Score. \textbf{Rechenfrage:} 1771 Kombinationen über Stars-and-bars erklären. Siehe F098 und R11.

\topic{Abbildung A.8: Eingangsbuffer-Fehler}
\begin{center}\includegraphics[width=\linewidth,height=0.60\textheight,keepaspectratio]{figures/thesis-page-226.pdf}\end{center}
Originalseite der Dissertation: PDF-Seite 226. Die folgende Interpretation bezieht sich auf die bezeichnete Abbildung.

\textbf{Frage:} Warum können Peaks groß sein, während ΔMAE61 klein bleibt?

\textbf{Antwort:} Der Peak betrifft einen kurzen Moment, während der MAE-Kontrast 61 Samples mittelt. Ein wieder abklingender Zustandseffekt kann lokal relevant sein und global wenig beitragen. Die Fault-Clean-Differenz und die Änderung des Referenzfehlers sind verschiedene Metriken. \textbf{Grenze:} Nur ein bestimmtes Mantissenbit und begrenzte Events wurden getestet. SOH-Recoveryprozente beziehen sich auf 60 Sekunden, nicht auf den 24-h-Test in Kapitel 6. Siehe F097.

\topic{Abbildung A.9: Komplexitätsskalierung}
\begin{center}\includegraphics[width=\linewidth,height=0.60\textheight,keepaspectratio]{figures/thesis-page-227.pdf}\end{center}
Originalseite der Dissertation: PDF-Seite 227. Die folgende Interpretation bezieht sich auf die bezeichnete Abbildung.

\textbf{Frage:} Warum steigen die MAC-Kurven quadratisch und die Reduktionskurve nähert sich einer Grenze?

\textbf{Antwort:} Recurrent-Matrizen enthalten H mal H-Verbindungen pro Gate. Bei festem Input und Kopf dominiert dieser Term für große H. Entfernt man einen Anteil p, beträgt die asymptotische Einsparung 2p-p². Die endlichen tatsächlich gewählten Architekturen liegen darunter. \textbf{Grenze:} Analytische MACs sind keine direkte Hardwarezeit und enthalten Activation Functions nicht. Siehe F077 und R07.

\topic{Abbildung A.10: Reichweite des L2-Kriteriums}
\begin{center}\includegraphics[width=\linewidth,height=0.60\textheight,keepaspectratio]{figures/thesis-page-228.pdf}\end{center}
Originalseite der Dissertation: PDF-Seite 228. Die folgende Interpretation bezieht sich auf die bezeichnete Abbildung.

\textbf{Frage:} Warum nicht Gradient oder Aktivierung statt Gewichtsnorm?

\textbf{Antwort:} Der gewählte Score ist deterministisch aus den trainierten Gewichten berechenbar und passt zur Kanalstruktur. Gradientenkriterien benötigen Daten und Backward Pass, Activation-Kriterien einen repräsentativen Kalibrierstream und eine zeitliche Aggregation. \textbf{Grenze:} Die Tabelle vergleicht Informationsbedarf und Methodenscope, nicht experimentell gemessene Überlegenheit. Keine Aussage L2 ist am besten ohne Vergleichsexperiment. Siehe F075 und F076.

\topic{Abbildung A.11: Mixed-Precision-Grenze}
\begin{center}\includegraphics[width=\linewidth,height=0.60\textheight,keepaspectratio]{figures/thesis-page-229.pdf}\end{center}
Originalseite der Dissertation: PDF-Seite 229. Die folgende Interpretation bezieht sich auf die bezeichnete Abbildung.

\textbf{Frage:} Weshalb verbleibt ein erheblicher FP32-Anteil nach Quantization?

\textbf{Antwort:} Nur Wih und Whh werden als INT8 gespeichert. Kopf, Bias, Skalen und sämtliche dynamischen Zustände bleiben Float. Der gestapelte Speichervergleich macht diese Grenze sichtbar. \textbf{Grenze:} Die Grafik betrifft Modellkonstanten, nicht zwangsläufig denselben Gesamtumfang wie Firmware-Flash in Abbildung 7.12. Beide Größen nicht gegeneinander als Widerspruch lesen. Siehe F080 und F081.

\topic{Abbildung A.12: Statische Operationen und gemessene Zeit}
\begin{center}\includegraphics[width=\linewidth,height=0.60\textheight,keepaspectratio]{figures/thesis-page-230.pdf}\end{center}
Originalseite der Dissertation: PDF-Seite 230. Die folgende Interpretation bezieht sich auf die bezeichnete Abbildung.

\textbf{Frage:} Warum ist der statische Quantized/\allowbreak{}Base-Faktor etwa 1,81, aber der Zeitfaktor nicht?

\textbf{Antwort:} Das vereinfachte Zählen zusätzlicher Skalierungs-Multiplikationen bildet nicht alle Instruktionen, Konversionen, Speicherzugriffe, Schleifen und Activation Functions ab. Eine MAC und eine zusätzliche Multiplikation sind auch nicht zwingend zeitlich identisch teuer. \textbf{Grenze:} Ohne Profiling kann die genaue Diskrepanz nicht einer einzelnen Ursache zugeordnet werden. Das Bild belegt gerade die Grenze einer naiven FLOP-zu-Latenz-Umrechnung. Siehe F082 und F083.

\topic{14. Abschließende Selbstprüfung}
Du bist nicht dann vorbereitet, wenn du alle Zahlen auswendig kannst, sondern wenn du jede wichtige Zahl korrekt einordnest. Für die Prüfung sollten mindestens die nominalen SOC-Rangfolgen, die abweichende Recovery-Rangfolge, der Offsetmechanismus, die lange SOH-Filterantwort, die Inferenzmodi und die Mixed-Precision-Grenze ohne Hilfsmittel erklärbar sein.

\textbf{Abschlussaufgabe:} Wähle zufällig eine Abbildung aus Kapitel 5, eine aus Kapitel 6 und eine aus Kapitel 7. Erkläre bei jeder Input, Output, Referenz, Beobachtungseinheit, Intervention, Kennzahl und Grenze. Verbinde sie anschließend mit einer der drei Forschungsfragen. Wenn du dabei unterschiedliche Modelle oder Fehlerdefinitionen vermischst, wiederhole die betreffende Querverbindung statt nur den Plot auswendig zu lernen.

\textbf{Wichtigste sachliche Schlussantwort:} Die Dissertation liefert keine universelle Garantie für neuronale Batteriezustandsschätzung. Sie liefert eine konkrete, messbare Vorgehensweise, um Repräsentation, nominalen Fehler, Störungsmechanismen und Embedded-Ausführung gemeinsam zu beurteilen. Genau diese Verbindung ist verteidigbar, solange die Grenzen der einzelnen Studien transparent bleiben.

\topic{15. Weitere Grundlagenfragen und Tabellenprüfung}
Dieser Teil ergänzt die Fragen F001 bis F100 um häufige Nachfragen zur Methodenwahl und um eine vollständige Orientierung durch die 28 nummerierten Tabellen. Seiten beziehen sich wieder auf die PDF, nicht auf die gedruckte Seitenzählung.

\question{F101. Was unterscheidet ReLU, Sigmoid, tanh und Softmax?}
\textbf{Antwort:} ReLU gibt bei positiven Eingängen den Wert weiter und bei negativen null. Sigmoid bildet auf 0 bis 1 ab und wird für Gates oder begrenzte Ausgaben verwendet. tanh bildet auf -1 bis 1 ab und kann positive und negative Kandidatenzustände darstellen. Softmax normalisiert mehrere Ausgänge auf eine Summe von eins und passt zu konkurrierenden Klassenwahrscheinlichkeiten.

\textbf{Kritische Nachfrage:} Warum wurde Softmax für SOH erwähnt? Bei nur einem Softmax-Ausgang ist die Ausgabe immer eins. Eine sinnvolle Konfiguration müsste sich auf Hidden Layers oder mehrere Ausgänge beziehen. Der Text nennt eine Activation-Vergleichsstudie, aber die konkrete Platzierung muss aus dem ursprünglichen Tuning hervorgehen. Ohne diesen Nachweis nicht behaupten, die Arbeit belege allgemeine Untauglichkeit von Softmax.

\question{F102. Warum überhaupt skalieren und ist RobustScaler eine Robustheitsgarantie?}
\textbf{Antwort:} Skalierung verhindert, dass numerisch große Features wie kumulierte Zeit das Training allein durch ihre Größenordnung dominieren. Min-Max verwendet Bereichsgrenzen, RobustScaler typischerweise Median und Interquartilsabstand. Beide müssen mit ausschließlich Entwicklungsdaten bestimmt und unverändert auf Test und MCU angewendet werden.

\textbf{Grenze:} RobustScaler ist robust gegenüber manchen extremen Trainingswerten bei der Lage-/\allowbreak{}Skalenbestimmung. Er macht das neuronale Modell nicht automatisch robust gegen Sensorfehler oder neue Betriebsbereiche. Ein außerhalb des Trainingsbereichs liegendes Feature bleibt ein Extrapolationsproblem. \textbf{Beleg:} Kapitel 5.2.5 und 7.2.

\question{F103. Warum MSE als Loss, wenn Sie später MAE berichten?}
\textbf{Antwort:} MSE ist glatt und gewichtet große Residuen stärker, was für Training praktisch sein kann. MAE bleibt in derselben Einheit wie die Zielgröße und ist gut interpretierbar. Ein MSE-minimierender Kandidat muss nicht der MAE-Minimierer sein.

\textbf{Vertiefung:} Die Modellwahl sollte das spätere Anwendungsziel spiegeln. Will man insbesondere Maximalfehler, Recovery oder Störungsresistenz verbessern, reicht nominale MSE möglicherweise nicht. Eine geänderte Loss wäre eine neue Trainingsentscheidung und müsste validation-basiert erfolgen. RMSE enthält dieselbe Ordnung wie MSE auf identischen Samples, aber nicht notwendigerweise nach unterschiedlichen Aggregationen über Zellen.

\question{F104. Was unterscheidet AdamW von Adam, und was macht Gradient Clipping?}
\textbf{Antwort:} Adam passt Aktualisierungsschritte anhand von Momentenschätzungen der Gradienten an. AdamW trennt Gewichtszerfall konzeptionell von der adaptiven Gradientenaktualisierung. Gradient Clipping begrenzt große Gradienten, etwa über ihre Norm, und kann instabile Trainingsschritte verhindern.

\textbf{Grenze:} Diese Verfahren garantieren weder Generalisierung noch stabile Onlinezustände. Die Dissertation setzt sie als Trainingswerkzeuge ein, nicht als eigenständige wissenschaftliche Neuerung. Warm Restarts ändern die Lernratenplanung, nicht den Controllerzustand. \textbf{Recherche:} Die in Kapitel 7 zitierten Originalarbeiten von Loshchilov/\allowbreak{}Hutter und Pascanu sowie die tatsächlich gespeicherten Trainingskonfigurationen lesen.

\question{F105. Was unterscheidet globales und lokales Pruning?}
\textbf{Antwort:} Lokales Pruning legt Auswahl oder Rate schichtweise fest. Globales Pruning vergleicht Strukturen über mehrere Schichten und kann Kapazität ungleich verteilen. Scores verschiedener Schichten müssen dafür sinnvoll vergleichbar sein.

\textbf{Grenze:} Die in Kapitel 7 untersuchten Modelle besitzen einen einzelnen recurrent layer und einen spezifischen Gate-Gruppenscore. Daraus folgt keine Aussage über global optimales Pruning eines tiefen Netzes. In den Grundlagen aufgeführte Strategien sind ein Überblick, keine sämtlich durchgeführten Experimente. \textbf{Beleg:} Tabellen 3.3 und 3.4.

\question{F106. Was unterscheidet iterative Auswahl, Pruning-aware Training und Sparse Training?}
\textbf{Antwort:} One-shot entfernt Strukturen einmal, häufig mit anschließendem Fine-Tuning. Iteratives Pruning entfernt mehrfach kleinere Anteile. Pruning-aware Training berücksichtigt die gewünschte Struktur bereits während der Optimierung. Sparse Training kann von Beginn an mit beschränkter Konnektivität arbeiten und Strukturen dynamisch ändern.

\textbf{Grenze:} Diese Ansätze haben verschiedene Trainingskosten und führen nicht automatisch zu gleicher Deploymentrepräsentation. Die Dissertation nutzt in Kapitel 7 einen One-shot-Schnitt mit kurzer Nachanpassung. Nicht versehentlich die Grundlagenübersicht als eigenen Methodenvergleich darstellen.

\question{F107. Was ist der Unterschied zwischen PTQ und QAT?}
\textbf{Antwort:} Post-training quantization transformiert ein bereits trainiertes Modell, gegebenenfalls mithilfe von Kalibrierdaten. Quantization-aware training simuliert relevante Quantization Effects während der Gewichtsanpassung. QAT kann das Netz auf die späteren Rundungs- und Clippingeffekte vorbereiten, verursacht aber zusätzlichen Trainingsaufwand.

\textbf{Grenze:} Die hier untersuchte zeilenweise Max-Skalierung der Gewichte benötigt kein repräsentatives Activation Calibration Set wie eine vollständige Activation Quantization. Sie ist trotzdem nicht allgemein optimal. Die Speicher- und Laufzeitgrenze hängt vom Exportkernel ab, nicht allein von PTQ versus QAT. \textbf{Beleg:} Tabellen 3.6 und 3.7 sowie Kapitel 7.5.

\question{F108. Warum Per-Row statt Per-Tensor oder Per-Group?}
\textbf{Antwort:} Ein gemeinsamer Skalenfaktor für die gesamte Matrix kann von einer großen Ausreißerzeile bestimmt werden und kleine Zeilen grob auflösen. Zeilenweise Skalen passen sich jedem Ausgangskanal an. Kleinere Gruppen erlauben noch feinere Anpassung, benötigen aber mehr Skalen und komplexeren Zugriff.

\textbf{Grenze:} Der hier gewählte Modus ist ein nachvollziehbarer Kompromiss, kein experimenteller Sieg über jede andere Granularität. Für eine Hardwareentscheidung zählen außerdem Skalenanwendung, Vektorisierung und Datentransfer. \textbf{Beleg:} Tabelle 3.6 und Abbildung 7.5.

\question{F109. Warum kann eine einzelne große Gewichtszahl Quantization verschlechtern?}
\textbf{Antwort:} Bei Max-Abs-Skalierung bestimmt das größte Gewicht die Schrittweite der ganzen Zeile. Viele kleine Gewichte können dadurch auf denselben Code oder null runden. Ein größerer darstellbarer Bereich bedeutet bei festem Bitbudget gröbere Auflösung.

\textbf{Vertiefung:} Clipping könnte die Mehrheit genauer repräsentieren, verzerrt aber große Gewichte. Ob das günstiger ist, muss mit Daten und finaler Task-Metrik geprüft werden. Die Dissertation nutzt einen einfach reproduzierbaren Max-Abs-Pfad. Die Halbschrittgrenze gilt für genau diesen Pfad ohne zusätzliches Clipping, nicht für jede Quantization Strategy.

\question{F110. Kann man Pruning und Quantization kombinieren?}
\textbf{Antwort:} Prinzipiell ja, weil sie unterschiedliche Stellgrößen verändern: Topologie und numerische Repräsentation. Man könnte zuerst Kanäle auswählen und fine-tunen und anschließend das kleinere Netz quantisieren. Reihenfolge und gemeinsame Nachanpassung können die Güte beeinflussen.

\textbf{Grenze:} Die berichteten drei Hauptvarianten sind Base, Pruned und Quantized. Ohne eine vierte kombinierte Evaluation darf kein additiver oder multiplikativer Gewinn versprochen werden. Kompressionseffekte können interagieren. Eine neue Variante braucht denselben vollständigen Accuracy-, Transienten-, Speicher- und Laufzeittest.

\question{F111. Warum STM32 statt GPU, FPGA oder ASIC?}
\textbf{Antwort:} Der Zielanwendungsfall benötigt lokale, autonome BMS-Ausführung mit begrenztem Speicher, Kosten und Energie. Ein Mikrocontroller passt zu Messung, Schutz und Kommunikation. GPU, FPGA und ASIC bieten andere Parallelitäts-, Flexibilitäts- und Entwicklungsprofile.

\textbf{Grenze:} Der Vergleich in Tabelle 3.9 ist funktional, keine eigene Messkampagne über diese Plattformen. Eine GPU kann Training beschleunigen, ohne für das spätere Modul geeignet zu sein. Ein FPGA kann deterministische Datenpfade liefern, aber seine Entwicklung und Verifikation sind andere Aufgaben. Keine pauschale Energiehierarchie ohne Arbeitslast und Messung behaupten.

\question{F112. Warum manuelle C-Kerne statt automatischem Konverter?}
\textbf{Antwort:} Sie machen recurrent states, Gateordnung, Skalierung und Quantization-Grenzen explizit und erlauben eine nachvollziehbare Referenzimplementierung. Die Arbeit kann so Operationsfolge und Speicher besser kontrollieren.

\textbf{Grenze:} Manuelle Umsetzung ist fehleranfällig und nicht automatisch schneller. Numerische Äquivalenztests sind daher notwendig. Automatische Toolchains können hochoptimierte Operatoren bereitstellen, deren Zustandssemantik und Modellabdeckung aber geprüft werden müssen. Tabellen 3.5 und 3.8 sind Softwareübersichten, keine Garantie der Unterstützung genau des eigenen Modellgraphen in jeder Version.

\question{F113. Was ist ein sinnvoller Fallback für ein neuronales BMS-Modell?}
\textbf{Antwort:} Ein unabhängig überprüfter einfacher Schätzer, Plausibilitätsgrenzen und konservative Betriebsgrenzen könnten bei unplausiblen Eingängen oder Ausgaben aktiviert werden. Der Wechsel müsste Zustandskonsistenz und Übergangssprünge berücksichtigen.

\textbf{Grenze:} Das ist ein Entwicklungsvorschlag und kein bereits validiertes Ergebnis der Dissertation. Ein Fallback-Coulomb-Counter kann denselben Stromsensorfehler teilen. Diversität der Modelle allein garantiert daher keine unabhängige Fehlerabsicherung. Schutzfunktionen müssen unabhängig von einer optimistischen SOC-Prognose wirken.

\question{F114. Welche zusätzliche Untersuchung hätte den größten Erkenntnisgewinn?}
\textbf{Antwort:} Je nach Ziel: Für Klassenübertragbarkeit mehr unabhängige Zellen und Betriebsdomänen. Für schnelle DD-Deploymentaussagen die wichtigsten Störungstests im Continuous-Modus. Für die SOH-Netzleistung eine eindeutig dokumentierte Raw-/\allowbreak{}Filter-Ablation mit bekannten Anfangsbedingungen. Für Energie echte synchronisierte Leistungsmessung. Für INT8 optimierte, identisch instrumentierte Kernel.

\textbf{Priorisierung:} Nicht alles gleichzeitig versprechen. Eine gute Prüfungsantwort nennt den konkreten offenen Schluss und den kleinsten Versuch, der ihn beantwortet. Für den behaupteten Übergang vom robusten Softwaremodell zum schnellen Controller ist die Prüfung der Inferenzsemantik besonders unmittelbar.

\question{Tabellen 1.1 und 8.1: Beitrag und Synthese}
\textbf{PDF-Seiten 33 und 185.} Tabelle 1.1 ordnet Studien, Beiträge und Publikationsbasis zu. Tabelle 8.1 verknüpft die Beiträge mit Anforderungen. \textbf{Prüfungsfrage:} Welche Zelle der Synthesetabelle ist direkt gemessen und welche indirekt relevant? \textbf{Antwort:} MLP-Repräsentation ist primär eine Schätzstudie, kein dortiger vollständiger MCU-Test. Kapitel 6 und 7 enthalten direkte Kernmessungen, Energie in Kapitel 7 aber nur Proxy. \textbf{Offen:} Publikationsstatus ist zeitabhängig. Für die tatsächliche Verteidigung den dann gültigen Status prüfen, nicht aus der September-PDF dauerhaft übernehmen.

\question{Tabellen 3.1 und 3.2: Chemie und Schätzerprinzipien}
\textbf{PDF-Seiten 40 und 52.} Die Tabellen bündeln typische Eigenschaften von NMC/\allowbreak{}LFP sowie Kalman-/\allowbreak{}Neuronalschätzung. \textbf{Prüfungsfrage:} Sind das harte Ausschlusskriterien? \textbf{Antwort:} Nein, es sind typische Stärken und Grenzen unter den jeweils notwendigen Voraussetzungen. Ein neuronales Netz kann Physikfeatures nutzen, ein Observer adaptive oder gelernte Komponenten. Beobachtbarkeit und Trainingsabdeckung sind unterschiedliche Arten begrenzter Information. \textbf{Lernauftrag:} Zu jeder behaupteten Stärke ein Gegenbeispiel oder eine Bedingung nennen.

\question{Tabellen 3.3 bis 3.5: Pruningmethoden und Werkzeuge}
\textbf{PDF-Seiten 56, 61 und 63.} Granularität, Optimierungsverfahren und Softwareunterstützung sind drei unterschiedliche Achsen. \textbf{Prüfungsfrage:} Erzeugt eine hohe Sparse-Rate schon eine schnelle Firmware? \textbf{Antwort:} Nein, Format, Operator und Kernel müssen die Sparsity ausnutzen. \textbf{Lernauftrag:} One-shot/\allowbreak{}local/\allowbreak{}gate-group/\allowbreak{}structured/\allowbreak{}dense-export als genaue Beschreibung der eigenen Kapitel-7-Variante zusammensetzen. Frameworkfähigkeiten sind versionsabhängig und dürfen nicht ungeprüft als aktuelle Vollabdeckung ausgegeben werden. Siehe F105, F106 und F112.

\question{Tabellen 3.6 bis 3.8: Quantization Methods und Werkzeuge}
\textbf{PDF-Seiten 68, 70 und 72.} Die Tabellen behandeln Granularität, Strategie und Software. \textbf{Prüfungsfrage:} An welcher Stelle liegen Kalibrierungsdaten, Rundung, Integerarithmetik und Deploymentoperator? \textbf{Antwort:} Das sind getrennte Designentscheidungen. Der eigene Versuch ist zeilenweise symmetrische PTQ nur für Recurrent-Gewichte mit FP32-Zuständen. \textbf{Grenze:} Nicht alle aufgeführten Modi wurden selbst gemessen. Siehe F107 bis F109.

\question{Tabelle 3.9: Rechenplattformen}
\textbf{PDF-Seite 75.} \textbf{Prüfungsfrage:} Welche Plattform würden Sie bei härteren Latenzgrenzen wählen? \textbf{Antwort:} Erst Rechenlast, Parallelität, Energie, Stückzahl, Entwicklungsaufwand und Sicherheitsanforderungen festlegen. Die Dissertation misst nur die gewählten STM32-Pfade. Ein Tabellenvergleich der Prinzipien ersetzt keine Portierungsstudie. Siehe F111.

\question{Tabellen 4.1 bis 4.3: Daten, NMC-Szenarien und Zellgrenzen}
\textbf{PDF-Seiten 80, 81 und 86.} \textbf{Prüfungsfrage:} Sind Zellen, Szenarien und Samples dasselbe n? \textbf{Antwort:} Nein. 14 NMC-Zellen stehen sieben Szenarien gegenüber. Beim LFP-Test sind sechs Zellen die primären unabhängigen Einheiten, viele Samples beschreiben deren Verlauf. Die vollständige Split-Tabelle ist verbindlicher als eine unklare Farblegende. \textbf{Lernauftrag:} Für C09/\allowbreak{}C13/\allowbreak{}C15 ähnliche mittlere Lastgruppe, für C29 explorative High-Gruppe und für C27 fehlende Alterungsfenster erklären. Siehe F016 bis F025.

\question{Tabelle 5.1: Finale MLP-Größe}
\textbf{PDF-Seite 105.} \textbf{Prüfungsfrage:} Wie viele trainierbare Parameter hat diese Tabelle? \textbf{Antwort:} Erst tatsächliche Eingangsgröße klären, dann pro Schicht Eingänge mal Ausgänge plus Bias summieren. Bei 48 Inputs ergibt die gedruckte Breitenfolge 434561 Parameter. \textbf{Grenze:} Das ist eine bedingte Rechnung, keine Zählung aus einem geladenen Checkpoint. Die schematische Grafik und der anfängliche Acht-Neuronen-Kandidat sind kein Ersatz für die Tabelle. Siehe F030.

\question{Tabellen 6.1 und 6.2: Repräsentanten und Interventionen}
\textbf{PDF-Seiten 115 und 121 bis 123.} \textbf{Prüfungsfrage:} Welche Teile der Klasse wurden nicht untersucht? \textbf{Antwort:} Zum Beispiel alternative ECM-Ordnungen, Hysteresezustände, adaptive Kovarianzen, andere neuronale Architekturen und kombinierte Sensorfehler. \textbf{Lernauftrag:} Jede Interventionsstärke mit Einheit nennen und physikalischen Gain von Offset unterscheiden. Die Parameter sind kontrollierte Stressdefinitionen, keine geschätzte Wahrscheinlichkeit realer Sensorfehler. Siehe F037 bis F060.

\question{Tabelle 6.3: Empfehlung nach Priorität}
\textbf{PDF-Seite 150.} \textbf{Prüfungsfrage:} Welchen Vertreter würden Sie tatsächlich einsetzen? \textbf{Antwort:} Bei primärer nominaler Genauigkeit und geringer systematischer Stromfehlersensitivität spricht der Versuch für DD. Bei Recovery HECM, bei minimalem Rechenbudget DM/\allowbreak{}HDM. Die schnelle DD-Empfehlung verwendet Continuous und verlangt zusätzliche Störungsvalidierung dieser Semantik. \textbf{Grenze:} Eine Prioritätsempfehlung ist keine pauschale Systemfreigabe.

\question{Tabellen 7.1 und 7.2: Hardware-KPIs}
\textbf{PDF-Seite 167.} \textbf{Prüfungsfrage:} Sind Latency und Inference doppelte Messungen? \textbf{Antwort:} Nein, Host Latency und Kernzeit haben unterschiedliche Grenzen. Flash ist gelinkter Speicher, RAM enthält statische Daten und beobachteten Stack, Eest ist aus Inference berechnet. \textbf{Lernauftrag:} SOC Base/\allowbreak{}Pruned/\allowbreak{}Quantized mit 1,40/\allowbreak{}0,80/\allowbreak{}6,99 ms und SOH mit 22,73/\allowbreak{}12,72/\allowbreak{}29,21 ms erklären. Keine unabhängige Energiebestätigung aus der letzten Spalte ableiten. Siehe F087 bis F091.

\question{Tabelle 7.3: Kompressionstrade-off}
\textbf{PDF-Seite 178.} \textbf{Prüfungsfrage:} Warum ist negativer ΔMAE gut, negativer ΔFlash ebenfalls gut und U kleiner eins günstig? \textbf{Antwort:} Es sind Änderungen gegenüber Base beziehungsweise normierte Kosten. Die Vorzeichen verschiedener Tabellen dürfen nicht mit den höher-ist-besser-Robustheitsscores aus Kapitel 6 verwechselt werden. \textbf{Grenze:} Gerundete Tabellenwerte können kleine Unterschiede zur Rechnung aus ungerundeten Daten erzeugen. Mehrere Prozentänderungen sind nicht einfach zu addieren.

\question{Tabelle A.1: Lookupvariation}
\textbf{PDF-Seite 219.} \textbf{Prüfungsfrage:} Was bedeuten eckige Angaben wie 6/\allowbreak{}22? \textbf{Antwort:} Die Caption definiert Anzahlen der Subfälle mit null-ausschließendem Interaktionsintervall, nicht Grenzen eines einzigen Intervalls. Der Spaltenkopf mit CI kann deshalb missverständlich wirken. \textbf{Wichtig:} Die Tabelle enthält tatsächlich nicht nur Widerstand und OCV, sondern auch Zeitkonstanten. -10 Prozent Widerstand erzeugt 5,10 h Recovery-or-censor und 17 Prozent zensierte Zellen. Siehe F059 und F060.

\question{Tabelle A.2: Nominale Unsicherheiten}
\textbf{PDF-Seite 219.} \textbf{Prüfungsfrage:} Was ist P95 und warum nicht derselbe Wert wie RMSE? \textbf{Antwort:} P95 ist ein Quantil der absoluten Fehler, RMSE ein quadratisches Mittel aller Fehler. Beide beschreiben verschiedene Aspekte. DD-P95 beträgt hier 0,0508, während MAE 0,0258 ist. \textbf{Grenze:} Ein quantilbezogener Makrowert ist nicht automatisch das Quantil aller gepoolten Samples. Aggregationsreihenfolge beachten.

\question{Tabelle A.3: Gepaarte Tests}
\textbf{PDF-Seite 220.} \textbf{Prüfungsfrage:} Warum durchweg pHolm = 0,1875? \textbf{Antwort:} Diskrete exakte Tests mit sechs Zellen und Korrektur über sechs Vergleiche liefern eine grobe Ergebnisauflösung. Die Tabelle enthält etwa HECM gegen DD mit Mittelwertdifferenz -0,0079 und einem punktweisen Intervall unter null, aber exaktem p = 0,0938. \textbf{Grenze:} Standardisiertes dz kann bei kleiner Stichprobe instabil sein. Rohe Differenzen und Einheiten ebenfalls nennen. Siehe F063 und F064.

\question{Tabelle A.4: Ausgewählte Störungspenalties}
\textbf{PDF-Seite 220.} \textbf{Prüfungsfrage:} Ist -0,0000 eine reale Verbesserung? \textbf{Antwort:} Es ist ein gerundeter sehr kleiner negativer Wert. Für Interpretation ungerundete Daten und Unsicherheit prüfen. Die Tabelle fasst globale MAE-Kontraste zusammen und ersetzt keine lokale Spikeanalyse. Die Offsetwerte sind zusätzliche MAE, nicht gesamte MAE. Siehe F048 und F053.

\question{Tabelle A.5: Score-Sensitivität}
\textbf{PDF-Seite 220.} \textbf{Prüfungsfrage:} Warum ändern sich Abstände stark mit Gewichtung? \textbf{Antwort:} Familien, Severity-Level und Einzelszenarien erhalten unterschiedliche relative Bedeutung. DD bleibt im dargestellten Kandidatensatz führend, was begrenzte Rankingstabilität zeigt. \textbf{Grenze:} Kein frei gewähltes Gewichtungssystem ist eine gemessene Häufigkeit von Feldfehlern. Siehe F068.

\question{Tabelle A.6: Recovery und Hardware zusammen}
\textbf{PDF-Seite 221.} \textbf{Prüfungsfrage:} Warum ist DM-Relapse null, obwohl DM nicht recovern kann? \textbf{Antwort:} Ohne qualifizierten Bandeintritt gibt es auch keinen anschließenden Rückfall. Niedriger Relapse allein ist deshalb nicht automatisch gut. DM/\allowbreak{}HDM haben hohe Zensur. \textbf{Zweiter Prüfpunkt:} Frühe bereits erfüllte Endpunkte sind konservative obere Zeitgrenzen, weil der gemeinsame Scorebeginn später liegt. DD-Hardwarewerte betreffen Continuous. Siehe F042 bis F046 und F085.

\question{Tabelle A.7: Fehleränderung nach Bitflip}
\textbf{PDF-Seite 226.} \textbf{Prüfungsfrage:} Warum ist der SOH-Median negativ und trotzdem kein Robustheitsgewinn? \textbf{Antwort:} Ein einzelner Fault kann eine bestehende Zielabweichung zufällig kompensieren. Die gepaarte Outputabweichung, Peak und Recovery bleiben zusätzliche Kriterien. P95 der ereignisweisen MAE-Änderung ist nicht P95 aller instantaneous errors. \textbf{Grenze:} Ein bestimmtes Fraction-Bit und wenige Ereignisse erfassen keine vollständige Fehlertoleranz. Siehe F097.

\topic{Modellübersicht: Architektur, Training und Ergebnisse}
\subsection*{MODELLÜBERSICHT · 1/\allowbreak{}5}
\section*{\color{accent}Welche Modelle habe ich benutzt?}
Lernübersicht zu den Kapiteln 5, 6 und 7. Kapitel 8 fasst die Ergebnisse zusammen. Die drei Untersuchungen beantworten unterschiedliche Fragen.

{\fontsize{9.5}{12}\selectfont\setlength{\tabcolsep}{2mm}\renewcommand{\arraystretch}{1.2}\begin{longtable}{>{\raggedright\arraybackslash}p{35.47mm}>{\raggedright\arraybackslash}p{126.53mm}}
\toprule
\textbf{Kapitel /\allowbreak{} Modell} & \textbf{Aufbau und Aufgabe} \\
\midrule\endhead
5 · SOH-MLP & Feedforward-Netz ohne rekurrentes Gedächtnis. Spannung, Temperatur und kumulierter Strom werden als Lag-Vektor eingegeben. Neun Dense Hidden Layers mit ReLU, danach ein linearer SOH-Ausgang. Untersucht den Nutzen zeitlicher Input Features. \\
6 · DM → SOC & Coulomb Counting mit fester Nennkapazität; Stromintegration und Rücksetzen bei erkannter Vollladung. Kein neuronales SOC-Netz. \\
6 · HDM → SOC & Wie DM, aber effektive Kapazität = Nennkapazität × geschätzter SOH. Nutzt das gemeinsame SOH-LSTM. \\
6 · HECM → SOC & Ersatzschaltbild mit zwei RC-Zweigen und Kalman-Korrektur anhand der Spannung. Zustände: SOC und zwei Polarisationsspannungen. Tabellenparameter werden nach SOC und SOH angepasst. \\
6 · DD /\allowbreak{} GRU → SOC & Eine GRU-Schicht: ursprünglich 96, nach Structured Pruning 67 Einheiten. MLP-Kopf mit 96 Hidden Units und einem Sigmoid-Ausgang. Acht Input Features. \\
6 · gemeinsames SOH-LSTM & 20 stündliche Features → zweistufiges Embedding (Breite 128) → 2 LSTM-Schichten mit je 112 Einheiten → 3 residuale MLP-Blöcke → Kopf mit Breite 160 → SOH. Liefert SOH für HDM, HECM und DD. \\
7 · SOC-LSTM + MLP & Base: 1 LSTM-Schicht, Hidden Size 64; MLP-Kopf 64 → 1 mit ReLU/\allowbreak{}Sigmoid. Sechs Input Features, sekündliche Verarbeitung. \\
7 · SOH-LSTM + MLP & Base: 1 LSTM-Schicht, Hidden Size 128; MLP-Kopf 128 → 1 mit ReLU/\allowbreak{}linearem Ausgang. Sechs Input Features, sekündliche Verarbeitung. \\
\bottomrule
\end{longtable}}
MLP-Schichtbreiten in Kapitel 5: 128 → 256 → 256 → 256 → 128 → 512 → 128 → 256 → 256 → 1. Bei 16 Lag-Schritten × 3 Features entstehen 48 Eingabewerte.

Kapitel 7 vergleicht je Aufgabe Base (FP32), strukturell Pruned (FP32) und Quantized (rekurrente Gewichte INT8). Quantized ist kein vollständig ganzzahliges Netz; Zustände und weitere Berechnungen bleiben in Gleitkomma. Quellen: [1], [2], [3], [4].


\clearpage
\subsection*{MODELLÜBERSICHT · 2/\allowbreak{}5}
\section*{\color{accent}Training, Fenster und Daten}
Window Length = Zeitspanne eines Training Sample. Batch Size = Anzahl solcher Beispiele in einem Verarbeitungsschritt. Beides ist von der Zahl der Zellen zu unterscheiden.

{\fontsize{9.5}{12}\selectfont\setlength{\tabcolsep}{2mm}\renewcommand{\arraystretch}{1.2}\begin{longtable}{>{\raggedright\arraybackslash}p{25.51mm}>{\raggedright\arraybackslash}p{54.07mm}>{\raggedright\arraybackslash}p{50.12mm}>{\raggedright\arraybackslash}p{24.3mm}}
\toprule
\textbf{Modell} & \textbf{Eingaben pro Zeitschritt} & \textbf{Training Window} & \textbf{Batch} \\
\midrule\endhead
Kap. 5 SOH-MLP & U, T, kumulierter Strom; flach als Lag-Vektor & Beste berichtete Kombination: 16 × 10 min = 160 min (2 h 40 min) & Nicht bestätigt; gefundenes MLP-Skript derzeit nur als Cloud-Datei verfügbar \\
Kap. 6 SOC-GRU /\allowbreak{} DD & U, I, T, SOH, Qc, dU/\allowbreak{}dt, dI/\allowbreak{}dt, Δt & 2.024 Punkte bei nominal 1 Hz; rund 34 min & 512; je 2 Batches pro Update laut Konfiguration \\
Kap. 6 SOH-LSTM & U, I, T, EFC, Qc: jeweils Mittelwert, Std., Min., Max. pro Stunde = 20 Features & 192 Stundenwerte = 8 Tage & 128 \\
Kap. 7 SOC-LSTM & U, I, T, Qc, dU/\allowbreak{}dt, dI/\allowbreak{}dt & 2.024 Punkte bei 1 Hz; rund 34 min & 512 \\
Kap. 7 SOH-LSTM & Zeit, U, I, T, EFC, Qc & 2.048 Punkte bei 1 Hz; rund 34 min (0,57 h) & 128 \\
\bottomrule
\end{longtable}}
\subsection*{Welche Daten gehören wohin?}
{\fontsize{9.5}{12}\selectfont\setlength{\tabcolsep}{2mm}\renewcommand{\arraystretch}{1.2}\begin{longtable}{>{\raggedright\arraybackslash}p{26.84mm}>{\raggedright\arraybackslash}p{135.16mm}}
\toprule
\textbf{Studie} & \textbf{Datengrundlage und Aufteilung} \\
\midrule\endhead
Kapitel 5 & NMC-Datensatz, 14 Zellen. Vier vollständige Zellen für die unabhängige Prüfung zurückgehalten; 20 \% der Training Sequences für Validierung. Rund 15.000 Training Sequences im beschriebenen Tuning. \\
Kapitel 6 & LFP DoE Cycle Ageing Dataset. Training: C01, C03, C05, C11, C17, C23. Validierung: C07, C19, C21. Unabhängiger Holdout: C09, C13, C15, C25, C27, C29. SOC-Baseline: 16 ausgewählte Fenster aus diesen sechs Testzellen. \\
Kapitel 7 & LFP-Daten. Die geprüften SOC-/\allowbreak{}SOH-Konfigurationen verwenden dieselben sechs Trainings- und drei Validierungszellen wie oben. Ergebnisse aus vollständigem Streaming-Replay, nicht aus denselben 16 Robustheitsfenstern. Die lokale SOH-Filteranalyse nutzt C07, das in diesen Konfigurationen zur Validierung gehört. \\
\bottomrule
\end{longtable}}
U = Spannung; I = Strom; T = Temperatur; Qc = kumulierter Ladungszustand der Integration; EFC = äquivalente Vollzyklen; Δt = Zeitabstand. Fensterdauern sind die üblichen nominalen Angaben N × Abtastintervall. Batchwerte stammen aus gespeicherten Konfigurationen, nicht aus einem neuen Training. DM/\allowbreak{}HDM/\allowbreak{}HECM besitzen keine eigene neuronale SOC-Batch Size. Quellen: [1]-[4], [6].


\clearpage
\subsection*{MODELLÜBERSICHT · 3/\allowbreak{}5}
\section*{\color{accent}Was kam bei SOC und MLP heraus?}
MAE und RMSE sind hier einheitlich in Prozentpunkten angegeben. Beispiel: 0,0258 auf der 0-bis-1-Skala entspricht 2,58 Prozentpunkten, nicht 97,42 \% Klassifikationsgenauigkeit.

\subsection*{Kapitel 6 · SOC ohne Messstörung}
{\fontsize{9.5}{12}\selectfont\setlength{\tabcolsep}{2mm}\renewcommand{\arraystretch}{1.2}\begin{longtable}{>{\raggedright\arraybackslash}p{27.34mm}>{\raggedright\arraybackslash}p{16.71mm}>{\raggedright\arraybackslash}p{16.71mm}>{\raggedright\arraybackslash}p{93.25mm}}
\toprule
\textbf{Modell} & \textbf{MAE} & \textbf{RMSE} & \textbf{Einordnung} \\
\midrule\endhead
DM & 6,90 & 7,40 & Feste Kapazität; größter mittlerer Baseline-Fehler. \\
HDM & 4,12 & 4,42 & SOH-basierte Kapazitätsanpassung verbessert den beobachteten Mittelwert. \\
HECM & 3,37 & 3,88 & Spannungsrückführung und dynamisches Ersatzschaltbild. \\
DD /\allowbreak{} GRU & 2,58 & 3,00 & Kleinster beobachteter mittlerer Baseline-Fehler. \\
\bottomrule
\end{longtable}}
Auswertung: erst Fenster je Zelle mitteln, dann die sechs Zellen gleich gewichten. Diese Rangfolge beschreibt die ungestörte Baseline. Sie gilt nicht automatisch für alle Störungen. Der paarweise Sign-Flip-Test mit Holm-Korrektur liefert hier keinen Signifikanznachweis.

\subsection*{Kapitel 7 · SOC im vollständigen Streaming-Replay}
{\fontsize{9.5}{12}\selectfont\setlength{\tabcolsep}{2mm}\renewcommand{\arraystretch}{1.2}\begin{longtable}{>{\raggedright\arraybackslash}p{27.34mm}>{\raggedright\arraybackslash}p{16.71mm}>{\raggedright\arraybackslash}p{16.71mm}>{\raggedright\arraybackslash}p{93.25mm}}
\toprule
\textbf{Variante} & \textbf{MAE} & \textbf{RMSE} & \textbf{Ergebnis} \\
\midrule\endhead
Base FP32 & 2,68 & 3,81 & Unkomprimiertes Referenznetz. \\
Pruned FP32 & 2,34 & 3,13 & In diesem Lauf kleinster Fehler; zusätzlich weniger Speicher und kürzere Kernel Time. \\
Quantized & 2,79 & 3,88 & Fehler ähnlich Base; geringerer Flash-Bedarf, aber hier längere Kernel Time. \\
\bottomrule
\end{longtable}}
\subsection*{Kapitel 5 · SOH mit MLP und Lag-Features}
{\fontsize{9.5}{12}\selectfont\setlength{\tabcolsep}{2mm}\renewcommand{\arraystretch}{1.2}\begin{longtable}{>{\raggedright\arraybackslash}p{43.13mm}>{\raggedright\arraybackslash}p{21.26mm}>{\raggedright\arraybackslash}p{22.78mm}>{\raggedright\arraybackslash}p{66.82mm}}
\toprule
\textbf{Ergebnis} & \textbf{MAE} & \textbf{MSE (pp²)} & \textbf{Bedeutung} \\
\midrule\endhead
Berichteter Durchschnitt & 1,10 & 2,95 & Mittlerer Fehler über die ausgewerteten Fälle. \\
Günstige Einzelfälle & 0,487 /\allowbreak{} 0,253 & 0,410 /\allowbreak{} 0,118 & Gute Ergebnisse einzelner Zellen, keine allgemeine Genauigkeitsgarantie. \\
Atypische Zelle 9 & 2,528 & 8,327 & Deutlich größere Abweichungen vom Referenz-SOH. \\
\bottomrule
\end{longtable}}
MSE ist der mittlere quadratische Fehler, nicht RMSE. Die MLP-Verlaufsbilder enthalten eine Rolling-Glättung; die hier gezeigten Zahlen werden so aus Abschnitt 5.5 übernommen. Der genaue Filterbezug dieser Kennzahlen ist damit nicht zusätzlich auditiert. Kein Architektur-Ranking zwischen Kapiteln: Daten, Ziele, Nachverarbeitung und Gewichtung unterscheiden sich. Quellen: [1], Abschnitte 5.5 und 7.7, Tabelle A.2.


\clearpage
\subsection*{MODELLÜBERSICHT · 4/\allowbreak{}5}
\section*{\color{accent}SOH: Ergebnisse sauber auseinanderhalten}
Die Zahlen 0,85 und 1,48 stammen aus unterschiedlichen Ausführungen. „Vor Filterung“ und „nach Filterung“ dürfen nur innerhalb derselben Rechnung direkt verglichen werden.

\subsection*{Robustheitsmodell · stündliches SOH-LSTM}
{\fontsize{9.5}{12}\selectfont\setlength{\tabcolsep}{2mm}\renewcommand{\arraystretch}{1.2}\begin{longtable}{>{\raggedright\arraybackslash}p{111.25mm}>{\raggedright\arraybackslash}p{23.37mm}>{\raggedright\arraybackslash}p{23.37mm}}
\toprule
\textbf{Auswertung laut Modelldokumentation} & \textbf{MAE} & \textbf{RMSE} \\
\midrule\endhead
Ungewichteter Mittelwert der Kennzahlen über sechs Holdout-Zellen & 2,4945 & 3,1557 \\
Alle 13.038 Holdout-Samples zusammengefasst (andere Gewichtung) & 2,1792 & 3,2326 \\
\bottomrule
\end{longtable}}
\subsection*{Embedded · ursprünglich gespeicherter Benchmark}
{\fontsize{9.5}{12}\selectfont\setlength{\tabcolsep}{2mm}\renewcommand{\arraystretch}{1.2}\begin{longtable}{>{\raggedright\arraybackslash}p{26.43mm}>{\raggedright\arraybackslash}p{18.22mm}>{\raggedright\arraybackslash}p{18.22mm}>{\raggedright\arraybackslash}p{91.12mm}}
\toprule
\textbf{Variante} & \textbf{MAE} & \textbf{RMSE} & \textbf{Ausgabestand} \\
\midrule\endhead
Base & 0,85 & 1,15 & Gespeicherte Endausgabe nach Nachverarbeitung \\
Pruned & 1,46 & 1,87 & Gespeicherte Endausgabe nach Nachverarbeitung \\
Quantized & 1,41 & 1,89 & Gespeicherte Endausgabe nach Nachverarbeitung \\
\bottomrule
\end{longtable}}
\subsection*{Embedded · separate lokale Filteranalyse, gleicher C07-Verlauf}
{\fontsize{9.5}{12}\selectfont\setlength{\tabcolsep}{2mm}\renewcommand{\arraystretch}{1.2}\begin{longtable}{>{\raggedright\arraybackslash}p{85.66mm}>{\raggedright\arraybackslash}p{22.78mm}>{\raggedright\arraybackslash}p{22.78mm}>{\raggedright\arraybackslash}p{22.78mm}}
\toprule
\textbf{MAE nach Verarbeitungsstand} & \textbf{Base} & \textbf{Pruned} & \textbf{Quantized} \\
\midrule\endhead
Vor zeitlicher Filterung; Startwert bereits ausgerichtet & 1,8496 & 1,9847 & 2,1064 \\
Nach Filterstufe 1 & 1,6771 & 1,6436 & 1,9259 \\
Nach Filterstufen 1 und 2 & 1,4763 & 1,6946 & 1,0279 \\
\bottomrule
\end{longtable}}
Belastbare Aussage: Im lokalen Base-Lauf verändert die Filterkette den MAE von 1,8496 über 1,6771 auf 1,4763 Prozentpunkte. Die 0,85 sind nicht der Endwert dieser Rechnung.

Offener Prüfpunkt: Die lokale Wiederholungsrechnung reproduziert den ursprünglichen Verlauf nicht exakt. Die Unterlagen dokumentieren außerdem unterschiedliche historische Filterdefinitionen. Die tatsächlichen Filterketten, Datenreihenfolge und numerischen Ausführungen beider Läufe müssen abgeglichen werden; die Ursache der Abweichung ist nicht isoliert nachgewiesen.

Alle Werte in Prozentpunkten. Lokaler Vergleich: 14.496.990 Samples nach 2.048 Warm-up-Punkten, mit Startwertskalierung; kein erneutes Training und kein neuer STM32-Lauf. Ein kleinerer Wert allein beweist keine Überlegenheit sekündlicher gegenüber stündlichen Eingaben. Quellen: [1], Abschnitte 7.7/\allowbreak{}7.8 und Abb. A.6; [3], [5], [6].


\clearpage
\subsection*{MODELLÜBERSICHT · 5/\allowbreak{}5}
\section*{\color{accent}So erkläre ich es in der Prüfung}
Training, gespeichertes Gedächtnis und Auswertung sind drei verschiedene Ebenen. Diese Begriffe helfen, die Tabellen richtig einzuordnen.

{\fontsize{9.5}{12}\selectfont\setlength{\tabcolsep}{2mm}\renewcommand{\arraystretch}{1.2}\begin{longtable}{>{\raggedright\arraybackslash}p{35.79mm}>{\raggedright\arraybackslash}p{126.21mm}}
\toprule
\textbf{Begriff} & \textbf{Bedeutung in deinen Modellen} \\
\midrule\endhead
Training Window /\allowbreak{} Chunk & Ein zusammenhängender Zeitabschnitt einer Zelle. Beim Embedded-LSTM bleiben seine Messpunkte chronologisch; die Fenster werden über Trainingszellen hinweg gemischt. \\
Batch & Mehrere Fenster, gegebenenfalls aus verschiedenen Trainingszellen. Jedes Fenster besitzt seinen eigenen Zustand. Beim Embedded-SOC: 512 Fenster; beim Embedded-SOH: 128. \\
Epoche & Ein Durchgang durch die Training Windows. Bei den geprüften Embedded-Loadern wird ein unvollständiger letzter Batch weggelassen. Validation Data und Test Data gehören nicht zur Trainingsepoche. \\
Gedächtnis im Training & Die geprüften Kapitel-7-Skripte starten Hidden State und Cell State für jedes Fenster neu. Die gelernten Gewichte bleiben erhalten. Backpropagation reicht innerhalb des Fensters zurück. \\
Rolling Window im Benchmark & Kapitel-6-DD: für jede Vorhersage die letzten 2.024 Punkte erneut verarbeiten; Zustand zwischen Fenstern zurücksetzen. Das ist kein dauerhaft weitergereichtes Gedächtnis. \\
Stateful Streaming & Kapitel-7-LSTMs: Zustände im Einsatz zwischen Messpunkten weitergeben. Kapitel-6-SOH-LSTM: Zustände zwischen stündlichen Updates weitergeben. Trainingsfensterlänge ist keine feste Gedächtnis-Löschfrist im Einsatz. \\
\bottomrule
\end{longtable}}
\textbf{Merksatz:} Kapitel 5: MLP mit expliziter Vergangenheit. Kapitel 6: Robustheitsvergleich mit SOC-GRU und gemeinsamem stündlichem SOH-LSTM. Kapitel 7: zwei sekündliche LSTMs für den Embedded-Kompressionsvergleich.

\subsection*{Quellen und Prüfstand}\begingroup\fontsize{9.5}{11.5}\selectfont\setlength{\parskip}{4pt}
[1] Dissertation main(1).pdf: Tabelle 5.1; Abschnitte 5.2-5.5, 6.2.2-6.2.3, 7.2-7.3 und 7.7-7.8; Tabelle A.2; Abbildung A.6.

[2] LFP\_\allowbreak{}SOC\_\allowbreak{}SOH\_\allowbreak{}Model/\allowbreak{}2\_\allowbreak{}models/\allowbreak{}SOC\_\allowbreak{}1.7.0.0: train\_\allowbreak{}soc.yaml sowie PrunedFT\_\allowbreak{}1.7.0.0\_\allowbreak{}s30\_\allowbreak{}struct/\allowbreak{}config/\allowbreak{}train\_\allowbreak{}soc.yaml. Architektur, Batch und Akkumulation aus gespeicherten Konfigurationen.

[3] LFP\_\allowbreak{}SOC\_\allowbreak{}SOH\_\allowbreak{}Model/\allowbreak{}2\_\allowbreak{}models/\allowbreak{}SOH\_\allowbreak{}JES2\_\allowbreak{}0.1.0: config/\allowbreak{}train\_\allowbreak{}soh.yaml und README.md. Fenster 192 Stunden; Architektur, Zellaufteilung und SOH-Holdout-Metriken.

[4] LFP\_\allowbreak{}LSTM\_\allowbreak{}MLP/\allowbreak{}1\_\allowbreak{}training: 1.5.0.0/\allowbreak{}config/\allowbreak{}train\_\allowbreak{}soc.yaml und scripts/\allowbreak{}train\_\allowbreak{}soc.py; 2.1.0.0/\allowbreak{}config/\allowbreak{}train\_\allowbreak{}soh.yaml und scripts/\allowbreak{}train\_\allowbreak{}soh.py. Belegt Fenster, Batch, gemischte Reihenfolge und fehlende Zustandsübergabe zwischen Training Windows.

[5] LATEX/\allowbreak{}EAAI/\allowbreak{}embedded\_\allowbreak{}AI\_\allowbreak{}optimization/\allowbreak{}review\_\allowbreak{}1/\allowbreak{}review\_\allowbreak{}analysis: KEY\_\allowbreak{}RESULTS.md und results/\allowbreak{}filter/\allowbreak{}soh\_\allowbreak{}filter\_\allowbreak{}compression\_\allowbreak{}local\_\allowbreak{}windows.csv. Für die Endfilterkette gilt die CSV-Zeile FinalSequential; ältere Zusammenfassungen enthalten auch alternative Filtervarianten.

[6] Derselbe review\_\allowbreak{}analysis-Ordner: LOCAL\_\allowbreak{}WINDOWS\_\allowbreak{}REEXECUTION.md und results/\allowbreak{}analysis\_\allowbreak{}provenance.json. Kennzeichnen die lokale Wiederholungsrechnung und die Unterschiede zur archivierten Auswertung.

Zusätzliche MLP-Codeprüfung: Im Ordner TUB/\allowbreak{}7\_\allowbreak{}Codes/\allowbreak{}MLP wurde SOH\_\allowbreak{}sliding\_\allowbreak{}window.py gefunden. Der Inhalt war wegen des nicht laufenden Synology-Clouddienstes nicht lesbar; daher ist weder seine Zuordnung zum finalen Modell noch eine Batch Size daraus bestätigt.

Quellen [2]-[6] relativ zu 1\_\allowbreak{}Scripts. Ergänzung geprüft am 23.09.2026. Vorhandene Fragenkatalogseiten bleiben erhalten. Kein neues Training oder Benchmark für diese Übersicht.

\endgroup\topic{Englische Fachbegriffe auf einen Blick}
Hidden Channel und Hidden State sind bewusst getrennt: Ein Channel ist eine Komponente, der State der gesamte Vektor.

{\fontsize{9.5}{12}\selectfont\begin{longtable}{>{\raggedright\arraybackslash}p{51mm}>{\raggedright\arraybackslash}p{109mm}}\toprule\textbf{Fachbegriff} & \textbf{Bedeutung} \\ \midrule\endhead
Hidden State h & Gesamter interner Ausgabe-/\allowbreak{}Gedächtnisvektor des LSTM oder der GRU. \\[4pt]
Cell State c & Zusätzlicher interner Speichervektor des LSTM. Nicht der Zustand einer Batteriezelle. \\[4pt]
Hidden Channel /\allowbreak{} Hidden Unit & Eine interne Dimension. Beim LSTM-Pruning werden die zugehörigen Komponenten von h und c und ihre Verbindungen gemeinsam entfernt. \\[4pt]
Hidden Size & Anzahl interner Dimensionen, zum Beispiel 64 beim SOC-LSTM. \\[4pt]
Hidden Layer & Eine ganze interne Netzschicht; nicht dasselbe wie ein einzelner Hidden Channel. \\[4pt]
Feature /\allowbreak{} Input Feature & Eingangsgröße, z. B. Spannung oder eine berechnete Ableitung. \\[4pt]
Training Chunk /\allowbreak{} Training Window & Zusammenhängender Abschnitt einer Zeitreihe, der als Trainingsbeispiel verarbeitet wird. \\[4pt]
Batch /\allowbreak{} Batch Size & Gruppe von Trainingsbeispielen /\allowbreak{} Anzahl der Beispiele in dieser Gruppe. \\[4pt]
Epoch /\allowbreak{} Epoche & Ein Durchgang durch die Trainingsbeispiele; Details wie drop\_\allowbreak{}last können Restbeispiele auslassen. \\[4pt]
Forward Pass /\allowbreak{} Backpropagation & Vorhersage berechnen /\allowbreak{} Fehlergradienten zur Anpassung der Gewichte zurückrechnen. \\[4pt]
Recurrent State & Oberbegriff: GRU hat h; LSTM hat h und c. \\[4pt]
Stateful Streaming & Zustände zwischen aufeinanderfolgenden Eingaben weitergeben. \\[4pt]
Rolling Window & Jeweils einen begrenzten zurückliegenden Abschnitt neu auswerten. \\[4pt]
Structured Pruning & Ganze Hidden Channels mit den zugehörigen Verbindungen entfernen. \\[4pt]
Quantization & Gewichte oder andere Zahlen mit geringerer numerischer Auflösung darstellen. \\[4pt]
Fine-Tuning & Ein bestehendes Modell durch weiteres Training anpassen. \\[4pt]
Loss Function /\allowbreak{} Learning Rate & Trainingsfehlerfunktion /\allowbreak{} Schrittweite für die Anpassung der Gewichte. \\[4pt]
Overfitting /\allowbreak{} Validation Data & Zu starke Anpassung an Trainingsdaten /\allowbreak{} Daten zur Modellkontrolle und Auswahl. \\[4pt]
Backward Difference & Änderungsrate aus aktuellem minus vorherigem Messwert, geteilt durch Zeitabstand. \\[4pt]
Cell-Macro /\allowbreak{} Coordinate-wise Median & Gleich gewichteter Mittelwert über Batteriezellen /\allowbreak{} Median separat für jedes Feature. \\[4pt]
\bottomrule\end{longtable}}

\topic{Methoden verstehen: Was tun sie und warum?}
\textbf{1/4 · Daten aufbereiten und darstellen}

\textbf{Leseschlüssel:} Der Einsatz der folgenden Bausteine ist in Modellcode und Konfigurationen nachvollziehbar. Die Erklärung ihres Zwecks ist eine fachliche Begründung. Ohne einen kontrollierten Vergleich mit und ohne Baustein ist damit kein einzelner Genauigkeitsgewinn nachgewiesen. Die Beispiele beziehen sich vor allem auf das stündliche SOH-Modell aus Abschnitt 6.2.3.

\question{Hourly Aggregation: Warum aus fünf Größen zwanzig Features werden}
\textbf{Was passiert?} Spannung, Strom, Temperatur, EFC und kumulierte Ladung werden pro Stunde jeweils durch Mittelwert, Standardabweichung, Minimum und Maximum beschrieben: fünf Größen mal vier Kennzahlen ergibt zwanzig Features.

\textbf{Warum hier?} SOH beschreibt langsame Alterung. Die Zusammenfassung reduziert die Zahl der Zeitschritte und macht lange Betriebszeiträume mit kürzeren Sequenzen verarbeitbar. Neben dem durchschnittlichen Niveau bleiben Schwankungen und Extremwerte als Hinweise auf die Belastung erhalten.

\textbf{Beispiel und Grenze:} Zwei Stunden können denselben mittleren Strom haben, aber unterschiedlich starke Stromspitzen. Standardabweichung und Maximum unterscheiden diese Fälle. Die genaue Reihenfolge der Ereignisse innerhalb der Stunde geht trotzdem verloren. Stündliche Daten sind deshalb nicht automatisch besser als sekündliche Daten.

\question{Input Scaling / RobustScaler: Warum vor dem Netz skalieren?}
\textbf{Was passiert?} Jedes der zwanzig Features wird mit festen, aus den Trainingsdaten bestimmten Kennwerten umgerechnet. Der verwendete RobustScaler zentriert standardmäßig am Median und skaliert mit dem Interquartilsabstand:
\[
x'_j=\frac{x_j-\operatorname{Median}_{\rm train}(x_j)}{Q_{75,\rm train}(x_j)-Q_{25,\rm train}(x_j)}.
\]
\textbf{Warum hier?} Spannung, Temperatur und kumulierte Größen haben sehr unterschiedliche Zahlenbereiche. Scaling erleichtert die Optimierung, weil reine Einheitenunterschiede nicht die Größenordnung der Eingangswerte bestimmen. Median und Quartile sind weniger empfindlich gegenüber einzelnen Extremwerten als Mittelwert und Standardabweichung.

\textbf{Beispiel und Grenze:} Ein Wert von 30 bei Median 20 und Interquartilsabstand 10 wird zu 1. Das entfernt weder einen Sensorfehler noch neue Ausreißer. Im Betrieb bleiben die Trainingskennwerte fest; sie werden nicht auf Testzellen neu angepasst.

\question{Feature Embedding: Warum noch ein kleines MLP vor dem LSTM?}
\textbf{Was passiert?} Zwei Linear Layers mit Aktivierungen bilden die skalierten zwanzig Features auf 128 interne Werte ab. Diese Gewichte werden gemeinsam mit dem restlichen Modell trainiert.

\textbf{Warum hier?} Das Netz kann bereits vor der zeitlichen Verarbeitung Kombinationen von Stundenkennzahlen lernen. Das LSTM erhält eine auf die Vorhersageaufgabe trainierte Darstellung. Die 128 Werte enthalten keine zusätzlichen Messinformationen; sie sind berechnete Merkmale. Die Breite 128 ist eine Architekturwahl, keine notwendige Eigenschaft eines Embeddings.

\clearpage
\section*{\color{accent}Methoden verstehen: Schichten im Embedding}
\textbf{2/4 · Kombinieren, normalisieren und nichtlinear verarbeiten}

\textbf{Der geprüfte Inferenzpfad:} 20 skalierte Features $\rightarrow$ Linear 20 auf 128 $\rightarrow$ LayerNorm $\rightarrow$ GELU $\rightarrow$ Linear 128 auf 128 $\rightarrow$ GELU. Dropout kommt nur während des Trainings hinzu.

\question{Linear Layer: Was wird hier eigentlich gelernt?}
Jeder Ausgabewert ist eine gewichtete Summe der Eingaben plus Bias:
\[
u_j=\sum_{i=1}^{D}w_{ji}x_i+b_j.
\]
\textbf{Warum hier?} Die Gewichte lernen, welche Eingaben in welcher Kombination zur SOH-Schätzung beitragen. Die erste Matrix hat 128 Zeilen und 20 Spalten: Jede der 128 Zeilen berechnet eine andere Kombination der gleichen zwanzig Eingaben. Der Bias erlaubt zusätzlich eine Verschiebung.

\textbf{Anschaulich:} Ein Zwischenwert könnte Spannung und Temperatur gemeinsam berücksichtigen. Welche Bedeutung ein einzelner gelernter Wert tatsächlich hat, ist ohne Analyse offen. Nur mehrere Linear Layers ohne dazwischenliegende Nichtlinearität wären wieder zu einer einzigen affinen Abbildung zusammenfassbar.

\question{LayerNorm: Warum zusätzlich zum Input Scaling?}
\textbf{Was passiert?} Nach der ersten Linear Layer werden Mittelwert und Varianz über die 128 Zwischenwerte desselben Stundenvektors berechnet. Jeder Wert wird normalisiert und danach mit gelernten Parametern skaliert und verschoben:
\[
y_j=\gamma_j\frac{u_j-\mu_u}{\sqrt{v_u+\varepsilon}}+\beta_j.
\]
Hier sind $\mu_u$ und $v_u$ Mittelwert und Varianz dieses Vektors; $\varepsilon$ stabilisiert die Division. $\gamma_j$ und $\beta_j$ werden gelernt.

\textbf{Warum hier?} Auch aus skalierten Inputs können die trainierbaren Gewichte sehr unterschiedlich große Aktivierungen erzeugen. LayerNorm kontrolliert deren Größenordnung und kann dadurch die Optimierung stabilisieren. Sie verwendet keine Statistik über zukünftige Stunden und benötigt keine große Batch zur Statistikschätzung.

\textbf{Abgrenzung:} Input Scaling nutzt feste Trainingskennwerte pro Eingangsfeature. LayerNorm nutzt die aktuellen Zwischenwerte innerhalb eines einzelnen Vektors. Es wird also nicht nochmals über die Stunde gemittelt und das Input Scaling nicht ersetzt. [L]

\question{GELU: Warum diese Aktivierung?}
\textbf{Was passiert?} GELU verändert jeden Wert nichtlinear und glatt. Anders als ReLU, die negative Werte auf null setzt, kann GELU kleine negative Ausgaben erhalten. Die genaue Definition lautet $\operatorname{GELU}(u)=u\Phi(u)$, wobei $\Phi$ die Verteilungsfunktion der Standardnormalverteilung ist.

\textbf{Warum hier?} Die Aktivierung ermöglicht eine nichtlineare Merkmalsaufbereitung vor dem LSTM. Der glatte Übergang ist eine plausible Motivation für GELU. Ein Vorteil gegenüber ReLU wurde für dieses konkrete SOH-Modell durch die hier geprüften Unterlagen aber nicht isoliert belegt. [G]

\clearpage
\section*{\color{accent}Methoden verstehen: Gedächtnis und Residual Blocks}
\textbf{3/4 · Warum mehrere unterschiedliche Netzbausteine?}

\question{Zwei LSTM Layers: Warum nicht nur ein großes MLP?}
\textbf{Was passiert?} Das erste LSTM verarbeitet die 128 Embedding-Werte einer Stunde und seinen bisherigen Zustand. Sein Hidden State mit 112 Werten wird zum Input des zweiten LSTM, das eigene Hidden und Cell States mit jeweils 112 Werten hat.

\textbf{Warum hier?} Ein einzelner Stundenvektor beschreibt die aktuelle Stunde, Alterung hängt aber von der bisherigen Nutzung ab. Recurrent States ermöglichen eine gelernte Zusammenfassung dieser Historie. Das zweite LSTM kann zeitliche Muster der ersten Darstellung weiterverarbeiten. Zwei Layers sind eine Wahl der Modellkapazität, kein Beleg, dass zwei immer besser als eines sind.

\textbf{Wichtig:} Hidden Size 112 bedeutet 112 interne Werte pro State und Layer, nicht 112 gespeicherte Stunden. Im stündlichen Betrieb werden die States zwischen den Updates weitergegeben.

\question{Residual MLP Block / Skip Connection: Warum die Eingabe addieren?}
\textbf{Was passiert?} Ein MLP-Zweig verarbeitet 112 Werte über eine Zwischenbreite von 160 zurück auf 112. Danach wird der ursprüngliche Input hinzugefügt und das Ergebnis normalisiert. Ohne Trainings-Dropout geschrieben:
\[
F(u)=W_2\operatorname{GELU}(W_1u+b_1)+b_2,
\qquad y=\operatorname{LayerNorm}(u+F(u)).
\]
\textbf{Warum hier?} Der Zweig kann eine Ergänzung zur schon vorhandenen Darstellung lernen, statt sie vollständig neu erzeugen zu müssen. Der zusätzliche direkte Rechenweg kann außerdem den Gradientenfluss beim Training erleichtern. Deshalb heißt es Residual: Der lernbare Zweig liefert eine Änderung zum Input. Im Modell stehen drei solche Blöcke hintereinander. [R]

\textbf{Zahlenbeispiel:} Input $(0{,}4;0{,}8)$ plus gelernte Änderung $(0{,}1;-0{,}2)$ ergibt $(0{,}5;0{,}6)$ vor LayerNorm. Für die Addition müssen beide Vektoren gleich lang sein. Deshalb geht der Block von 112 über 160 wieder auf 112 zurück.

\textbf{Grenze:} Wegen der anschließenden LayerNorm ist der vollständige Block bei $F(u)=0$ nicht zwingend die unveränderte Identität. Die Skip Connection ist auch kein zusätzlicher Speicher über die Zeit; diese Aufgabe übernehmen die LSTM-States.

\question{Output Head, ReLU, Sigmoid und linearer Output}
\textbf{Was passiert?} Der SOH-Head aus Kapitel 6 bildet die Darstellung über 112 auf 160 auf 160 auf einen SOH-Wert ab. GELU liegt zwischen den Linear Layers. Der letzte Output ist linear.

\textbf{Warum hier?} Der Head übersetzt die interne Darstellung in die gesuchte Zielgröße. Ein linearer SOH-Output erzwingt keine Sättigung nahe null oder eins, kann dafür aber auch außerhalb dieses Bereichs liegen. Die SOC-Heads verwenden Sigmoid, um die Ausgabe auf null bis eins zu begrenzen. Die einfacheren Embedded-Heads aus Kapitel 7 verwenden ReLU im Inneren. Das sind verschiedene Architekturentscheidungen, keine austauschbaren Namen für denselben Schritt.

\clearpage
\section*{\color{accent}Methoden verstehen: Training und Aussagegrenzen}
\textbf{4/4 · Warum die Ausführung genauso wichtig ist wie die Architektur}

\question{Dropout: Warum im Training absichtlich Werte ausblenden?}
\textbf{Was passiert?} Während des Trainings werden zufällig Aktivierungen auf null gesetzt; die übrigen werden entsprechend skaliert. Bei der Inferenz ist Dropout deaktiviert.

\textbf{Warum hier?} Das soll eine zu starke Abhängigkeit von einzelnen Aktivierungen reduzieren und Overfitting begrenzen. Es ist eine Regularisierung des Trainings. Es simuliert weder automatisch fehlende Sensorwerte noch löscht es dauerhaft Gewichte wie beim Pruning. Ob die konkrete Dropout-Rate optimal ist, erfordert einen Vergleich.

\question{Training Chunk, Batch und Backpropagation-Horizont}
\textbf{Was passiert?} Ein Chunk ist ein begrenzter zeitlicher Trainingsabschnitt. Eine Batch fasst mehrere solcher Beispiele für ein gemeinsames Gewichtsupdate zusammen. Beim stündlichen SOH-Modell ist die Chunk-Länge 192 Stunden, also acht Tage.

\textbf{Warum hier?} Begrenzte Sequenzen halten Speicherbedarf und Rechenaufwand beim Training beherrschbar. Der Backpropagation-Horizont beschreibt, über wie viele Zeitschritte der Fehler rückwärts verfolgt wird. Daraus folgt nicht automatisch, dass ein Zustand nach acht Tagen gelöscht werden muss. Ein Zustandsreset ist eine separate Entscheidung der Implementierung.

\question{Stateful Streaming und Rolling Window: Warum unterscheiden?}
\textbf{Stateful Streaming:} Jeder neue Input wird mit dem bisherigen State verarbeitet. Beim stündlichen SOH-Modell werden Hidden und Cell States fortgeführt. Das nutzt die vorherige Rechnung und vermeidet vollständige Wiederholungen.

\textbf{Rolling Window:} Beim SOC-DD-Hauptbenchmark wird für jede Schätzung das letzte Fenster von 2.024 Samples neu verarbeitet; zwischen Fenstern wird der State zurückgesetzt. Das erhält die trainierte Fenstersemantik, verursacht aber deutlich mehr Rechenarbeit. Schnelleres Streaming ist deshalb nicht automatisch numerisch derselbe Schätzer.

\question{So begründe ich die Auswahl, ohne zu viel zu behaupten}
\textbf{Belegt:} Die Architektur enthält Input Scaling, ein gelerntes Embedding, LayerNorm, GELU, zwei LSTM Layers und Residual Blocks. Jeder Baustein besitzt eine nachvollziehbare technische Funktion.

\textbf{Plausible Motivation:} Stundenaggregation erschließt lange Zeiträume; das Embedding bereitet Features auf; Normalisierung und Residual Connections können die Optimierung erleichtern; die LSTMs verarbeiten Historie.

\textbf{Nicht allein dadurch belegt:} Dass jeder Baustein den MAE senkt oder die beste Wahl ist. Dafür wäre eine Ablation nötig: einen Baustein gezielt ändern, unter vergleichbaren Bedingungen neu trainieren und auf denselben unabhängigen Daten auswerten. Die Auswahl erfolgt auf Trainings-/Validierungsdaten, nicht durch wiederholtes Optimieren auf dem Holdout.

\textbf{Quellen und Codebezug.} Dissertation Abschnitt 6.2.3, gedruckte Seite 89; Modellkonfiguration und README von SOH\_JES2\_0.1.0; SOH-Modellklasse in train\_soh.py (0.1.2.3); shared\_soh\_trace.py für Stundenfeatures, Scaling und State-Weitergabe. Funktionsprüfung: 25.09.2026. Keine neue Ablationsstudie.

\textbf{Methodenliteratur:} [L] Ba et al., \href{https://arxiv.org/abs/1607.06450}{Layer Normalization (2016)}. [G] Hendrycks und Gimpel, \href{https://arxiv.org/abs/1606.08415}{Gaussian Error Linear Units (2016)}. [R] He et al., \href{https://arxiv.org/abs/1512.03385}{Deep Residual Learning for Image Recognition (2015)}; allgemeines Residual-Prinzip, kein Wirksamkeitsnachweis für dieses SOH-Modell.

\end{document}
